% VestaRV: Castalia analog chapter of the TRM, hand-written.
% Analog content is per-silicon, so this chapter belongs to the implementation and
% not to the shared platform sources; a chip supplies its own
% implementations/asic/<chip>/analog/ and nothing else in the TRM changes.
% LatexUserGuide.CopyAnalogChapter copies this directory to latex/TRM/include/analog/
% at `make chip` time, so paths here are relative to latex/TRM/ (include/analog/...),
% and a chip with no analog/ directory gets no analog chapter.
% The per-block characterisation files (fig_*.tex, tab_*.tex, data/*.dat and the
% BiasGenCascWideSwing.tex master) are generated from Maestro results by
% platform/common/tools/maestro2tex: edit its config and regenerate, never the file.

\section{Analog Circuitry} \label{s:analog}

\AsicNameForUserGuide{} carries a mixed-signal front end alongside the digital
subsystem documented in the preceding chapters. This chapter describes the analog
blocks themselves and reports their simulated performance; the memory-mapped
registers through which software reaches them are documented with the
peripherals\ifcqanalog{} in Section~\ref{ss:cqanalog-regs}, a placeholder map until
the analog IP is integrated\fi.
% ss:cqanalog-regs lives in the generated CqAnalog chapter, which the master template
% inputs only under \ifcqanalog. A config that ships the analog IP with cqAfeStubs
% false (penta_wound_afe.json) renders this chapter without that one, so an unguarded
% \ref is undefined and trm-lint turns that into a build failure.

Unless stated otherwise, every figure and table in this chapter is taken from
schematic-level simulation of the block's own testbench (pre-layout, with no
extracted parasitics), and is therefore a design-intent characterisation rather
than silicon data. Two places run on an extracted netlist and say so where they
do: the SAR ADC section and the code-level probe of the impedance mode.
Each block section names the testbench it came from and the
process, temperature and supply points that were simulated, so the coverage behind
each number is explicit. Numbers are regenerated directly from the simulator's
results database rather than transcribed by hand.

\subsection*{Blocks documented in this revision}

The analog front end is being brought up incrementally. The electrochemical models
the benches use are defined once in Section~\ref{ss:analog-models}, and the nonlinear
electrode model that cyclic voltammetry needs in Section~\ref{ss:analog-electrode-bc}.
This revision documents
the local bias generator, the class-AB amplifier that serves as both the
channel's transimpedance amplifier (A2) and its potentiostat control amplifier
(A1), the complete potentiostat pixel system built from a pair of those
amplifiers around a three-electrode cell, the 10-bit SAR analog-to-digital
converter that closes the channel, the programmable feedback resistor that sets
the channel's transimpedance gain, and the 12-bit R-2R reference DAC with the
supply block that feeds it. Section~\ref{ss:eis} documents not a block but a
second measurement mode of the channel: impedance spectroscopy performed by the
potentiostat itself, with no impedance hardware.
Each block section closes with the register settings
that observe that block on a real die, and Section~\ref{ss:calibration} lists the
per-chip constants those measurements produce. Both are design intent: the analog
test bus and the electrode isolation switches they use are not yet in the
schematics.

\subsection*{Characterisation coverage}

\begin{center}
\begin{tabular}{llccl}
\toprule
Block & Section & Benches & Corners & Monte Carlo \\
\midrule
Bias generator        & \ref{ss:biasgen-local} & 3 & \checkmark$^{\dagger}$ & n/a \\
Class-AB amplifier    & \ref{ss:ampab}         & 6 & \checkmark$^{\dagger}$ & \checkmark \\
Potentiostat pixel system & \ref{ss:pixel}     & 6 & \checkmark$^{\dagger}$ & \checkmark \\
Impedance spectroscopy & \ref{ss:eis}          & 1 & \checkmark$^{\dagger}$ & \checkmark \\
SAR ADC               & \ref{ss:adc}          & 2 & \checkmark & \checkmark \\
Feedback resistor     & \ref{ss:tiarprog}      & 5 & \checkmark & \checkmark \\
Reference DAC         & \ref{ss:dacr2r12}      & 6 & \checkmark$^{\dagger}$ & \checkmark \\
Four-channel macro    & \ref{ss:quad}          & 1 &            &            \\
\bottomrule
\multicolumn{5}{l}{\footnotesize $^{\dagger}$\,Corner run with the passive model
  rows pinned at typical: a lower bound. See below.} \\
\end{tabular}
\end{center}

\noindent Every block now carries a corner run, but a corner run is not the same
claim in every row: the five rows marked \(\dagger\) are cornered over the
transistors alone, for the reason set out immediately below. The impedance row is
not a block but a mode that adds no devices, and its corner and Monte Carlo runs
cover the small-signal measurement only.
The four-channel row is nominal only:
the macro adds no device the block rows do not already corner, and its own claim
is that the four channels are independent, which a corner shifts equally on all
four.
\emph{n/a} marks a block whose
section reports no Monte Carlo because the devices at issue carry no
per-instance mismatch terms: the bias generator's start-up branch sits on model
cards that give a mismatch
run a spread of zero from exactly the devices under test. The converter's Monte
Carlo has been run and is reported as a negative result for a related reason
(Section~\ref{sss:adc-mc}); nothing in this chapter quotes a converter gain or
offset \(\sigma\).
Corner runs cover process, temperature and supply at the five points listed in
each block section. Monte Carlo runs cover process and mismatch together unless a
table states mismatch-only; a corner figure
and a Monte Carlo \(\sigma\) are different quantities and are reported separately
throughout. Design
targets, where quoted alongside a measurement, are stated in the block section that
reports the measurement.

% The pinned-passive caveat, and the single authored place where this defect is
% explained: every affected caption points here in one clause ("pinned-passive caveat,
% page~\pageref{p:pinned-passives}") rather than restating it. Source:
% ~/vestarv/docs/afe_rev2_corner_audit.html and ~/chips/castalia/simruns/
% cornaudit_20260825/ (README.md S3 the headline table, results_tables.txt the four
% measured tables, corner_scan.txt the sweep over every setup); also caught in
% simruns/bgladder_20260825 (why tab_biasgen_dcop_free exists) and
% simruns/eiscorner_20260825 S2.5 (why note_eis_intro.tex carries its own numbers).
% Nothing here re-measures a published table: the affected numbers stand as they were,
% marked as a lower bound. Model section and kit file names are exact here, unlike
% elsewhere in the chapter, because the recommendation is only actionable that way.

\paragraph{The pinned-passive caveat.}\label{p:pinned-passives}
The corner set that serves this chapter moves the sixteen transistor model rows
only. The resistor, capacitor and metal-capacitor rows hold their typical
sections in every corner, so poly sheet resistance, a
\(\pm\SI{35}{\percent}\) quantity, does not move at all. It sets the
programmable feedback resistor of Section~\ref{ss:tiarprog} directly, and it
sets the reference resistor of the local bias generator, and through that the
tail current, transconductance, bandwidth and supply current of every amplifier
in this chapter. Every corner column marked \(\dagger\) above is therefore the
spread of the transistors alone, and is a lower bound on the true spread.

Measured on the same benches with the passive rows freed: the control
amplifier's unity-gain frequency spreads \SIrange{15.89}{36.09}{\mega\hertz}
against the published \SIrange{22.36}{22.91}{\mega\hertz}, a box \(36\times\)
wider, and its supply current \SIrange{81.5}{193.8}{\micro\ampere} against
\SIrange{110.0}{136.7}{\micro\ampere}, \(4.2\times\); worst-case gain margin
falls from \SI{8.10}{\decibel} to \SI{7.60}{\decibel}. On the transimpedance
loop at gain code 63 the feedback resistance spreads \(-33.4\,\%\) and
\(+35.5\,\%\) against the published \(+0.44\,\%\), a box \(157\times\) wider;
its \SI{-3}{\decibel} frequency spreads \SIrange{8.53}{16.14}{\kilo\hertz}
against \SIrange{10.90}{11.32}{\kilo\hertz}, \(18\times\); and worst-case loop
phase margin falls from \SI{71.69}{\degree} to \SI{61.52}{\degree}.

The two mixed corners cannot be repaired. The foundry ships the passive families
with fast, slow and typical sections only (there is no mixed-corner section
for a resistor or a capacitor), so fs and sf had to pin those rows, and the
pinned table was then cloned into ff and ss, where a real section does exist.
That is why the tables look partly right: the bias generator's fs and sf columns
agree with the freed matrix to three digits while its ff and ss columns diverge
by \SI{47}{\percent} and \SI{25}{\percent}, because the two columns that were
always going to be pinned are the two that are still correct. Only ff and ss can
carry the true spread.

\textbf{The Monte Carlo has been the only honest process bound in this chapter
all along.} Its statistical model section varies the resistors, and the
\SI{11.45}{\percent} process \(\sigma\) it measures on the feedback resistance is
\(3\sigma \approx \SI{34}{\percent}\): the box the corner set was missing.
Every Monte Carlo table and figure in this chapter therefore stands as published.
So does the converter of Section~\ref{ss:adc}, whose extracted array carries no
capacitor model card and whose poly strips are fill with no dc path; the
programmable feedback resistor of
Section~\ref{ss:tiarprog}, whose corner axis is the resistor family by
construction and which is the one section that already shows the
\(\pm\SI{35}{\percent}\) band; the offset and tracking tables, which are
matching-limited and move \emph{less} when the passives are freed than the
published corner box already allows; and the reference DAC's zero-scale,
full-scale and LSB columns, which are ratiometric and identical to six digits
pinned or freed.

\paragraph{The corrected corner set.} Define it once as a file rather than as a
table typed into each setup, so that a corrected table cannot be cloned with its
defects: one section per corner, each a complete model-row list. Five sections,
named so that the pin is visible from a corner label:
\texttt{castalia\_tt}, \texttt{castalia\_ff\_all} and \texttt{castalia\_ss\_all}
with every row moved, plus \texttt{castalia\_fs\_mos} and
\texttt{castalia\_sf\_mos}, which cannot move the passives and never will.
Verify from the netlist that executed, never from the setup:

\begin{quote}\small
\verb+grep -E 'cor_(res|mim|rtmom)\.scs' <run>/input.scs+
\end{quote}

\noindent must show the corner letter on the passive rows. One gap survives the
fix: the statistical model section omits the MIM statistics, so MIM capacitor
variation is bounded by neither the corner set nor the Monte Carlo today; the
\SI{22}{\pico\farad} compensation capacitor of Section~\ref{ss:ampab} is the
device that makes that matter.

% Generated block sections. Each is produced by tools/maestro2tex from a
% Maestro run; see that tool's README for how to regenerate or add one.
% Keep a block after whatever it depends on: the amplifier section quotes the
% bias rails that the bias generator section defines, and refers to it by \ref.

% HAND-WRITTEN fragment -- NOT generated by maestro2tex. Input directly from
% AnalogChapter.tex, before the block sections, because every block refers to it.
%
% The component values below are the single source of truth for the electrochemical
% models. Do NOT restate them in figure captions: a caption copy of r_ct drifted to
% 10k against the bench's 30k and shipped in the TRM before this section existed.
% Values read from: castalia/randles_cell and castalia/randles_electrode (OA), and
% confirmed against the instance line in
% simarchive/anatop_tb_controlamp/cm_2026-07-27/Interactive.156/1/CellDriveTB/netlist/.

\subsection{Electrochemical Models}\label{ss:analog-models}

Every electrode measurement in this chapter uses one of the two lumped models of
Figures~\ref{fig:randles-cell} and~\ref{fig:randles-electrode}, valued in
Table~\ref{tab:analog-models}. The electrode impedance sets the feedback factor
of both the potentiostat and the transimpedance loop, so those blocks' stability
figures depend on it directly.

The three-electrode cell is what the potentiostat drives: into the working
electrode, \(R_{\mathrm{ct}} \parallel C_{\mathrm{dl}}\) in series with
\(R_u\), about \SI{31}{\kilo\ohm} at DC and falling toward
\(R_u = \SI{1}{\kilo\ohm}\) above roughly \SI{5}{\hertz}. No bench reported here
uses the single-interface model.

Both models are linear and frequency-independent, with no Warburg term and no
potential dependence, so they give small-signal behaviour about a bias point
rather than a full electrochemical response. Both are also nominal: no element
carries a spread, so nothing in this chapter accounts for electrode-to-electrode
variation, which on a real wound interface is likely to exceed the process
spread reported beside it.

\begin{table}[htbp]
  \centering
  \caption[{Component values of the electrochemical models.}]{Component values of
  the electrochemical models. The three-electrode values are the subcircuit
  defaults; the transimpedance benches override them, and the values used are
  stated with each measurement.}
  \label{tab:analog-models}
  \begin{tabular}{llrl}
    \toprule
    Model & Element & Value & Represents \\
    \midrule
    three-electrode cell      & $R_{\mathrm{ce}}$  & \SI{1}{\kilo\ohm}    & solution resistance, CE to RE \\
                                & $R_u$              & \SI{1}{\kilo\ohm}    & uncompensated resistance, RE to WE surface \\
                                & $R_{\mathrm{ct}}$  & \SI{30}{\kilo\ohm}   & charge transfer at the WE interface \\
                                & $C_{\mathrm{dl}}$  & \SI{1}{\micro\farad} & double layer at the WE interface \\
    \addlinespace
    single interface & $R_s$              & \SI{10}{\kilo\ohm}   & series solution resistance \\
                                & $R_{\mathrm{ct}}$  & \SI{1}{\mega\ohm}    & charge transfer \\
                                & $C_{\mathrm{dl}}$  & \SI{10}{\nano\farad} & double layer \\
    \bottomrule
  \end{tabular}
\end{table}

The closed-loop measurements of Section~\ref{ss:ampab} override the defaults with
$R_u = \SI{100}{\ohm}$ or \SI{1}{\kilo\ohm}, $R_{\mathrm{ct}} =
\SI{10}{\kilo\ohm}$ or \SI{1}{\mega\ohm}, and $C_{\mathrm{dl}} =
\SI{100}{\nano\farad}$ or \SI{1}{\micro\farad}; the value used is stated with
each measurement.

%% Hand-authored circuitikz figure -- NOT generated by maestro2tex.
%% The cell is drawn as the abstracted symbol of sch_cell_symbol.tex, the same
%% circle-with-three-leads used in the Virtuoso schematics. The internal network
%% is deliberately NOT shown: the element values in note_models.tex define the
%% model. Do not restate them here.
\providecommand{\MaestroRoot}{include/analog/}
\input{\MaestroRoot sch_cell_symbol.tex}
\begin{figure}[htbp]
  \centering
\begin{circuitikz}[
    every node/.append style={font=\footnotesize},
    nlab/.style={font=\scriptsize, inner sep=1.5pt},
  ]
\CellSymThree{0,0}{1.0}
\end{circuitikz}
  \caption[{Three-electrode cell model.}]{Three-electrode cell model, the load the potentiostat actually drives, shown as the symbol used throughout this chapter. It is a lumped linear network between the counter, reference and working electrodes; its elements and their values are given in Table~\ref{tab:analog-models}. The model is frequency-independent (no Warburg diffusion term and no potential dependence), so it represents small-signal behaviour about a bias point rather than a full electrochemical response.}
  \label{fig:randles-cell}
\end{figure}

%% Hand-authored circuitikz figure -- NOT generated by maestro2tex.
%% Abstracted symbol (sch_cell_symbol.tex), two leads: the single interface is a
%% two-terminal element. Internals deliberately not shown; values live in
%% note_models.tex.
\providecommand{\MaestroRoot}{include/analog/}
\input{\MaestroRoot sch_cell_symbol.tex}
\begin{figure}[htbp]
  \centering
\begin{circuitikz}[
    every node/.append style={font=\footnotesize},
    nlab/.style={font=\scriptsize, inner sep=1.5pt},
  ]
\CellSymTwo{0,0}{1.0}{\texttt{a}}{\texttt{b}}
\end{circuitikz}
  \caption[{Single-interface model.}]{Single electrode--electrolyte interface, presented as a two-terminal element and shown as the same abstracted symbol. Its elements and their values are given in Table~\ref{tab:analog-models}; they describe a much smaller double layer and a much larger charge-transfer resistance than the three-electrode cell. None of the benches reported in this chapter uses it.}
  \label{fig:randles-electrode}
\end{figure}


% Nonlinear (Butler-Volmer + diffusion) electrode model, used for cyclic voltammetry.
% Hand-written like the linear models above: a simulation model shared by the benches,
% not a block, so it sits with them.
\FloatBarrier
%% HAND-WRITTEN fragment -- NOT generated by maestro2tex. Input directly from
%% AnalogChapter.tex, immediately after note_models.tex: this is the nonlinear
%% counterpart of the linear cell defined there, and both are shared by the
%% benches rather than owned by one block.
%%
%% Sources of the statements below, so they can be re-checked:
%%   model source + equation-to-reference map:
%%     ic/castalia/electrode_bc/veriloga/veriloga.va (header)
%%   grid constants (nseg, h0, beta, delta): same header, "GRID:" line
%%   wrapper topology + parameter pass-through:
%%     simruns/cv_bc_first/.../ps_cv/netlist/input.scs, subckt randles_cell_bc
%%   plan/handoff: ~/vestarv/docs/afe_rev2_cv_electrode.html
%% Parameter defaults live in Table~\ref{tab:electrode-params} -- do not
%% restate them in prose or captions.

\subsection{Voltammetric Electrode Model}\label{ss:analog-electrode-bc}

This section documents a \emph{simulation model}, not a circuit block and not
silicon: \texttt{electrode\_bc} is a Verilog-A element in the \texttt{castalia}
library that stands in for the electrode--electrolyte interface when a bench
sweeps potential and expects a voltammetric peak. Everything reported here is
either a standalone Spectre run of the model against closed-form
electrochemistry, or one transient of the model driven by the rev-2 channel.

The linear models of Section~\ref{ss:analog-models} cannot produce such a peak.
\(R_u + R_{\mathrm{ct}} \parallel C_{\mathrm{dl}}\) is a linear network, so
under a triangular sweep every branch current is a linear functional of the
excitation: a ramp through \(R_{\mathrm{ct}}\) plus a constant
\(C_{\mathrm{dl}}\,\mathrm{d}v/\mathrm{d}t\) box that changes sign at the
vertices. Changing element values moves slopes and never peak positions, so the
extrema stay pinned at the switching potentials
(Figure~\ref{fig:electrode-validation}, dashed). A peak needs two
nonlinearities acting together: Butler--Volmer kinetics, exponential in
overpotential about the couple's formal potential \(E^{0\prime}\), and
diffusional depletion, which collapses the surface concentration as current
flows and turns the exponential rise over into a \(t^{-1/2}\) decay.
Butler--Volmer alone, with the surface concentration pinned at bulk, gives a
sigmoid rather than a peak.

\texttt{electrode\_bc} implements both. Two coupled one-dimensional diffusion
ladders, one per species, share a single Butler--Volmer surface flux. Each
ladder is discretised by the box (finite-volume) method on an exponentially
expanding grid, \(h_j = h_0\beta^{\,j}\) with 28 slabs per species,
\(h_0 = \SI{0.5}{\micro\metre}\) and \(\beta = 1.280\); one slab is one internal
electrical node whose potential carries the slab concentration normalised by
\SI{1e-6}{\mole\per\cubic\centi\metre}, because raw
\si{\mole\per\cubic\centi\metre} sits below the simulator's default
\SI{1}{\micro\volt} voltage tolerance and solves to noise. The outer boundary is
held at bulk concentration at a finite Nernst-layer thickness
\(\delta = h_0(\beta^{28}-1)/(\beta-1) = \SI{1.79}{\milli\metre}\). Parameters
are listed in Table~\ref{tab:electrode-params}; they are in the CGS units of the
source literature (\si{\centi\metre\per\second},
\si{\square\centi\metre\per\second}, \si{\mole\per\cubic\centi\metre}), not SI.
The element emits anodic-positive current, which is the polarity the
transimpedance amplifier measures and the opposite of the cathodic-positive
convention of the reference text.

The benches instantiate it through the wrapper \texttt{randles\_cell\_bc}, which
is the topology of Figure~\ref{fig:randles-cell} with the charge-transfer
resistor replaced: \(R_{\mathrm{ce}}\), \(R_u\) and \(C_{\mathrm{dl}}\) remain
linear elements and \texttt{electrode\_bc} sits between
\(\mathrm{WE_s}\) and WE. The linear cell of Section~\ref{ss:analog-models}
stays the correct model for small-signal impedance work; only voltammetry needs
this one.

%% HAND-WRITTEN table -- NOT generated by maestro2tex.
%% Transcribed from the parameter declarations of
%% ic/castalia/electrode_bc/veriloga/veriloga.va. This table is the single
%% source of truth for these defaults; do not restate them elsewhere.
\begin{table}[htbp]
  \centering
  \caption[{Parameters of the electrode model.}]{Parameters of the
  Verilog-A electrode model, with the defaults the standalone
  validation runs of Table~\ref{tab:electrode-peaks} use. Units are CGS, as in
  the source literature: a bulk concentration of \SI{1}{\milli\mole\per\litre} is
  \SI{1e-6}{\mole\per\cubic\centi\metre} and an area of \SI{1}{\square\milli\metre}
  is \SI{0.01}{\square\centi\metre}. The last two rows are numerical guards, not
  physical quantities; the exponent clamp is never reached over the
  \(\pm\SI{0.8}{\volt}\) window of this front end. The grid constants
  (28~slabs per species, \(h_0\), \(\beta\), \(\delta\)) are compile-time
  constants of the module, not parameters, and are given in
  Section~\ref{ss:analog-electrode-bc}.}
  \label{tab:electrode-params}
  \begin{tabular}{llrll}
    \toprule
    Parameter & Symbol & Default & Unit & Quantity \\
    \midrule
    \texttt{E0}    & \(E^{0\prime}\)      & 0.0     & \si{\volt}                        & formal potential of the couple \\
    \texttt{nel}   & \(n\)                & 1       & ---                               & electrons transferred \\
    \texttt{alpha} & \(\alpha\)           & 0.5     & ---                               & transfer coefficient \\
    \texttt{k0}    & \(k^0\)              & 0.1     & \si{\centi\metre\per\second}       & standard heterogeneous rate constant \\
    \texttt{Dif}   & \(D\)                & \num{5e-6} & \si{\square\centi\metre\per\second} & diffusion coefficient, both species \\
    \texttt{Area}  & \(A\)                & 0.01    & \si{\square\centi\metre}           & electrode area \\
    \texttt{Cox}   & \(C_{\mathrm{O}}^{*}\) & \num{1e-6} & \si{\mole\per\cubic\centi\metre} & bulk concentration, oxidised species \\
    \texttt{Cred}  & \(C_{\mathrm{R}}^{*}\) & 0       & \si{\mole\per\cubic\centi\metre} & bulk concentration, reduced species \\
    \texttt{Temp}  & \(T\)                & 298.15  & \si{\kelvin}                      & interface temperature \\
    \addlinespace
    \texttt{expmax} & ---                 & 40      & ---                               & exponent clamp on the kinetic terms \\
    \texttt{gpar}   & ---                 & \num{1e-15} & \si{\siemens}                 & DC shunt across the terminals \\
    \bottomrule
  \end{tabular}
\end{table}

%% HAND-WRITTEN fragment -- NOT generated by maestro2tex. The engine has no
%% signal-vs-signal (XY) figure kind and these runs are raw Spectre, not a
%% Maestro results database, so the traces are built from .dat files instead.
%%
%% Data: data/electrode_cv_bc.dat, data/electrode_cv_randles.dat,
%% data/electrode_rs_sim.dat, data/electrode_rs_line.dat, all written by
%% ~/chips/castalia/skill/electrode_bc_fig.py from the CSV exports in
%% ~/chips/castalia/electrode_models/csv/ (m3_bvdiff_v*.csv, m1_randles_v*.csv;
%% columns t, E_applied, E_interface, i_WE). Peak values and the
%% Randles-Sevcik slope come from electrode_models/report.txt, the same numbers
%% as Table~\ref{tab:electrode-peaks}. Re-run that script after any re-run of
%% electrode_models/runme.sh.
\providecommand{\MaestroRoot}{include/analog/}
\providecolor{mtxA}{HTML}{2A78D6}
\providecolor{mtxB}{HTML}{EB6834}
\begin{figure}[htbp]
  \centering
  {\pgfplotsset{compat=1.18}%
  \pgfplotsset{
    elecax/.style={
      width=0.50\linewidth, height=5.6cm,
      grid=both,
      major grid style={line width=0.2pt, draw=black!14},
      minor grid style={line width=0.2pt, draw=black!7},
      minor tick num=1,
      axis line style={line width=0.4pt, draw=black!45},
      tick style={line width=0.4pt, draw=black!45},
      tick align=outside,
      label style={font=\small}, tick label style={font=\small},
      legend style={font=\footnotesize, draw=none, fill=white, fill opacity=0.85,
                    text opacity=1, at={(0.03,0.97)}, anchor=north west,
                    row sep=1pt},
      every axis plot/.append style={line join=round},
    }}%
  \begin{tikzpicture}
    \begin{axis}[elecax,
      xlabel={Interface potential \(E-E^{0\prime}\)~[\si{\volt}]},
      ylabel={Electrode current~[\si{\micro\ampere}]},
      enlarge x limits=0.04, enlarge y limits=0.10,
    ]
      \addplot[mtxB, line width=0.9pt, dashed, no marks]
        table {\MaestroRoot data/electrode_cv_randles.dat};
      \addlegendentry{linear cell}
      \addplot[mtxA, line width=1.1pt, solid, no marks]
        table {\MaestroRoot data/electrode_cv_bc.dat};
      \addlegendentry{\texttt{electrode\_bc}}
    \end{axis}
  \end{tikzpicture}\hfill
  \begin{tikzpicture}
    \begin{axis}[elecax,
      xlabel={\(\sqrt{v}\)~[\(\sqrt{\si{\volt\per\second}}\)]},
      ylabel={Peak current \(|i_p|\)~[\si{\micro\ampere}]},
      xmin=0, ymin=0, enlarge x limits=false, enlarge y limits=0.06,
    ]
      \addplot[black, line width=0.8pt, dashed, no marks]
        table {\MaestroRoot data/electrode_rs_line.dat};
      \addlegendentry{Randles--\v{S}ev\v{c}\'ik}
      \addplot[mtxA, only marks, mark=*, mark size=1.9pt]
        table {\MaestroRoot data/electrode_rs_sim.dat};
      \addlegendentry{\texttt{electrode\_bc}}
    \end{axis}
  \end{tikzpicture}}
  \caption[{Validation of \texttt{electrode\_bc} against closed-form
  voltammetry.}]{Standalone validation of \texttt{electrode\_bc} against
  closed-form voltammetry, ideal potential source, model defaults of
  Table~\ref{tab:electrode-params} with \(R_s = \SI{500}{\ohm}\) and
  \(C_{\mathrm{dl}} = \SI{0.2}{\micro\farad}\) in series and in parallel
  respectively, \SI{25}{\celsius}. Left: current against interface potential
  over a \(\pm\SI{0.6}{\volt}\) triangular sweep at \SI{50}{\milli\volt\per\second},
  second cycle, for the nonlinear model (solid) and for the linear cell of
  Section~\ref{ss:analog-models} driven identically (dashed, \(R_{\mathrm{ct}} =
  \SI{1}{\mega\ohm}\)); the linear response is a ramp with its extrema pinned at
  the switching potentials --- its \(C_{\mathrm{dl}}\) box is only
  \SI{20}{\nano\ampere} wide at this scan rate and invisible here --- while the
  nonlinear one peaks \SI{30}{\milli\volt} either side of \(E^{0\prime}\) and
  decays afterwards. Right: peak cathodic current of the
  \emph{first} scan against \(\sqrt{v}\) over
  \SIrange{25}{400}{\milli\volt\per\second}, against the Randles--\v{S}ev\v{c}\'ik
  line, of slope \SI{6.015}{\micro\ampere} per \(\sqrt{\si{\volt\per\second}}\);
  the two panels use different cycles because
  Randles--\v{S}ev\v{c}\'ik describes the first scan into an undepleted bulk,
  while the second cycle is the repeatable steady response and sits
  \SIrange{3}{5}{\percent} below it. Numbers in Table~\ref{tab:electrode-peaks}.}
  \label{fig:electrode-validation}
\end{figure}

%% HAND-WRITTEN table -- NOT generated by maestro2tex.
%% Numbers transcribed from ~/chips/castalia/electrode_models/report.txt:
%%   ip(R-S), ip(Cdl=200n), ip(Cdl=1p), err_on%, err_off%  -- section
%%     "Cdl ABLATION vs RANDLES-SEVCIK" (first cathodic scan)
%%   dEp -- main peak table, rows m3_bvdiff_v*, column "c1 dEp"
%% Machine-readable equivalents: electrode_models/{peaks,ablate}.json.
%% Same run set as Figure~\ref{fig:electrode-validation}.
\begin{table}[htbp]
  \centering
  \caption[{Peak current and peak separation of \texttt{electrode\_bc} against
  Randles--\v{S}ev\v{c}\'ik.}]{Peak metrics of \texttt{electrode\_bc} against
  the Randles--\v{S}ev\v{c}\'ik law, first cathodic scan, conditions as
  Figure~\ref{fig:electrode-validation}. The last-but-one column repeats the
  measurement with \(C_{\mathrm{dl}}\) removed: the remaining
  \SIrange{-0.45}{-0.25}{\percent} is a scan-rate-independent discretisation
  bias of the 28-slab grid, so the whole scan-rate dependence of the error is
  double-layer charging, which the ablation reproduces as
  \(C_{\mathrm{dl}}v/i_p\) to two decimals. Peak separation is inflated by
  \(iR\) drop across \(R_s = \SI{500}{\ohm}\), which grows with scan rate; with
  \(R_s\) and \(C_{\mathrm{dl}}\) removed it is \SI{57.04}{\milli\volt} at every
  rate, against Nicholson's reversible \(2.218RT/F = \SI{56.9}{\milli\volt}\).}
  \label{tab:electrode-peaks}
  \begin{tabular}{rrrrrr}
    \toprule
    Scan rate & \(i_p\) Randles--\v{S}ev\v{c}\'ik & \(i_p\) model & Error & Error, \(C_{\mathrm{dl}}\to 0\) & \(\Delta E_p\) \\
    \si{\milli\volt\per\second} & \si{\micro\ampere} & \si{\micro\ampere} & \si{\percent} & \si{\percent} & \si{\milli\volt} \\
    \midrule
     25 & 0.951 & 0.952 & $+0.08$ & $-0.45$ & 58.78 \\
     50 & 1.345 & 1.349 & $+0.32$ & $-0.42$ & 59.85 \\
    100 & 1.902 & 1.915 & $+0.67$ & $-0.38$ & 60.41 \\
    200 & 2.690 & 2.721 & $+1.16$ & $-0.33$ & 62.33 \\
    400 & 3.804 & 3.875 & $+1.86$ & $-0.25$ & 64.74 \\
    \bottomrule
  \end{tabular}
\end{table}

%% HAND-WRITTEN fragment -- NOT generated by maestro2tex.
%% Every number here comes from ~/chips/castalia/electrode_models/report.txt:
%%   sqrt(v) scaling      -- ratio of the 400 and 25 mV/s rows of tab_electrode_peaks
%%   k0 / Lambda rows     -- "GRID / TIMESTEP / KINETICS CONVERGENCE" and the
%%                           Matsuda-Ayabe block under it
%%   nseg rows            -- "B. INTERNAL-NODE SCALING"
%%   timestep rows        -- t_ms* rows of the convergence table
%%   concentration floor  -- "CONCENTRATION vs TIA RANGE"
%% Machine-readable: electrode_models/{results,converge,conc,probes}.json.

\paragraph{Validation.} Peak current scales \(4.07\times\) between
\SI{25}{\milli\volt\per\second} and \SI{400}{\milli\volt\per\second} against the
\(4.00\times\) of the \(\sqrt{v}\) law, running high against
Randles--\v{S}ev\v{c}\'ik by \SIrange{0.08}{1.86}{\percent}
(Table~\ref{tab:electrode-peaks}). Kinetics walk the Matsuda--Ayabe progression
correctly: at \SI{50}{\milli\volt\per\second},
\(k^0 = \SI{0.1}{\centi\metre\per\second}\) gives \(\Lambda = 18.1\), reversible
(\(\Delta E_p = \SI{60.0}{\milli\volt}\), \(i_p\) within \SI{0.11}{\percent}),
and \(k^0 = \SI{1e-3}{\centi\metre\per\second}\) gives \(\Lambda = 0.181\),
\(\Delta E_p = \SI{167}{\milli\volt}\), \(i_p\) \SI{15.5}{\percent} low.

Discretisation, not the solver, is the accuracy limit: eight slabs per species
is unusable (\(i_p\) high by \SI{232}{\percent}), 16 and above agree within
\SI{0.35}{\percent}, and 28 ships. Dropping the maximum time step from
\SI{20}{\milli\second} to \SI{0.5}{\milli\second} moves \(i_p\) below the fourth
significant figure. Roughly \num{60} sweeps ran without a convergence failure,
and the exponentials survive \(\pm\SI{12}{\volt}\) of overpotential with the
clamp removed, so the clamp and shunt conductance are insurance.

One design limit follows. At \SI{50}{\milli\volt\per\second} on
\SI{1}{\square\milli\metre}, double-layer charging is \SI{0.74}{\percent} of the
peak at \SI{1}{\milli\mole\per\litre} but \SI{37}{\percent} of a
\SI{27}{\nano\ampere} peak at \SI{0.02}{\milli\mole\per\litre}: below roughly
\SI{0.1}{\milli\mole\per\litre} the faradaic peak is buried in charging current
rather than in the front end's noise, and the remedy is a slower scan, not more
gain.

%% HAND-WRITTEN fragment -- NOT generated by maestro2tex.
%% Source of every number: the archived headless run
%%   ~/chips/castalia/simruns/cv_bc_37C/   (electrode Temp=310.15, matching temp=37)
%%     summary_ps_cv.txt                                  -> i_peak
%%     castalia/tb_anatop_pixel/maestro/results/maestro/Ocean.0/1/ps_cv/
%%       netlist/input.scs                                -> topology, tran line,
%%                                                           temp, ECS params
%%       psf/variables_file                               -> design variables
%% Predicted peak: Randles-Sevcik at the ECS parameters plus Cdl*dv/dt.
%% Sign convention measured on this run; the TB-07 formula in
%% ~/vestarv/docs/afe_rev2_simplan.html is inverted.

\paragraph{First measurement with the rev-2 channel in the loop.} The
cyclic-voltammetry test is the
potentiostat bench of Section~\ref{ss:pixel} with the cell master swapped from
the linear cell to the nonlinear electrode cell and nothing else changed:
the class-AB amplifier of Section~\ref{ss:ampab} as both A1 and A2, the local
bias generator of Section~\ref{ss:biasgen-local}, an ideal
\(R_f = \SI{35}{\kilo\ohm}\) rather than the programmable resistor of
Section~\ref{ss:tiarprog}, and \(C_{\mathrm{dl}} = \SI{0.2}{\micro\farad}\),
\(R_u = R_{\mathrm{ce}} = \SI{1}{\kilo\ohm}\) around an electrode at
\(E^{0\prime} = \SI{0.3}{\volt}\), \SI{1}{\milli\mole\per\litre},
\(k^0 = \SI{0.1}{\centi\metre\per\second}\). The commanded potential is a
\SIrange{0.75}{1.75}{\volt} triangle of \SI{0.5}{\second} per ramp, i.e.
\SI{2}{\volt\per\second}; transient to \SI{1.2}{\second} at
\texttt{errpreset=conservative}, typical process point, \SI{37}{\celsius} for
silicon and electrode alike (the model's \texttt{Temp} parameter is set to
\SI{310.15}{\kelvin} to match the simulator temperature).

The measured peak is \SI{8.666}{\micro\ampere} against
\SI{8.74}{\micro\ampere} predicted --- \SI{8.34}{\micro\ampere} of
Randles--\v{S}ev\v{c}\'ik faradaic current evaluated at \SI{310.15}{\kelvin}
plus \SI{0.40}{\micro\ampere} of \(C_{\mathrm{dl}}\,\mathrm{d}v/\mathrm{d}t\)
--- a \SI{-0.8}{\percent} discrepancy, so the nonlinearity survives the loop
rather than being reshaped by it. Current is read from the output voltage as
\(i_{\mathrm{we}} = (v_{\mathrm{out}} - V_{\mathrm{CM}})/R_f\): A2 holds WE at
\(V_{\mathrm{CM}}\) and the whole electrode current returns through \(R_f\), so
\(v_{\mathrm{out}} = V_{\mathrm{CM}} + i_{\mathrm{we}}R_f\). The peak therefore
moves the output \SI{0.30}{\volt} from \(V_{\mathrm{CM}} = \SI{1.25}{\volt}\),
\SI{38}{\percent} of the \(\pm\SI{0.80}{\volt}\) transimpedance window at this
gain, so the sweep is measured unclipped. The four wound-panel analytes run as
Maestro corners on the same test (each corner overrides
\(E^{0\prime}\), \(k^0\), \(D\), \(n\) and the bulk concentrations);
Figure~\ref{fig:electrode-panel} shows the resulting voltammograms.

\paragraph{What these numbers do not cover.} Five limits, each of which changes
how a figure above must be read.

\begin{itemize}[nosep]
  \item The in-loop results are nominal transients: one process point, one
    temperature, no process-corner set and no Monte Carlo. The analyte corners
    of Figure~\ref{fig:electrode-panel} vary only the electrode parameters and
    all run at the nominal process point.
  \item A Verilog-A element carries no mismatch statistics. The electrode model
    contributes nothing to a Monte Carlo run and must never be the thing one
    measures; every Monte Carlo \(\sigma\) in this chapter holds the electrode
    exactly nominal and reports circuit variation alone.
  \item The standalone validation of Table~\ref{tab:electrode-peaks} runs at
    \SI{25}{\celsius} throughout; the in-loop run above is the only result at
    \SI{37}{\celsius}. \texttt{Temp} is a plain model parameter, not tied to
    the simulator temperature --- any new bench must set it explicitly or the
    interface kinetics silently revert to \SI{25}{\celsius}. \(D\) is held at
    its \SI{25}{\celsius} value in both cases; the \SIrange{25}{37}{\celsius}
    increase in a real aqueous \(D\) (\(\sim\)30\,\%) is not modelled.
  \item Randles--\v{S}ev\v{c}\'ik assumes semi-infinite diffusion; the model
    pins bulk concentration at \(\delta = \SI{1.79}{\milli\metre}\). The
    difference is negligible while \(\delta \gg \sqrt{Dt_{\mathrm{scan}}}\),
    which holds by \(16\times\) at \SI{50}{\milli\volt\per\second}
    (\(\sqrt{Dt} = \SI{0.11}{\milli\metre}\)) and by far more at
    \SI{2}{\volt\per\second}. It does not hold for a thin unstirred film: a
    \SI{93}{\micro\metre} variant of the model exists for that case and
    attenuates the reverse peak to \(0.65\) of the forward one. Quote \(\delta\)
    with any peak taken from this model.
  \item Peak position and height identify and quantify an analyte; \(k^0\)
    only shapes the peak, and it is the least-supported parameter of the set.
    No heterogeneous rate constant in \si{\centi\metre\per\second} is published
    on carbon for any of the four wound markers the rev-2 panel targets --- the
    values in use are fitted to reproduce a measured \(\Delta E_p\) or
    \(\mathrm{d}E_p/\mathrm{d}\log v\). That is a gap in the literature, not in
    the model.
\end{itemize}

%% HAND-WRITTEN fragment -- NOT generated by maestro2tex (XY figure kind).
%%
%% Data: data/electrode_panel_{ua,aa,h2o2,pyo}.dat, written by
%% ~/chips/castalia/skill/cv_panel_fig.py --dat from the Maestro corner run
%% archived at ~/chips/castalia/simruns/cv_panel (tb_anatop_pixel/CvBvTB, test
%% ps_cv, analyte corners UricAcid/Ascorbate/H2O2_PB/Pyocyanin defined in that
%% maestro view; 5/5 completed, summary_ps_cv.txt). Columns: E = V(RE)-V_CM [V],
%% i_we = (V_CM - V_OUT)/R_f [uA] (anodic positive). Same data as ISCAS27
%% Fig. 4. Re-run the corner set and the script to regenerate.
\providecommand{\MaestroRoot}{include/analog/}
\providecolor{mtxA}{HTML}{2A78D6}
\providecolor{mtxB}{HTML}{EB6834}
\providecolor{mtxC}{HTML}{1BAF7A}
\providecolor{mtxD}{HTML}{EDA100}
\begin{figure}[htbp]
  \centering
  {\pgfplotsset{compat=1.18}%
  \begin{tikzpicture}
    \begin{axis}[
      width=0.78\linewidth, height=6.4cm,
      grid=both,
      major grid style={line width=0.2pt, draw=black!14},
      minor grid style={line width=0.2pt, draw=black!7},
      minor tick num=1,
      axis line style={line width=0.4pt, draw=black!45},
      tick style={line width=0.4pt, draw=black!45},
      tick align=outside,
      label style={font=\small}, tick label style={font=\small},
      legend style={font=\footnotesize, draw=none, fill=white,
                    fill opacity=0.85, text opacity=1,
                    at={(0.03,0.97)}, anchor=north west, row sep=1pt},
      every axis plot/.append style={line join=round},
      xlabel={\(E = V_{\mathrm{RE}} - V_{\mathrm{CM}}\)~[\si{\volt}]},
      ylabel={\(i_{\mathrm{we}}\)~[\si{\micro\ampere}]},
      enlarge x limits=0.03, enlarge y limits=0.08,
    ]
      \addplot[mtxA, line width=1.0pt, no marks]
        table {\MaestroRoot data/electrode_panel_ua.dat};
      \addlegendentry{uric acid, \SI{500}{\micro\mole\per\litre}}
      \addplot[mtxB, line width=1.0pt, no marks]
        table {\MaestroRoot data/electrode_panel_aa.dat};
      \addlegendentry{ascorbate, \SI{100}{\micro\mole\per\litre}}
      \addplot[mtxC, line width=1.0pt, no marks]
        table {\MaestroRoot data/electrode_panel_h2o2.dat};
      \addlegendentry{H$_2$O$_2$/PB, \SI{60}{\micro\mole\per\litre}}
      \addplot[mtxD, line width=1.0pt, no marks]
        table {\MaestroRoot data/electrode_panel_pyo.dat};
      \addlegendentry{pyocyanin, \SI{50}{\micro\mole\per\litre}}
      \addplot[black!45, line width=0.4pt, domain=-0.42:0.82, samples=2,
        forget plot] {0};
    \end{axis}
  \end{tikzpicture}}
  \caption[{Four-analyte voltammograms through the rev-2 channel.}]{%
  Voltammograms of the four-analyte wound panel through the rev-2 channel, one
  electrode parameter set per channel, from the cyclic-voltammetry
  Maestro corner run of Section~\ref{ss:analog-electrode-bc}
  (\SI{2}{\volt\per\second}, \(R_f = \SI{35}{\kilo\ohm}\), conditions of the
  in-loop run above). Peak potential identifies the analyte and peak height its
  concentration; each trace is one full cycle started at that analyte's
  non-faradaic vertex. \(E^{0\prime}\) and concentrations are
  literature-anchored; \(k^{0}\) is fitted (no measured value exists on carbon
  for any of the four), and reversal peaks appear for the chemically
  irreversible couples because following chemistry is not modelled. The two
  cathodic markers span \SIrange{-1.4}{-2.4}{\micro\ampere} --- unmeasurable by
  the unidirectional rev-1 front end.}
  \label{fig:electrode-panel}
\end{figure}

%% HAND-WRITTEN fragment -- NOT generated by maestro2tex.
%% Source of every number: ~/chips/castalia/simruns/dpv_panel (57 standalone
%% Spectre transients of the CvBvTB pixel netlist with a pwl DPV source;
%% orchestrator ~/work/ieee/ISCAS27/latek/graphics-data/sim_figdata.py, engine
%% ~/chips/castalia/skill/dpv_engine.py, summary_first_full_run.txt).
%% Cell names deliberately kept out of the rendered prose.

\paragraph{Differential-pulse voltammetry through the channel.} The same bench
runs a staircase-plus-pulse program at the published assay envelope
(Figure~\ref{fig:electrode-dpv}), every \SI{4}{\milli\volt} tread quantised to the 14-bit DAC
LSB of \SI{152.6}{\micro\volt} and current sampled \SI{5}{\milli\second} before
each edge. Four reporters give four resolved peaks, apexes within
\SI{2}{\milli\volt} and \SI{1}{\percent} of the assay targets. Three findings
bound how the numbers read:

\begin{itemize}[nosep]
  \item \emph{No residual charging baseline.} A blank gives under
    \SI{1}{\pico\ampere} at every step, \SI{45}{\milli\second} of a
    \SI{50}{\milli\second} pulse being far beyond the
    \SI{200}{\nano\farad}/\SI{1}{\kilo\ohm} time constant, so
    \(\Delta I = I(S_2)-I(S_1)\) cancels charging entirely.
  \item \emph{The \SI{560}{\kilo\ohm} tap slew-limits on every edge.} A
    \SI{50}{\milli\volt} step into \SI{200}{\nano\farad} is
    \SI{10}{\nano\coulomb} against a \(\sim\)\SI{2.2}{\micro\ampere} output
    stage, leaving the applied potential wrong for
    \(\sim\)\SI{20}{\milli\second}. Settled samples are unaffected, but the
    transfer turns superlinear above \(\sim\)\SI{350}{\nano\ampere} (up to
    \SI{30}{\percent} at the top of the tap; proportional within
    \SI{3}{\percent} below it), which is why
    Figure~\ref{fig:electrode-dose} uses a \emph{measured} transfer.
  \item \emph{Full scale is consumed by \(I(S_2)\), not \(\Delta I\).} The
    largest sampled current is \(1.02\)--\(1.04\times\) the differential peak
    and reaches \SI{84}{\percent} of the \(\pm\SI{2.06}{\micro\ampere}\) tap;
    do headroom accounting on \(I(S_2)\).
\end{itemize}

%% HAND-WRITTEN fragment -- NOT generated by maestro2tex (XY figure kind).
%%
%% Data: data/electrode_dpv_{exc,samp,ca125,he4,ca199,cea,peaks}.dat, converted
%% from the simulated CSVs written by
%% ~/work/ieee/ISCAS27/latek/graphics-data/sim_figdata.py (engine:
%% ~/chips/castalia/skill/dpv_engine.py; runs archived at
%% ~/chips/castalia/simruns/dpv_panel, 57 transients, summary in
%% summary_first_full_run.txt). Bench: the CvBvTB pixel netlist with
%% VSRC_PATTERN replaced by a 14-bit-quantised pwl DPV staircase; Rf = 560k,
%% c_dl = 200n, temp 37. Excitation excerpt window is t = 10..12 s of the
%% 42.6 s scan. Columns: exc/samp = t [s], E [V]; voltammogram = E_base [V],
%% dI [nA] per reporter; peaks = apex coordinates.
\providecommand{\MaestroRoot}{include/analog/}
\providecolor{mtxA}{HTML}{2A78D6}
\providecolor{mtxB}{HTML}{EB6834}
\providecolor{mtxC}{HTML}{1BAF7A}
\providecolor{mtxD}{HTML}{EDA100}
\begin{figure}[htbp]
  \centering
  {\pgfplotsset{compat=1.18}%
  \pgfplotsset{
    elecdpvax/.style={
      width=0.50\linewidth, height=5.6cm,
      grid=both,
      major grid style={line width=0.2pt, draw=black!14},
      minor grid style={line width=0.2pt, draw=black!7},
      minor tick num=1,
      axis line style={line width=0.4pt, draw=black!45},
      tick style={line width=0.4pt, draw=black!45},
      tick align=outside,
      label style={font=\small}, tick label style={font=\small},
      legend style={font=\footnotesize, draw=none, fill=white,
                    fill opacity=0.85, text opacity=1,
                    at={(0.03,0.97)}, anchor=north west, row sep=1pt},
      every axis plot/.append style={line join=round},
    }}%
  \begin{tikzpicture}
    \begin{axis}[elecdpvax,
      xlabel={Time~[\si{\second}]},
      ylabel={\(E = V_{\mathrm{RE}} - V_{\mathrm{CM}}\)~[\si{\volt}]},
      enlarge x limits=0.02, enlarge y limits=0.10,
    ]
      \addplot[mtxA, line width=0.8pt, no marks]
        table {\MaestroRoot data/electrode_dpv_exc.dat};
      \addlegendentry{excitation}
      \addplot[mtxB, only marks, mark=*, mark size=1.9pt]
        table {\MaestroRoot data/electrode_dpv_samp.dat};
      \addlegendentry{\(S_1\), \(S_2\)}
    \end{axis}
  \end{tikzpicture}\hfill
  \begin{tikzpicture}
    \begin{axis}[elecdpvax,
      xlabel={Base potential \(E\)~[\si{\volt}]},
      ylabel={\(\Delta I = I(S_2)-I(S_1)\)~[\si{\nano\ampere}]},
      enlarge x limits=0.03, enlarge y limits=0.08,
    ]
      \addplot[mtxA, line width=1.0pt, no marks]
        table {\MaestroRoot data/electrode_dpv_ca125.dat};
      \addlegendentry{CA-125}
      \addplot[mtxB, line width=1.0pt, no marks]
        table {\MaestroRoot data/electrode_dpv_he4.dat};
      \addlegendentry{HE4}
      \addplot[mtxC, line width=1.0pt, no marks]
        table {\MaestroRoot data/electrode_dpv_ca199.dat};
      \addlegendentry{CA19-9}
      \addplot[mtxD, line width=1.0pt, no marks]
        table {\MaestroRoot data/electrode_dpv_cea.dat};
      \addlegendentry{CEA}
      \addplot[black, only marks, mark=o, mark size=1.9pt, forget plot]
        table {\MaestroRoot data/electrode_dpv_peaks.dat};
    \end{axis}
  \end{tikzpicture}}
  \caption[{Differential-pulse voltammetry through the rev-2 channel.}]{%
  Differential-pulse voltammetry through the rev-2 channel, four-reporter
  immunoassay panel, one electrode parameter set per reporter
  (\(R_f = \SI{560}{\kilo\ohm}\), \(C_{\mathrm{dl}} = \SI{200}{\nano\farad}\),
  \SI{37}{\celsius}). Top: a \SI{2}{\second} excerpt of the excitation ---
  \SI{4}{\milli\volt} treads, \SI{50}{\milli\volt} / \SI{50}{\milli\second}
  pulses at a \SI{200}{\milli\second} period, every level quantised to the
  14-bit DAC LSB --- with the two current-sampling instants \SI{5}{\milli\second}
  before each edge. Bottom: the differential current over the full
  \SIrange{-0.300}{0.548}{\volt} scan (213 steps, \SI{42.6}{\second});
  peak potential separates the reporters, peak height carries concentration
  (circles mark the apexes: \SIlist{855;1300;1666;938}{\nano\ampere}).
  Reporter kinetics as in the cyclic-voltammetry runs; the panel differs from
  Figure~\ref{fig:electrode-panel} in readout, gain tap and reporter set.}
  \label{fig:electrode-dpv}
\end{figure}

%% HAND-WRITTEN fragment -- NOT generated by maestro2tex (XY figure kind).
%%
%% Data: data/electrode_dose_{ca125,he4,ca199,cea,cut}.dat, converted from the
%% simulated CSVs of ~/work/ieee/ISCAS27/latek/graphics-data/sim_figdata.py.
%% The binding (concentration -> surface reporter activity) math is the
%% published-assay model of that script; the current transfer through it is
%% MEASURED from the DPV runs (11 concentrations per reporter at each
%% reporter's own peak potential, inverted as a table -- linearity was checked,
%% not assumed: <3 % below ~350 nA, up to ~30 % superlinear at the top of the
%% 560k tap; see note_electrode_dpv.tex). Columns: c [marker-specific clinical
%% unit], i [nA]. cut = the four clinical decision points through the same
%% chain.
\providecommand{\MaestroRoot}{include/analog/}
\providecolor{mtxA}{HTML}{2A78D6}
\providecolor{mtxB}{HTML}{EB6834}
\providecolor{mtxC}{HTML}{1BAF7A}
\providecolor{mtxD}{HTML}{EDA100}
\begin{figure}[htbp]
  \centering
  {\pgfplotsset{compat=1.18}%
  \begin{tikzpicture}
    \begin{axis}[
      width=0.72\linewidth, height=6.2cm,
      xmode=log, ymode=log,
      grid=both,
      major grid style={line width=0.2pt, draw=black!14},
      minor grid style={line width=0.2pt, draw=black!7},
      axis line style={line width=0.4pt, draw=black!45},
      tick style={line width=0.4pt, draw=black!45},
      tick align=outside,
      label style={font=\small}, tick label style={font=\small},
      legend style={font=\footnotesize, draw=none, fill=white,
                    fill opacity=0.85, text opacity=1,
                    at={(0.97,0.03)}, anchor=south east, row sep=1pt},
      every axis plot/.append style={line join=round},
      xlabel={Analyte concentration [marker-specific clinical unit]},
      ylabel={DPV peak current \(\Delta I_p\)~[\si{\nano\ampere}]},
    ]
      \addplot[mtxA, line width=1.0pt, no marks]
        table {\MaestroRoot data/electrode_dose_ca125.dat};
      \addlegendentry{CA-125}
      \addplot[mtxB, line width=1.0pt, no marks]
        table {\MaestroRoot data/electrode_dose_he4.dat};
      \addlegendentry{HE4}
      \addplot[mtxC, line width=1.0pt, no marks]
        table {\MaestroRoot data/electrode_dose_ca199.dat};
      \addlegendentry{CA19-9}
      \addplot[mtxD, line width=1.0pt, no marks]
        table {\MaestroRoot data/electrode_dose_cea.dat};
      \addlegendentry{CEA}
      \addplot[black, only marks, mark=*, mark size=1.9pt, forget plot]
        table {\MaestroRoot data/electrode_dose_cut.dat};
    \end{axis}
  \end{tikzpicture}}
  \caption[{Dose response of the four-reporter panel through the channel.}]{%
  Dose response of the four-reporter panel: published-assay binding curves
  composed with the current transfer \emph{measured} from the DPV runs of
  Figure~\ref{fig:electrode-dpv} (11 simulated concentrations per reporter,
  \SI{560}{\kilo\ohm} tap). Filled circles are the clinical decision points
  through the same chain: CA-125 35~U/mL \(\to\) \SI{517}{\nano\ampere}, HE4
  \SI{70}{\pico\mole\per\litre} \(\to\) \SI{409}{\nano\ampere}, CA19-9
  37~U/mL \(\to\) \SI{313}{\nano\ampere}, CEA
  \SI{5}{\nano\gram\per\milli\litre} \(\to\) \SI{185}{\nano\ampere}. The concentration axis carries each marker's own
  clinical unit; only the current axis is common. The binding model is
  analytic; the electrode and channel behind the current axis are simulated.}
  \label{fig:electrode-dose}
\end{figure}

%% HAND-WRITTEN table -- NOT generated by maestro2tex.
%% Mirrors the equation-to-source map in the header of
%% ic/castalia/electrode_bc/veriloga/veriloga.va, and section 8 of
%% ~/vestarv/docs/afe_rev2_cv_electrode.html. The TRM carries no bibliography,
%% so the full citations are given in the prose paragraph that follows.
\begin{table}[htbp]
  \centering
  \caption[{Equations implemented by the electrode model and their
  sources.}]{Every relation the electrode model implements, in the form it is
  implemented, with the source it is taken from; the same map is carried in the
  model's own header. The lower group is not implemented --- those are the
  closed forms the model is validated against in Table~\ref{tab:electrode-peaks}.
  Here \(\eta = E - E^{0\prime}\) and \(f = F/RT\). Equation numbers are those of
  the second edition of Bard and Faulkner; the third edition renumbers.}
  \label{tab:electrode-refs}
  \small
  \begin{tabular}{>{\raggedright\arraybackslash}p{0.21\linewidth}%
                  >{\raggedright\arraybackslash}p{0.45\linewidth}l}
    \toprule
    Element & Form & Source \\
    \midrule
    \multicolumn{3}{l}{\emph{Implemented}} \\
    Equivalent circuit & \(R_s + C_{\mathrm{dl}}\parallel Z_f\); the model is \(Z_f\) & Randles 1947 \\
    Faradaic current & \(i = nFAk^0\big[C_{\mathrm{R}}(0,t)e^{(1-\alpha)nf\eta}\) \(-\;C_{\mathrm{O}}(0,t)e^{-\alpha nf\eta}\big]\) & B\&F Eq.~3.3.11 \\
    Mass transport & \(\partial C/\partial t = D\,\partial^2 C/\partial x^2\) & B\&F Eq.~4.4.2 \\
    Surface boundary & \(i/(nFA) = D\,\partial C_{\mathrm{O}}/\partial x |_{x=0}\) & B\&F Eq.~4.4.29 \\
    Outer boundary & \(C(\delta,t) = C^{*}\), Nernst diffusion layer & B\&F Sec.~1.4.2 \\
    Discretisation & box (finite-volume) method & Feldberg 1969; B\&S Ch.~3 \\
    Space grid & \(h_j = h_0\beta^{\,j}\) & B\&S Ch.~7 \\
    \addlinespace
    \multicolumn{3}{l}{\emph{Validated against}} \\
    Peak current & \(i_p = (\num{2.69e5})\,n^{3/2}AD^{1/2}C^{*}v^{1/2}\) & B\&F Eq.~6.2.19 \\
    Peak separation & \(\Delta E_p = 2.218\,RT/nF\) & Nicholson 1965 \\
    Reversibility & \(\Lambda = k^0/\sqrt{\pi D n f v}\) & Matsuda \& Ayabe 1955 \\
    \bottomrule
  \end{tabular}
\end{table}

%% HAND-WRITTEN fragment -- NOT generated by maestro2tex.
%% The TRM has no bibliography and no \cite anywhere, so the sources abbreviated
%% in Table~\ref{tab:electrode-refs} are spelled out here in prose. Keep this
%% list in step with the header of ic/castalia/electrode_bc/veriloga/veriloga.va.

\paragraph{Sources.} The abbreviations of Table~\ref{tab:electrode-refs} are:
B\&F, A.~J. Bard and L.~R. Faulkner, \emph{Electrochemical Methods:
Fundamentals and Applications}, 2nd ed., Wiley, 2001; B\&S, D.~Britz and
J.~Strutwolf, \emph{Digital Simulation in Electrochemistry}, 4th ed., Springer,
2016; Randles 1947, J.~E.~B. Randles, \emph{Discuss. Faraday Soc.} 1:11, 1947;
Feldberg 1969, S.~W. Feldberg, \emph{Electroanal. Chem.} 3:199, 1969;
Nicholson 1965, R.~S. Nicholson, \emph{Anal. Chem.} 37:1351, 1965;
Matsuda \& Ayabe 1955, H.~Matsuda and Y.~Ayabe, \emph{Z. Elektrochem.} 59:494,
1955. The peak-current law is due independently to Randles
(\emph{Trans. Faraday Soc.} 44:327, 1948) and \v{S}ev\v{c}\'ik
(\emph{Coll. Czech. Chem. Commun.} 13:349, 1948); the numerical form quoted in
Table~\ref{tab:electrode-refs} is the \SI{25}{\celsius} one of B\&F
Eq.~6.2.19.

\FloatBarrier

% Each block master is followed by a hand-written note_meas_*.tex: the register
% settings that observe that block on a real die through the analog test bus or the
% converter's own input multiplexer. They live outside the masters because the masters
% are maestro2tex output and are overwritten on every regeneration.

% maestro2tex -- generated master include for BiasGenCascWideSwing.
%% Regenerate with tools/maestro2tex; do not edit by hand.
%%
%% Usage from the TRM (which builds with cwd = latex/TRM/):
%     % maestro2tex -- generated master include for BiasGenCascWideSwing.
%% Regenerate with tools/maestro2tex; do not edit by hand.
%%
%% Usage from the TRM (which builds with cwd = latex/TRM/):
%     % maestro2tex -- generated master include for BiasGenCascWideSwing.
%% Regenerate with tools/maestro2tex; do not edit by hand.
%%
%% Usage from the TRM (which builds with cwd = latex/TRM/):
%     \input{include/analog/BiasGenCascWideSwing.tex}
%
%% Needs pgfplots, booktabs, siunitx and placeins -- all already in
%% packages-commands.tex. \providecolor needs xcolor (also present).

\providecommand{\MaestroRoot}{include/analog/}

% Categorical palette: slots 1-5 of the validated reference palette,
% in documented order. Dash patterns are a secondary encoding for print.
\providecolor{mtxA}{HTML}{2A78D6}
\providecolor{mtxB}{HTML}{EB6834}
\providecolor{mtxC}{HTML}{1BAF7A}
\providecolor{mtxD}{HTML}{EDA100}
\providecolor{mtxE}{HTML}{E87BA4}

% Keep this block's floats inside this block (see placeins).
\FloatBarrier

\subsection{Local Wide-Swing Cascode Bias Generator}\label{ss:biasgen-local}

The local bias generator (schematic cell \texttt{anatop\_biasgen\_local}, drawn from
\texttt{BiasGenCascWideSwing}) is the cell each analog block instantiates to develop
its own cascode bias rails, as distinct from the single chip-level
\texttt{BiasGenCascWideSwingGlobal} that drives the global \texttt{bn!}/\texttt{bp!}
nets. It presents four outputs, the NMOS mirror and cascode gates
($V_{\mathrm{BNM}}$, $V_{\mathrm{BNC}}$) and the PMOS mirror and cascode gates
($V_{\mathrm{BPM}}$, $V_{\mathrm{BPC}}$), generated from a resistor-degenerated
reference core in the wide-swing cascode arrangement, so the cascode gates sit only
one saturation voltage away from the mirror gates and the mirrors keep their output
devices in saturation over a wide compliance range. An \texttt{En} input gates the
cell, and a start-up branch guarantees it leaves the degenerate zero-current state at
power-up. A 6-bit trim word, \texttt{BiasAdj\textless5:0\textgreater}, moves the
reference current so the operating point can be corrected after silicon.

The characterisation below is from schematic-level simulation (no extracted
parasitics) of the \texttt{BiasGenCascWideSwing\allowbreak\_tb} testbench, over the process,
temperature and supply points of
Table~\ref{tab:BiasGenCascWideSwing-corners}. All rail voltages are referred to
\texttt{vss!}, and the cell is enabled throughout.

\input{\MaestroRoot sch_biasgen_core.tex}
\input{\MaestroRoot sch_biasgen_startup.tex}
\input{\MaestroRoot tab_BiasGenCascWideSwing_corners.tex}
\input{\MaestroRoot tab_biasgen_dcop.tex}
\input{\MaestroRoot fig_biasgen_rails_vs_temp.tex}
\input{\MaestroRoot tab_biasgen_temp_stats.tex}
\input{\MaestroRoot tab_biasgen_tempco.tex}
\input{\MaestroRoot fig_biasgen_vbnm_vs_temp.tex}
\input{\MaestroRoot fig_biasgen_idd_vs_temp.tex}
\input{\MaestroRoot tab_biasgen_line.tex}
\input{\MaestroRoot fig_biasgen_vbnm_vs_vdd.tex}
\input{\MaestroRoot fig_biasgen_idd_vs_vdd.tex}
\input{\MaestroRoot tab_biasgen_trim.tex}
\input{\MaestroRoot fig_biasgen_vbnm_vs_code.tex}
\input{\MaestroRoot fig_biasgen_idd_vs_code.tex}
\input{\MaestroRoot tab_biasgen_startup_time.tex}
\input{\MaestroRoot fig_biasgen_startup.tex}
\input{\MaestroRoot fig_biasgen_startup_corners.tex}

\FloatBarrier

%
%% Needs pgfplots, booktabs, siunitx and placeins -- all already in
%% packages-commands.tex. \providecolor needs xcolor (also present).

\providecommand{\MaestroRoot}{include/analog/}

% Categorical palette: slots 1-5 of the validated reference palette,
% in documented order. Dash patterns are a secondary encoding for print.
\providecolor{mtxA}{HTML}{2A78D6}
\providecolor{mtxB}{HTML}{EB6834}
\providecolor{mtxC}{HTML}{1BAF7A}
\providecolor{mtxD}{HTML}{EDA100}
\providecolor{mtxE}{HTML}{E87BA4}

% Keep this block's floats inside this block (see placeins).
\FloatBarrier

\subsection{Local Wide-Swing Cascode Bias Generator}\label{ss:biasgen-local}

The local bias generator (schematic cell \texttt{anatop\_biasgen\_local}, drawn from
\texttt{BiasGenCascWideSwing}) is the cell each analog block instantiates to develop
its own cascode bias rails, as distinct from the single chip-level
\texttt{BiasGenCascWideSwingGlobal} that drives the global \texttt{bn!}/\texttt{bp!}
nets. It presents four outputs, the NMOS mirror and cascode gates
($V_{\mathrm{BNM}}$, $V_{\mathrm{BNC}}$) and the PMOS mirror and cascode gates
($V_{\mathrm{BPM}}$, $V_{\mathrm{BPC}}$), generated from a resistor-degenerated
reference core in the wide-swing cascode arrangement, so the cascode gates sit only
one saturation voltage away from the mirror gates and the mirrors keep their output
devices in saturation over a wide compliance range. An \texttt{En} input gates the
cell, and a start-up branch guarantees it leaves the degenerate zero-current state at
power-up. A 6-bit trim word, \texttt{BiasAdj\textless5:0\textgreater}, moves the
reference current so the operating point can be corrected after silicon.

The characterisation below is from schematic-level simulation (no extracted
parasitics) of the \texttt{BiasGenCascWideSwing\allowbreak\_tb} testbench, over the process,
temperature and supply points of
Table~\ref{tab:BiasGenCascWideSwing-corners}. All rail voltages are referred to
\texttt{vss!}, and the cell is enabled throughout.

\input{\MaestroRoot sch_biasgen_core.tex}
\input{\MaestroRoot sch_biasgen_startup.tex}
\input{\MaestroRoot tab_BiasGenCascWideSwing_corners.tex}
\input{\MaestroRoot tab_biasgen_dcop.tex}
\input{\MaestroRoot fig_biasgen_rails_vs_temp.tex}
\input{\MaestroRoot tab_biasgen_temp_stats.tex}
\input{\MaestroRoot tab_biasgen_tempco.tex}
\input{\MaestroRoot fig_biasgen_vbnm_vs_temp.tex}
\input{\MaestroRoot fig_biasgen_idd_vs_temp.tex}
\input{\MaestroRoot tab_biasgen_line.tex}
\input{\MaestroRoot fig_biasgen_vbnm_vs_vdd.tex}
\input{\MaestroRoot fig_biasgen_idd_vs_vdd.tex}
\input{\MaestroRoot tab_biasgen_trim.tex}
\input{\MaestroRoot fig_biasgen_vbnm_vs_code.tex}
\input{\MaestroRoot fig_biasgen_idd_vs_code.tex}
\input{\MaestroRoot tab_biasgen_startup_time.tex}
\input{\MaestroRoot fig_biasgen_startup.tex}
\input{\MaestroRoot fig_biasgen_startup_corners.tex}

\FloatBarrier

%
%% Needs pgfplots, booktabs, siunitx and placeins -- all already in
%% packages-commands.tex. \providecolor needs xcolor (also present).

\providecommand{\MaestroRoot}{include/analog/}

% Categorical palette: slots 1-5 of the validated reference palette,
% in documented order. Dash patterns are a secondary encoding for print.
\providecolor{mtxA}{HTML}{2A78D6}
\providecolor{mtxB}{HTML}{EB6834}
\providecolor{mtxC}{HTML}{1BAF7A}
\providecolor{mtxD}{HTML}{EDA100}
\providecolor{mtxE}{HTML}{E87BA4}

% Keep this block's floats inside this block (see placeins).
\FloatBarrier

\subsection{Local Wide-Swing Cascode Bias Generator}\label{ss:biasgen-local}

The local bias generator (schematic cell \texttt{anatop\_biasgen\_local}, drawn from
\texttt{BiasGenCascWideSwing}) is the cell each analog block instantiates to develop
its own cascode bias rails, as distinct from the single chip-level
\texttt{BiasGenCascWideSwingGlobal} that drives the global \texttt{bn!}/\texttt{bp!}
nets. It presents four outputs, the NMOS mirror and cascode gates
($V_{\mathrm{BNM}}$, $V_{\mathrm{BNC}}$) and the PMOS mirror and cascode gates
($V_{\mathrm{BPM}}$, $V_{\mathrm{BPC}}$), generated from a resistor-degenerated
reference core in the wide-swing cascode arrangement, so the cascode gates sit only
one saturation voltage away from the mirror gates and the mirrors keep their output
devices in saturation over a wide compliance range. An \texttt{En} input gates the
cell, and a start-up branch guarantees it leaves the degenerate zero-current state at
power-up. A 6-bit trim word, \texttt{BiasAdj\textless5:0\textgreater}, moves the
reference current so the operating point can be corrected after silicon.

The characterisation below is from schematic-level simulation (no extracted
parasitics) of the \texttt{BiasGenCascWideSwing\allowbreak\_tb} testbench, over the process,
temperature and supply points of
Table~\ref{tab:BiasGenCascWideSwing-corners}. All rail voltages are referred to
\texttt{vss!}, and the cell is enabled throughout.

\input{\MaestroRoot sch_biasgen_core.tex}
\input{\MaestroRoot sch_biasgen_startup.tex}
\input{\MaestroRoot tab_BiasGenCascWideSwing_corners.tex}
\input{\MaestroRoot tab_biasgen_dcop.tex}
\input{\MaestroRoot fig_biasgen_rails_vs_temp.tex}
\input{\MaestroRoot tab_biasgen_temp_stats.tex}
\input{\MaestroRoot tab_biasgen_tempco.tex}
\input{\MaestroRoot fig_biasgen_vbnm_vs_temp.tex}
\input{\MaestroRoot fig_biasgen_idd_vs_temp.tex}
\input{\MaestroRoot tab_biasgen_line.tex}
\input{\MaestroRoot fig_biasgen_vbnm_vs_vdd.tex}
\input{\MaestroRoot fig_biasgen_idd_vs_vdd.tex}
\input{\MaestroRoot tab_biasgen_trim.tex}
\input{\MaestroRoot fig_biasgen_vbnm_vs_code.tex}
\input{\MaestroRoot fig_biasgen_idd_vs_code.tex}
\input{\MaestroRoot tab_biasgen_startup_time.tex}
\input{\MaestroRoot fig_biasgen_startup.tex}
\input{\MaestroRoot fig_biasgen_startup_corners.tex}

\FloatBarrier

% Start-up verification, hand-written: it comes from a bench built after the
% maestro2tex config was frozen (castalia/tb_anatop_bg_startup, one test bg_su, 101
% corners), so it is input here rather than from the generated master. Source
% ~/chips/castalia/simruns/bgladder_20260825/README.md, Task A.
% HAND-WRITTEN fragment -- NOT generated by maestro2tex. Input from
% AnalogChapter.tex immediately after BiasGenCascWideSwing.tex, which is
% generated and must not be hand-edited; the bench behind this fragment was
% built after that master's config was frozen.
%
% Source: ~/chips/castalia/simruns/bgladder_20260825/README.md, Task A, with
% tables/bg_startup_matrix.tsv (101 rows) and bg_startup/summary_bg_su.txt.
% Bench castalia/tb_anatop_bg_startup view BgStartupTB, maestro view, one test
% bg_su, 101 corners + nominal; schematic level, no extracted parasitics.
% The bench differs from BiasGenCascWideSwing_tb in exactly two ways: the
% supply is a 7-breakpoint vpwl (a vpulse starts AT its brown-out level, i.e.
% it answers the question by assuming it) and En is unstrapped from vdd! so the
% enable can be exercised independently.  Cross-checked at tt/27 C/BiasAdj 33:
% 612.75 / 775.84 / 1817.09 / 1589.34 mV and 24.4602 uA.
%
% Numbers used below:
%  101/101 corners completed, 0 simulation errors, 0 spec failures.
%  Ramp 100n/1u/10u/100u/1m/10m x {tt,ff,fs,sf,ss} at 25 C, nominal trim.
%  bnm_f and idd_f spread over the six ramp rates = 0.000000 at every corner.
%  t_set at tt: 211n/232n, 741n/994n, 4.40u/9.94u, 43.5u/99.4u, 436u/994u,
%    4.36m/9.94m for V_BNM / V_BPC; ratios 2.11/2.32 then 0.741/0.994 then
%    0.436-0.440 / 0.994 above 10 us.  Worst circuit-limited floor 487 ns /
%    641 ns at ss.  Slowest observed overall wc_slow_ss_hot_adj0, 7.28 ms on a
%    10 ms ramp.
%  Temperature 0/25/37/85 C x 5 corners; BiasAdj 0/33/63 x 5 corners; six
%    deliberate worst-case cells.  All start.
%  Worst case over all 101: V_BNM 508.19-735.14, V_BNC 694.32-953.90,
%    V_BPM 1691.96-1920.44, V_BPC 1316.84-1669.79 mV, I_DD 12.81-84.33 uA
%    (extremes trim-driven: BiasAdj 63 at ss/hot, 0 at ff/cold).
%  Brown-out depths 2.0 down to 0.0 V at tt/25 C plus 1.2/1.0/0.8/0.6 V at
%    ss/85 C and ff/0 C; 20 us, 1 ms and 100 ms holds; slow (1 ms) recovery
%    edge.  bnm_ratio = 1.000000 to six digits everywhere.  V_BNM collapse at
%    the dip: 0.16 % at 2.0 V, 1.38 % at 1.0 V, 6.02 % at 0.8 V, 19.0 % at
%    0.6 V, 72.4 % at 0.0 V/20 us, 83.8 % at 0.0 V/1 ms, 96.8 % (tt) and
%    99.96 % (ss 85 C, to 0.27 mV) at 0.0 V/100 ms.
%  Enable: En low 20 us with the rail at 2.5 V, tt/ss/ff x 25/85 C.  V_BNM =
%    0.00000 V exactly, V_SU 3.7-4.1 mV; bnm_ratio = idd_ratio = 1.000000 on
%    re-assert at all six cells.
%  Mismatch caveat: the four start-up devices sit on raw nch_25 cards with no
%    mismatchflag parameter at all; the two kick devices' Vth is what sets how
%    hard the kick is.  Corner/temperature/trim/ramp/brown-out results are
%    unaffected -- the corner is applied through model-card parameters that the
%    raw card consumes.
% Library cell names are kept out of the rendered prose.

\subsubsection{Start-up}\label{sss:biasgen-startup}

Start-up is verified across 101 corner points --- supply ramp rate over five
decades, five process corners, four temperatures and three trim codes, plus a
brown-out depth sweep and an enable cycle --- and all 101 start, none into the
degenerate state. The settled operating point is bit-identical across all five
decades of ramp rate: the spread of \(V_{\mathrm{BNM}}\) and of \(I_{\mathrm{DD}}\) over
the six ramp rates is zero to six digits at every process corner. Ramp rate
changes when the generator arrives, not where.

\begin{table}[htbp]
  \centering
  \caption[{Settling time of the bias rails against supply ramp rate.}]{Settling time of the two extreme bias rails against supply ramp rate, at tt and \SI{25}{\celsius}, defined as the last instant outside a \SI{1}{\percent} band around the value reached; the ratio is settling time over ramp time. Above \SI{10}{\micro\second} the block is ramp-limited and the rails settle at a fixed fraction of the ramp; below \SI{1}{\micro\second} it is circuit-limited, and the floor is \SI{211}{\nano\second} and \SI{232}{\nano\second} here, \SI{487}{\nano\second} and \SI{641}{\nano\second} at the slow corner.}
  \label{tab:biasgen-ramp}
  \begin{tabular}{lrrrr}
    \toprule
    Supply ramp & \(t_{\mathrm{set}}(V_{\mathrm{BNM}})\) & ratio & \(t_{\mathrm{set}}(V_{\mathrm{BPC}})\) & ratio \\
    \midrule
    \SI{100}{\nano\second}  & \SI{211}{\nano\second}  & 2.11  & \SI{232}{\nano\second}  & 2.32 \\
    \SI{1}{\micro\second}   & \SI{741}{\nano\second}  & 0.741 & \SI{994}{\nano\second}  & 0.994 \\
    \SI{10}{\micro\second}  & \SI{4.40}{\micro\second}& 0.440 & \SI{9.94}{\micro\second}& 0.994 \\
    \SI{100}{\micro\second} & \SI{43.5}{\micro\second}& 0.435 & \SI{99.4}{\micro\second}& 0.994 \\
    \SI{1}{\milli\second}   & \SI{436}{\micro\second} & 0.436 & \SI{994}{\micro\second} & 0.994 \\
    \SI{10}{\milli\second}  & \SI{4.36}{\milli\second}& 0.436 & \SI{9.94}{\milli\second}& 0.994 \\
    \bottomrule
  \end{tabular}
\end{table}

\paragraph{No brown-out floor exists.} The supply is ramped up, held, dipped and
returned, at thirteen depths from \SI{2.0}{\volt} down to \SI{0}{\volt} and at
three hold times. Every dip recovers exactly: the ratio of \(V_{\mathrm{BNM}}\)
after to before is \num{1.000000} to six digits at every depth and every corner,
and the supply current returns to its pre-dip value. How far the rails actually
fall is the honest caveat --- at \SI{0}{\volt} the only discharge path for the
rail decoupling is subthreshold leakage, so a short dip is not a power-down and
\(V_{\mathrm{BNM}}\) sags only \SI{72}{\percent}. Held \SI{100}{\milli\second}
at \SI{0}{\volt} the rails do collapse, to \SI{0.27}{\milli\volt} at the slow
hot corner, and recovery is still exact. A recovery edge slower than
\SI{1}{\milli\second} is untested.

The enable input is the stronger test, and it passes. De-asserted, it pulls both
NMOS rails to ground and holds the start-up node down, so the cell is forced
into the true zero-current state --- \(V_{\mathrm{BNM}} =
\SI{0.00000}{\volt}\), against the \SI{96.8}{\percent} collapse a
\SI{100}{\milli\second} power-down manages. On re-assert both the rails and the
supply current return to \num{1.000000} of their previous values at all six
cells tested.

One limit on the coverage. The devices that do the starting sit on model cards
that carry no per-instance mismatch terms, so a mismatch Monte Carlo of start-up
as the cell stands would report zero contribution from exactly the mechanism
under test, and none has been run. The results above are corner-based and do not
depend on it: the process corner reaches those devices through the model-card
parameters they do consume.

% HAND-WRITTEN fragment -- NOT generated by maestro2tex. Input from
% AnalogChapter.tex immediately after BiasGenCascWideSwing.tex, which is
% generated and must not be hand-edited.
%
% Sources: analog test-bus fabric, slot map and register fields from
% ~/vestarv/docs/afe_rev2_debug_mux.html §6 (fields) and §11 (slots) --
% chip-level slots 2-5 = bn!/bnc!/bp!/bpc!, per-channel slots 10/11 = vbn_ch/
% vbp_ch (observe through atb_o only), chip-level slot 8 = iref_ovr.
% Rail values from tab_biasgen_dcop; trim span from tab_biasgen_trim; MC
% sigma/mu 12.5 % from bg_mc (Plan.4.Run.2). Bring-up ordering from §9.
% Library cell names are kept out of the rendered prose.

\subsubsection{Measuring this block on chip}

Set \texttt{AFE\_DBGOWNSEL}~=~5 and step \texttt{AFE\_DBGGSEL} over slots 2 to 5
to reach the four global rails; a single channel's own rails are per-channel
slots 10 and 11 under \texttt{AFE\_DBGSEL}. Read all of them on the buffered pad
\texttt{atb\_o} with \texttt{AFE\_DBGBUFEN}~=~1 and never on the direct pad:
these are gate nets with picoamperes of drive, and the direct path moves what it
measures. Expect the tt column of Table~\ref{tab:biasgen-dcop} ---
\SI{612.8}{\milli\volt} and \SI{775.8}{\milli\volt} on the NMOS pair,
\SI{1817.1}{\milli\volt} and \SI{1589.3}{\milli\volt} on the PMOS pair, each
inside a corner spread of \SIrange{117}{142}{\milli\volt}; a rail or a near-zero
on any of the four means the start-up branch did not fire, and no measurement
elsewhere on the die is meaningful until it does. Trim from the supply pin
rather than through the multiplexer: step the 6-bit code and keep the one whose
total analog supply current is nearest design, \SI{2.7}{\percent} per code over
a \SI{39.9}{\micro\ampere} span (Table~\ref{tab:biasgen-trim}) against a
\SI{12.5}{\percent} Monte Carlo spread the corner runs do not show, so the code
is per die and \(\pm\tfrac{1}{2}\) code leaves \(\pm1.4\,\%\). Chip-level slot 8
drives the reference-current node from \texttt{atb\_d} with
\texttt{AFE\_DBGDIREN}~=~1, which is the last path to a working amplifier on a
die whose generator did not start. This test fabric is design intent; it is not
yet in the schematics.

% maestro2tex -- generated master include for AmpAB.
%% Regenerate with tools/maestro2tex; do not edit by hand.
%%
%% Usage from the TRM (which builds with cwd = latex/TRM/):
%     % maestro2tex -- generated master include for AmpAB.
%% Regenerate with tools/maestro2tex; do not edit by hand.
%%
%% Usage from the TRM (which builds with cwd = latex/TRM/):
%     % maestro2tex -- generated master include for AmpAB.
%% Regenerate with tools/maestro2tex; do not edit by hand.
%%
%% Usage from the TRM (which builds with cwd = latex/TRM/):
%     \input{include/analog/AmpAB.tex}
%
%% Needs pgfplots, booktabs, siunitx and placeins -- all already in
%% packages-commands.tex. \providecolor needs xcolor (also present).

\providecommand{\MaestroRoot}{include/analog/}

% Categorical palette: slots 1-5 of the validated reference palette,
% in documented order. Dash patterns are a secondary encoding for print.
\providecolor{mtxA}{HTML}{2A78D6}
\providecolor{mtxB}{HTML}{EB6834}
\providecolor{mtxC}{HTML}{1BAF7A}
\providecolor{mtxD}{HTML}{EDA100}
\providecolor{mtxE}{HTML}{E87BA4}

% Keep this block's floats inside this block (see placeins).
\FloatBarrier

\input{\MaestroRoot note_ampab_intro.tex}
\input{\MaestroRoot sch_TiaAmpAB_topology.tex}
\input{\MaestroRoot tab_AmpAB_corners.tex}
\input{\MaestroRoot sch_tb_unityfb.tex}
\input{\MaestroRoot tab_ampab_unity.tex}
\input{\MaestroRoot tab_ampab_openloop.tex}
\input{\MaestroRoot tab_ampab_offset.tex}
\input{\MaestroRoot tab_ampab_drive.tex}
\input{\MaestroRoot tab_ampab_power.tex}
\input{\MaestroRoot note_ampab_device_mc.tex}
\input{\MaestroRoot fig_ampab_mc_voffset.tex}
\input{\MaestroRoot tab_ampab_mc_offset.tex}
\input{\MaestroRoot tab_ampab_mc_unity.tex}
\input{\MaestroRoot tab_ampab_mc_openloop.tex}
\input{\MaestroRoot fig_ampab_mc_power.tex}
\input{\MaestroRoot tab_ampab_mc_power.tex}
\input{\MaestroRoot tab_ampab_mc_drive.tex}
\input{\MaestroRoot note_ampab_tia.tex}
\input{\MaestroRoot sch_tb_tiafb.tex}
\input{\MaestroRoot tab_TiaAmpAB_corners.tex}
\input{\MaestroRoot tab_tiaab_accuracy.tex}
\input{\MaestroRoot fig_tiaab_zt.tex}
\input{\MaestroRoot tab_tiaab_transfer.tex}
\input{\MaestroRoot fig_tiaab_loopgain.tex}
\input{\MaestroRoot tab_tiaab_stability.tex}
\input{\MaestroRoot tab_tiaab_noise.tex}
\input{\MaestroRoot note_ampab_tia_mc.tex}
\input{\MaestroRoot fig_tiaab_mc_voff.tex}
\input{\MaestroRoot tab_tiaab_mc_offset.tex}
\input{\MaestroRoot tab_tiaab_mc_accuracy.tex}
\input{\MaestroRoot fig_tiaab_mc_pm.tex}
\input{\MaestroRoot tab_tiaab_mc_stability.tex}
\input{\MaestroRoot tab_tiaab_mc_bandwidth.tex}
\input{\MaestroRoot fig_tiaab_mc_inoise.tex}
\input{\MaestroRoot tab_tiaab_mc_noise.tex}
\input{\MaestroRoot note_ampab_ctrl.tex}
\input{\MaestroRoot sch_tb_celldrive.tex}
\input{\MaestroRoot tab_ampab_cellstab.tex}
\input{\MaestroRoot fig_ampab_cell_ce.tex}
\input{\MaestroRoot tab_ampab_tracking.tex}
\input{\MaestroRoot tab_ampab_compliance.tex}
\input{\MaestroRoot note_ampab_ctrl_mc.tex}
\input{\MaestroRoot fig_ampab_mc_cellaol.tex}
\input{\MaestroRoot tab_ampab_mc_cellstab.tex}
\input{\MaestroRoot tab_ampab_mc_tracking.tex}
\input{\MaestroRoot tab_ampab_mc_regulation.tex}
\input{\MaestroRoot tab_ampab_mc_compliance.tex}

\FloatBarrier

%
%% Needs pgfplots, booktabs, siunitx and placeins -- all already in
%% packages-commands.tex. \providecolor needs xcolor (also present).

\providecommand{\MaestroRoot}{include/analog/}

% Categorical palette: slots 1-5 of the validated reference palette,
% in documented order. Dash patterns are a secondary encoding for print.
\providecolor{mtxA}{HTML}{2A78D6}
\providecolor{mtxB}{HTML}{EB6834}
\providecolor{mtxC}{HTML}{1BAF7A}
\providecolor{mtxD}{HTML}{EDA100}
\providecolor{mtxE}{HTML}{E87BA4}

% Keep this block's floats inside this block (see placeins).
\FloatBarrier

\input{\MaestroRoot note_ampab_intro.tex}
\input{\MaestroRoot sch_TiaAmpAB_topology.tex}
\input{\MaestroRoot tab_AmpAB_corners.tex}
\input{\MaestroRoot sch_tb_unityfb.tex}
\input{\MaestroRoot tab_ampab_unity.tex}
\input{\MaestroRoot tab_ampab_openloop.tex}
\input{\MaestroRoot tab_ampab_offset.tex}
\input{\MaestroRoot tab_ampab_drive.tex}
\input{\MaestroRoot tab_ampab_power.tex}
\input{\MaestroRoot note_ampab_device_mc.tex}
\input{\MaestroRoot fig_ampab_mc_voffset.tex}
\input{\MaestroRoot tab_ampab_mc_offset.tex}
\input{\MaestroRoot tab_ampab_mc_unity.tex}
\input{\MaestroRoot tab_ampab_mc_openloop.tex}
\input{\MaestroRoot fig_ampab_mc_power.tex}
\input{\MaestroRoot tab_ampab_mc_power.tex}
\input{\MaestroRoot tab_ampab_mc_drive.tex}
\input{\MaestroRoot note_ampab_tia.tex}
\input{\MaestroRoot sch_tb_tiafb.tex}
\input{\MaestroRoot tab_TiaAmpAB_corners.tex}
\input{\MaestroRoot tab_tiaab_accuracy.tex}
\input{\MaestroRoot fig_tiaab_zt.tex}
\input{\MaestroRoot tab_tiaab_transfer.tex}
\input{\MaestroRoot fig_tiaab_loopgain.tex}
\input{\MaestroRoot tab_tiaab_stability.tex}
\input{\MaestroRoot tab_tiaab_noise.tex}
\input{\MaestroRoot note_ampab_tia_mc.tex}
\input{\MaestroRoot fig_tiaab_mc_voff.tex}
\input{\MaestroRoot tab_tiaab_mc_offset.tex}
\input{\MaestroRoot tab_tiaab_mc_accuracy.tex}
\input{\MaestroRoot fig_tiaab_mc_pm.tex}
\input{\MaestroRoot tab_tiaab_mc_stability.tex}
\input{\MaestroRoot tab_tiaab_mc_bandwidth.tex}
\input{\MaestroRoot fig_tiaab_mc_inoise.tex}
\input{\MaestroRoot tab_tiaab_mc_noise.tex}
\input{\MaestroRoot note_ampab_ctrl.tex}
\input{\MaestroRoot sch_tb_celldrive.tex}
\input{\MaestroRoot tab_ampab_cellstab.tex}
\input{\MaestroRoot fig_ampab_cell_ce.tex}
\input{\MaestroRoot tab_ampab_tracking.tex}
\input{\MaestroRoot tab_ampab_compliance.tex}
\input{\MaestroRoot note_ampab_ctrl_mc.tex}
\input{\MaestroRoot fig_ampab_mc_cellaol.tex}
\input{\MaestroRoot tab_ampab_mc_cellstab.tex}
\input{\MaestroRoot tab_ampab_mc_tracking.tex}
\input{\MaestroRoot tab_ampab_mc_regulation.tex}
\input{\MaestroRoot tab_ampab_mc_compliance.tex}

\FloatBarrier

%
%% Needs pgfplots, booktabs, siunitx and placeins -- all already in
%% packages-commands.tex. \providecolor needs xcolor (also present).

\providecommand{\MaestroRoot}{include/analog/}

% Categorical palette: slots 1-5 of the validated reference palette,
% in documented order. Dash patterns are a secondary encoding for print.
\providecolor{mtxA}{HTML}{2A78D6}
\providecolor{mtxB}{HTML}{EB6834}
\providecolor{mtxC}{HTML}{1BAF7A}
\providecolor{mtxD}{HTML}{EDA100}
\providecolor{mtxE}{HTML}{E87BA4}

% Keep this block's floats inside this block (see placeins).
\FloatBarrier

\input{\MaestroRoot note_ampab_intro.tex}
\input{\MaestroRoot sch_TiaAmpAB_topology.tex}
\input{\MaestroRoot tab_AmpAB_corners.tex}
\input{\MaestroRoot sch_tb_unityfb.tex}
\input{\MaestroRoot tab_ampab_unity.tex}
\input{\MaestroRoot tab_ampab_openloop.tex}
\input{\MaestroRoot tab_ampab_offset.tex}
\input{\MaestroRoot tab_ampab_drive.tex}
\input{\MaestroRoot tab_ampab_power.tex}
\input{\MaestroRoot note_ampab_device_mc.tex}
\input{\MaestroRoot fig_ampab_mc_voffset.tex}
\input{\MaestroRoot tab_ampab_mc_offset.tex}
\input{\MaestroRoot tab_ampab_mc_unity.tex}
\input{\MaestroRoot tab_ampab_mc_openloop.tex}
\input{\MaestroRoot fig_ampab_mc_power.tex}
\input{\MaestroRoot tab_ampab_mc_power.tex}
\input{\MaestroRoot tab_ampab_mc_drive.tex}
\input{\MaestroRoot note_ampab_tia.tex}
\input{\MaestroRoot sch_tb_tiafb.tex}
\input{\MaestroRoot tab_TiaAmpAB_corners.tex}
\input{\MaestroRoot tab_tiaab_accuracy.tex}
\input{\MaestroRoot fig_tiaab_zt.tex}
\input{\MaestroRoot tab_tiaab_transfer.tex}
\input{\MaestroRoot fig_tiaab_loopgain.tex}
\input{\MaestroRoot tab_tiaab_stability.tex}
\input{\MaestroRoot tab_tiaab_noise.tex}
\input{\MaestroRoot note_ampab_tia_mc.tex}
\input{\MaestroRoot fig_tiaab_mc_voff.tex}
\input{\MaestroRoot tab_tiaab_mc_offset.tex}
\input{\MaestroRoot tab_tiaab_mc_accuracy.tex}
\input{\MaestroRoot fig_tiaab_mc_pm.tex}
\input{\MaestroRoot tab_tiaab_mc_stability.tex}
\input{\MaestroRoot tab_tiaab_mc_bandwidth.tex}
\input{\MaestroRoot fig_tiaab_mc_inoise.tex}
\input{\MaestroRoot tab_tiaab_mc_noise.tex}
\input{\MaestroRoot note_ampab_ctrl.tex}
\input{\MaestroRoot sch_tb_celldrive.tex}
\input{\MaestroRoot tab_ampab_cellstab.tex}
\input{\MaestroRoot fig_ampab_cell_ce.tex}
\input{\MaestroRoot tab_ampab_tracking.tex}
\input{\MaestroRoot tab_ampab_compliance.tex}
\input{\MaestroRoot note_ampab_ctrl_mc.tex}
\input{\MaestroRoot fig_ampab_mc_cellaol.tex}
\input{\MaestroRoot tab_ampab_mc_cellstab.tex}
\input{\MaestroRoot tab_ampab_mc_tracking.tex}
\input{\MaestroRoot tab_ampab_mc_regulation.tex}
\input{\MaestroRoot tab_ampab_mc_compliance.tex}

\FloatBarrier

% HAND-WRITTEN fragment -- NOT generated by maestro2tex. Input from
% AnalogChapter.tex immediately after AmpAB.tex, which is generated.
%
% Salvaged from the retired note_calibration_seq.tex steps 4 and 5 and
% note_calibration_paths.tex C2, plus the fabric of
% ~/vestarv/docs/afe_rev2_debug_mux.html: per-channel slot 5 = ce_amp (A1
% output, inboard of the C1 isolation switch), slot 4 = vout, AFE_ADCSRCSEL = 4
% routes ATB into the SAR (§6.3). The "do not buffer it" argument is §4.1's
% chicken-and-egg: a test buffer built from the same 2.5 V devices carries an
% offset of the same class as A1 (sigma 5.20 mV) and A2 (sigma 8.91 mV).
% Numbers: tab_ampab_mc_offset, tab_tiaab_mc_offset. No cell names in prose.

\subsubsection{Measuring this block on chip}

Input offset is the term that decides yield, and a loopback is the only way to
see it. Open the reference-electrode isolation switch, close the unity-gain
loopback around A1, select per-channel slot 5 --- the amplifier output, inboard
of the counter-electrode isolation switch --- and read it either on
\texttt{atb\_d} with \texttt{AFE\_DBGDIREN}~=~1 or on-chip by routing the bus
into the converter with \texttt{AFE\_ADCSRCSEL}~=~4 and averaging. Repeat at
three reference codes across the working window: the intercept is the input
offset, \(\sigma = \SI{5.20}{\milli\volt}\)
(Table~\ref{tab:ampab-mc-offset}), and the slope is the finite-loop-gain
tracking error. Take the reading unbuffered --- a test buffer built from the
same devices carries an offset of the same class, so buffering substitutes one
unknown for another. For the transimpedance role, isolate the working electrode
and read slot 4 the same way; that reading is output-referred and carries
\(\sigma = \SI{8.91}{\milli\volt}\) (Table~\ref{tab:tiaab-mc-offset}), which
divided by the active gain is the channel's zero-current error. This test fabric
is design intent; it is not yet in the schematics.

% maestro2tex -- generated master include for Pixel.
%% Regenerate with tools/maestro2tex; do not edit by hand.
%%
%% Usage from the TRM (which builds with cwd = latex/TRM/):
%     % maestro2tex -- generated master include for Pixel.
%% Regenerate with tools/maestro2tex; do not edit by hand.
%%
%% Usage from the TRM (which builds with cwd = latex/TRM/):
%     % maestro2tex -- generated master include for Pixel.
%% Regenerate with tools/maestro2tex; do not edit by hand.
%%
%% Usage from the TRM (which builds with cwd = latex/TRM/):
%     \input{include/analog/Pixel.tex}
%
%% Needs pgfplots, booktabs, siunitx and placeins -- all already in
%% packages-commands.tex. \providecolor needs xcolor (also present).

\providecommand{\MaestroRoot}{include/analog/}

% Categorical palette: slots 1-5 of the validated reference palette,
% in documented order. Dash patterns are a secondary encoding for print.
\providecolor{mtxA}{HTML}{2A78D6}
\providecolor{mtxB}{HTML}{EB6834}
\providecolor{mtxC}{HTML}{1BAF7A}
\providecolor{mtxD}{HTML}{EDA100}
\providecolor{mtxE}{HTML}{E87BA4}

% Keep this block's floats inside this block (see placeins).
\FloatBarrier

\input{\MaestroRoot note_pixel_intro.tex}
\input{\MaestroRoot sch_tb_pixel.tex}
\input{\MaestroRoot tab_Pixel_corners.tex}
\input{\MaestroRoot note_pixel_accuracy.tex}
\input{\MaestroRoot fig_px_xover_ierr.tex}
\input{\MaestroRoot tab_px_acc35.tex}
\input{\MaestroRoot tab_px_hg1m_acc.tex}
\input{\MaestroRoot fig_px_hg1m_zoom.tex}
\input{\MaestroRoot note_pixel_hg.tex}
\input{\MaestroRoot tab_px_hg_demo.tex}
\input{\MaestroRoot note_pixel_ctrl.tex}
\input{\MaestroRoot fig_px_ctrl_loopgain.tex}
\input{\MaestroRoot tab_px_ctrl_stability.tex}
\input{\MaestroRoot tab_px_ctrl_tracking.tex}
\input{\MaestroRoot tab_px_ctrl_range.tex}
\input{\MaestroRoot tab_px_ctrl_drive.tex}
\input{\MaestroRoot note_pixel_grid.tex}
\input{\MaestroRoot tab_px_grid_pm.tex}
\input{\MaestroRoot note_pixel_noise.tex}
\input{\MaestroRoot tab_px_grid_noise.tex}
\input{\MaestroRoot note_pixel_transient.tex}
\input{\MaestroRoot fig_px_cv_voltammogram.tex}
\input{\MaestroRoot fig_px_cv_potential.tex}
\input{\MaestroRoot tab_px_cv_stats.tex}
\input{\MaestroRoot fig_px_step_i.tex}
\input{\MaestroRoot tab_px_step_stats.tex}
\input{\MaestroRoot note_pixel_mc.tex}
\input{\MaestroRoot fig_px_mc_izero_35k.tex}
\input{\MaestroRoot fig_px_mc_izero_560k.tex}
\input{\MaestroRoot tab_px_mc_izero.tex}
\input{\MaestroRoot tab_px_mc_izero_allgain.tex}
\input{\MaestroRoot tab_px_mc_rezero.tex}

\FloatBarrier

%
%% Needs pgfplots, booktabs, siunitx and placeins -- all already in
%% packages-commands.tex. \providecolor needs xcolor (also present).

\providecommand{\MaestroRoot}{include/analog/}

% Categorical palette: slots 1-5 of the validated reference palette,
% in documented order. Dash patterns are a secondary encoding for print.
\providecolor{mtxA}{HTML}{2A78D6}
\providecolor{mtxB}{HTML}{EB6834}
\providecolor{mtxC}{HTML}{1BAF7A}
\providecolor{mtxD}{HTML}{EDA100}
\providecolor{mtxE}{HTML}{E87BA4}

% Keep this block's floats inside this block (see placeins).
\FloatBarrier

\input{\MaestroRoot note_pixel_intro.tex}
\input{\MaestroRoot sch_tb_pixel.tex}
\input{\MaestroRoot tab_Pixel_corners.tex}
\input{\MaestroRoot note_pixel_accuracy.tex}
\input{\MaestroRoot fig_px_xover_ierr.tex}
\input{\MaestroRoot tab_px_acc35.tex}
\input{\MaestroRoot tab_px_hg1m_acc.tex}
\input{\MaestroRoot fig_px_hg1m_zoom.tex}
\input{\MaestroRoot note_pixel_hg.tex}
\input{\MaestroRoot tab_px_hg_demo.tex}
\input{\MaestroRoot note_pixel_ctrl.tex}
\input{\MaestroRoot fig_px_ctrl_loopgain.tex}
\input{\MaestroRoot tab_px_ctrl_stability.tex}
\input{\MaestroRoot tab_px_ctrl_tracking.tex}
\input{\MaestroRoot tab_px_ctrl_range.tex}
\input{\MaestroRoot tab_px_ctrl_drive.tex}
\input{\MaestroRoot note_pixel_grid.tex}
\input{\MaestroRoot tab_px_grid_pm.tex}
\input{\MaestroRoot note_pixel_noise.tex}
\input{\MaestroRoot tab_px_grid_noise.tex}
\input{\MaestroRoot note_pixel_transient.tex}
\input{\MaestroRoot fig_px_cv_voltammogram.tex}
\input{\MaestroRoot fig_px_cv_potential.tex}
\input{\MaestroRoot tab_px_cv_stats.tex}
\input{\MaestroRoot fig_px_step_i.tex}
\input{\MaestroRoot tab_px_step_stats.tex}
\input{\MaestroRoot note_pixel_mc.tex}
\input{\MaestroRoot fig_px_mc_izero_35k.tex}
\input{\MaestroRoot fig_px_mc_izero_560k.tex}
\input{\MaestroRoot tab_px_mc_izero.tex}
\input{\MaestroRoot tab_px_mc_izero_allgain.tex}
\input{\MaestroRoot tab_px_mc_rezero.tex}

\FloatBarrier

%
%% Needs pgfplots, booktabs, siunitx and placeins -- all already in
%% packages-commands.tex. \providecolor needs xcolor (also present).

\providecommand{\MaestroRoot}{include/analog/}

% Categorical palette: slots 1-5 of the validated reference palette,
% in documented order. Dash patterns are a secondary encoding for print.
\providecolor{mtxA}{HTML}{2A78D6}
\providecolor{mtxB}{HTML}{EB6834}
\providecolor{mtxC}{HTML}{1BAF7A}
\providecolor{mtxD}{HTML}{EDA100}
\providecolor{mtxE}{HTML}{E87BA4}

% Keep this block's floats inside this block (see placeins).
\FloatBarrier

\input{\MaestroRoot note_pixel_intro.tex}
\input{\MaestroRoot sch_tb_pixel.tex}
\input{\MaestroRoot tab_Pixel_corners.tex}
\input{\MaestroRoot note_pixel_accuracy.tex}
\input{\MaestroRoot fig_px_xover_ierr.tex}
\input{\MaestroRoot tab_px_acc35.tex}
\input{\MaestroRoot tab_px_hg1m_acc.tex}
\input{\MaestroRoot fig_px_hg1m_zoom.tex}
\input{\MaestroRoot note_pixel_hg.tex}
\input{\MaestroRoot tab_px_hg_demo.tex}
\input{\MaestroRoot note_pixel_ctrl.tex}
\input{\MaestroRoot fig_px_ctrl_loopgain.tex}
\input{\MaestroRoot tab_px_ctrl_stability.tex}
\input{\MaestroRoot tab_px_ctrl_tracking.tex}
\input{\MaestroRoot tab_px_ctrl_range.tex}
\input{\MaestroRoot tab_px_ctrl_drive.tex}
\input{\MaestroRoot note_pixel_grid.tex}
\input{\MaestroRoot tab_px_grid_pm.tex}
\input{\MaestroRoot note_pixel_noise.tex}
\input{\MaestroRoot tab_px_grid_noise.tex}
\input{\MaestroRoot note_pixel_transient.tex}
\input{\MaestroRoot fig_px_cv_voltammogram.tex}
\input{\MaestroRoot fig_px_cv_potential.tex}
\input{\MaestroRoot tab_px_cv_stats.tex}
\input{\MaestroRoot fig_px_step_i.tex}
\input{\MaestroRoot tab_px_step_stats.tex}
\input{\MaestroRoot note_pixel_mc.tex}
\input{\MaestroRoot fig_px_mc_izero_35k.tex}
\input{\MaestroRoot fig_px_mc_izero_560k.tex}
\input{\MaestroRoot tab_px_mc_izero.tex}
\input{\MaestroRoot tab_px_mc_izero_allgain.tex}
\input{\MaestroRoot tab_px_mc_rezero.tex}

\FloatBarrier

% HAND-WRITTEN fragment -- NOT generated by maestro2tex. Input from
% AnalogChapter.tex immediately after Pixel.tex, which is generated.
%
% Salvaged from the retired note_calibration_seq.tex step 5 (current zero) and
% the triage table of ~/vestarv/docs/afe_rev2_debug_mux.html §9/§9.1:
% per-channel slot 5 = ce_amp, slot 6 = re_buf (buffered at source), slot 7 =
% we_sum behind a series guard T-gate, slot 8 = we_pad. The discriminating pair
% is CE against WE; both faults read as a stuck code at the converter.
% Numbers: 3 sigma of A1 offset = 15.6 mV (tab_ampab_mc_offset); 0.824 nA/LSB at
% the top gain code (ss:tiarprog); zero-current sigma from sss:pixel-mc.
% Library cell names are kept out of the rendered prose.

\subsubsection{Measuring this block on chip}

Two nodes separate the four faults that all read as a stuck converter code.
Select per-channel slot 5 and confirm the counter-electrode drive sits mid-rail
and follows the reference code; then slot 6, the reference-electrode buffer
output, which must equal the commanded reference within \SI{16}{\milli\volt}
(\(3\sigma\) of A1's offset). A railed counter electrode with the summing node
at the common mode is a control-loop or electrode fault; an unrailed counter
electrode with the summing node off the common mode is a transimpedance fault;
the converter reports a rail for both, which is why this pair and not the
converter is the triage. Read the summing node itself at slot 7, through its
series guard and with the mid node clamped --- it is a virtual ground at
\SI{0.824}{\nano\ampere} per LSB at the top gain code, so an ungranted site must
stay doubly isolated from it. For the zero-current constant, open the
working-electrode isolation switch so the cell current is zero by construction,
read the transimpedance output at every gain code the channel will use, and
subtract the converter's zero code; the uncalibrated spread is \(\sigma =
12.1\)--\SI{167.3}{LSB} over the four taps (Section~\ref{sss:pixel-mc}). This
test fabric is design intent; it is not yet in the schematics.


% Assembled channel, hand-written throughout (no maestro2tex config; the numbers come
% from the pre-tapeout verification of the assembled cell,
% ~/vestarv/docs/afe_rev2_pixel_verify.html, and its runs under
% ~/chips/castalia/simruns/). It belongs immediately after the pixel section: the same
% channel with every ideal source replaced by silicon, read as a diff against it.
% The characterisation is re-baselined against ~/chips/castalia/simruns/rebase_20260824
% (dc, a 24-point stability grid, noise, power, five corners, a 200-sample Monte Carlo
% and a multiplexer-disturbance study) after the cell gained two unity-gain multiplexer
% buffers, lost both reference-DAC supply regulators, got real MIM devices for A1's
% compensation pair and moved four amplifier devices onto mismatch-enabled wrappers.
% Three runs were not repeated on the new OA and are flagged in the fragments that use
% them: the staircase voltammetry run, the five-corner stability run, and the Monte
% Carlo attribution/ablation and stability runs.
%
% The section is commented out, not deleted, to keep the chapter short. Every fragment
% below is on disk with its own provenance header, and restoring the section is
% uncommenting these fourteen \input lines plus the "Assembled channel" and "Input
% multiplexer" rows of the coverage table, the multiplexer clause of the pinned-passive
% paragraph, and the labels ss:pixtop, sss:pixtop-*, tab:pixtop-* and fig:pixtop-*.
% Four sentences elsewhere no longer reference those labels: the extracted-netlist
% sentence and the blocks-documented paragraph above, note_adc_intro's Coverage
% paragraph (pointing at ss:eis, which splices the same extracted macro), and
% note_calibration's Residual paragraph.
%
% \FloatBarrier
% % HAND-WRITTEN fragment -- NOT generated by maestro2tex. Input from
% AnalogChapter.tex immediately after note_meas_pixel.tex, i.e. after the flat
% pixel bench this section is contrasted against.
%
% Source of every statement: ~/vestarv/docs/afe_rev2_pixel_verify.html
% (revision of 2026-08-23, post owner fix session) --
%   pin list and internal inventory        S2, S3
%   bench and Maestro setup                S5, S6 (3 tests px_dc/px_zero/px_gain,
%                                          nominal, 37 C, section castalia_typ_25)
%   flat-bench-vs-assembled difference list S10
% Runs: ~/chips/castalia/simruns/pixtop2_{dc,zero,gain,cv,adcx_*}, 2026-08-23.
%
% 2026-08-24 RE-BASELINE, ~/chips/castalia/simruns/rebase_20260824/README.md.
% The owner edited anatop_pixel/schematic at 03:53 and anatop_tia_amp_ab at
% 04:04; the whole characterisation was re-run on the new OA.  Inventory changes
% reflected in the bullets below (S1.1 instance census, 17 instances):
%   + BUF_ADC   anatop_tia_amp_ab, unity-gain, ADC_MXY -> ADC_IN
%   + BUF_ATP   anatop_tia_amp_ab, unity-gain, ATPY -> ATP
%     -- so anatop_tia_amp_ab is now instantiated FOUR times per channel
%        (CONTROLAMP = A1, TIA.TIAAMP = A2, BUF_ADC, BUF_ATP)
%   + C0-C3     mimcap_sin lt=wt=15.58u ~ 492 fF, buffer compensation
%   ~ C_CMP1_CTRL / C_CMP2_CTRL now mimcap_sin (492 fF) instead of an ideal
%     capacitor c=cc_ab (500 fF)
%   - DAC_VP_PSU, DAC_VCM_PSU and the whole anatop_pixel_dac_psu subckt DELETED,
%     with anatop_inv_sm I1 -- F-4 is closed, and both ladders now run from the
%     2.5 V analog rail directly (which they effectively already did)
%   ~ M48/M49 (nch_25->nch_25_mac) and M57/M58 (pch_25->pch_25_mac) inside
%     anatop_tia_amp_ab now carry mismatchflag=1
% Every number in this subsection now comes from rebase_20260824.
% Library cell names are deliberately kept OUT of the rendered prose (owner
% request 2026-08-21); they live in these comments instead. The cell is
% castalia/anatop_pixel, bench castalia/tb_anatop_pixel_top view PixelTopTB.

\subsection{Assembled Channel}\label{ss:pixtop}

This is the cell intended for tapeout: one complete analog channel with every
ideal source of Section~\ref{ss:pixel} replaced by the silicon that will drive
it on die. The topology and the sign convention are unchanged --- A1 on the
counter electrode, A2 holding the working electrode at the common mode through
\(R_f\) --- so this section reports only what the assembly changes.

\begin{itemize}[nosep]
  \item The potential program and the common mode come from two 12-bit R-2R
    DACs (Section~\ref{ss:dacr2r12}), one per channel, at
    \SI{610.35}{\micro\volt} per code. Both ladders run from the
    \SI{2.5}{\volt} analog rail directly; the cell carries no reference-DAC
    supply regulator.
  \item \(R_f\) is the 320-code programmable ladder of
    Section~\ref{ss:tiarprog}, not an ideal resistor, and carries a
    \SI{22.0}{\pico\farad} compensation capacitor across it, without which the
    assembled loop does not close (Section~\ref{sss:pixtop-stability}).
  \item The 10-bit converter of Section~\ref{ss:adc} closes the channel behind
    a 16:1 input multiplexer whose slots select the transimpedance output, the
    common-mode rail, the commanded potential and the analog test point. A
    unity-gain buffer stands between that multiplexer and the converter input,
    and a second buffers the test-port multiplexer output, so the class-AB
    amplifier of Section~\ref{ss:ampab} appears four times per channel: A1, A2
    and the two followers. Each buffer carries its own
    \SI{492}{\femto\farad} metal--insulator--metal compensation pair, as does
    A1.
  \item Bias comes from the local generator of Section~\ref{ss:biasgen-local}
    at trim code 37.
\end{itemize}

The counter, reference and working electrodes leave the cell directly: there
are no isolation switches, no reference-electrode buffer and no on-chip gain
reference resistor, so the fabric several of the calibration steps of
Section~\ref{ss:calibration} assume is not yet present.

All data below is schematic-level, at the nominal process point unless a corner
is named, \SI{37}{\celsius} and a \SI{2.5}{\volt} analog supply. It comes from
three dc and transient tests --- a sweep of the potential code, a transient with
the electrode inert, and a sweep of all 320 gain codes --- a netlist-level
staircase voltammetry run (Section~\ref{sss:pixtop-cv}), and loop-gain, noise,
transient-noise, supply-current and 200-sample Monte Carlo runs on the netlist
those tests produce
(Sections~\ref{sss:pixtop-stability}--\ref{sss:pixtop-mc}).
Section~\ref{sss:pixtop-adc} is the exception and is stated there: the
converter's capacitor array exists only in layout, so every converter number
comes from a netlist in which the extracted converter replaces the schematic
one.

% %% HAND-WRITTEN table fragment (maestro2tex house style; nothing here is
%% generated). Every value transcribed from
%% ~/vestarv/docs/afe_rev2_pixel_verify.html, revision 2026-08-23:
%%   rows 1-4   S7.1 (run simruns/pixtop2_dc)
%%   rows 5-8   S8.2 (run simruns/pixtop2_gain), S7.3
%%   rows 9-12  S9.1-S9.3 (run simruns/pixtop2_cv, cv_pixtop.dat)
%%   rows 13-14 S7.5 (runs simruns/pixtop2_adcx_*, extracted converter)
%% Flat-bench comparison values (0.814 mV, 8.6658 uA, 6.0553 uA) are the same
%% document's S10 and S9.2 columns, i.e. tb_anatop_pixel/CvBvTB test ps_cv.
%%
%% 2026-08-23, post-C_F: rows 1-14 are UNCHANGED -- simruns/postfix_corner/
%% corner_table.txt reproduces every dc value at tt to the printed digits, and
%% the capacitor moves only the transient zero, in the 4th digit (1.6 uV,
%% attributed in simruns/postfix_cfattrib).  Four rows ADDED at the foot:
%%   PM 75.71 deg / 1.068 dB worst of 24, 72.37 deg at ss
%%       -- simruns/postfix_stb/stb_table.txt, postfix_corner_stb/
%%   f3dB 85.77 kHz (N=0) / 3415 Hz (N=319) at Cdl 100 nF, Rct 1 M -- same table
%%   sigma_tread 1.475-2.889 nA = 0.036-1.781 LSB over 16 codes
%%       -- simruns/trmnoise_gain/noise_vs_gain_table.txt
%%   1.6424 mW, of which 806.758 uW is the two DAC supply regulators
%%       -- simruns/postfix_pwr/power_table.txt (extracted converter)
%% The LSB row's conditions cell now states that 88.5 nA is the NOMINAL-ladder
%% figure; the measured in-loop Zt(dc) of 19.63 kOhm gives 81.2 nA, which is what
%% sss:pixtop-noise and fig:pixtop-noise use.
%%
%% 2026-08-24: two Monte Carlo rows ADDED at the foot.  A third row,
%% "Per-channel yield, calibrated 61.5 % / 0.50 % uncalibrated", was written and
%% then REMOVED the same day on the owner's instruction "do not talk about pixel
%% level yield yet"; yield_table.txt still holds it.  The tracking rows'
%% conditions cells say which of the two ranges is the specification (the
%% 0.4-2.1 V window, owner decision 2026-08-24, sss:pixtop-mc).
%%
%% 2026-08-24 RE-BASELINE against simruns/rebase_20260824 (the cell was edited
%% at 03:53/04:04 and everything re-measured).  Changed here:
%%   power  1.6424 -> 1.3228 mW  (pwr/power_table.txt, extracted converter:
%%          528.7005 uA on 2.5 V + 1.0391 uA on 1.0 V).  Attribution closes to
%%          0.05 uW: -806.758 uW of deleted DAC PSU regulators, +243.776 uW
%%          BUF_ADC, +243.347 uW BUF_ATP.  Corner 2.5 V rail 432.0 (ss) to
%%          913.6 (ff) uA, was 517.1 / 1075.2.
%%   MC     sigma 5.754 -> 5.488 mV (tia_offset) and 5.552 -> 5.287 mV
%%          (re_track), mc_sigma_table.txt, 200 samples, 1264 random variables.
%%          v_adcin row ADDED: mean 1.24815 V, sigma 8.040 mV.
%%   ADDED  converter-input zero row (dc_table.txt: v_adcin 1.2484987 V, i.e.
%%          -619.6 uV = BUF_ADC's own offset below V_OUT) and the calibration-
%%          slot disturbance row (mux/mux_disturbance_table.txt slot 1 NEW).
%%   CAPTION: the "regulators instantiated but drive nothing" clause is gone --
%%          they are deleted from the cell.  The tracking figures are still
%%          optimistic, for the same reason stated the other way round.
%% Every dc row above is UNCHANGED and re-verified: rebase_20260824/dc_table.txt
%% reproduces all of them to the printed digits (worst delta 3.0 nV), and
%% corner/corner_table_raw.txt still passes 20/20 at tt/ff/fs/sf/ss.  Rows 15-16
%% (stability, noise) are unchanged; stb/stb_table.txt gives the same 75.71 deg
%% worst of 24 and noise/noise_table.txt the same tread-averaged 1.475-2.889 nA.
%% The 72.37 deg slow-corner figure is NOT re-measured on this OA -- the
%% five-corner stability run predates the edit; the nominal grid moved +0.36 deg.
\begin{table}[!htb]
  \centering
  \caption[{Headline measurements on the assembled channel.}]{Headline
  measurements on the assembled channel, nominal process point,
  \SI{37}{\celsius}, \SI{2.5}{\volt} supply. Rows 1--9 are dc; rows 10--13 come
  from the staircase voltammetry run of Figure~\ref{fig:pixtop-cv}; rows 14--16
  and the power figure need the extracted converter, everything else in the
  channel staying at schematic level; rows 17--20 are the small-signal,
  noise and power runs of
  Sections~\ref{sss:pixtop-stability}--\ref{sss:pixtop-power}; the final three
  are Monte Carlo, and are narrower than the built channel's for the reason
  Section~\ref{sss:pixtop-open} gives. The two tracking
  figures are optimistic by construction: both ladders run from the
  \SI{2.5}{\volt} rail, so the \SI{25.8}{\milli\volt} of code-dependent
  reference movement Section~\ref{sss:dacr2r12-supply} measures on the supply
  block is not exercised here.}
  \label{tab:pixtop-summary}
  {\footnotesize
  \setlength{\tabcolsep}{5pt}
  \begin{tabular}{@{}%
    >{\raggedright\arraybackslash}p{0.315\textwidth}%
    >{\raggedleft\arraybackslash}p{0.185\textwidth}%
    >{\raggedright\arraybackslash}p{0.435\textwidth}@{}}
    \toprule
    Quantity & Value & Conditions \\
    \midrule
    Reference tracking error &
      \SI{0.672}{\milli\volt} &
      dc sweep over the specified \SIrange{0.4}{2.1}{\volt} window
      (\S\ref{sss:pixtop-mc}); \SI{0.814}{\milli\volt} with
      an ideal source (\S\ref{ss:pixel}) \\
    Reference tracking error, full range &
      \SI{1.419}{\milli\volt} &
      \SIrange{0}{2.42}{\volt}, all 4096 codes; reported, not specified \\
    Reference DAC step &
      \SI{610.35}{\micro\volt} &
      \(\SI{2.5}{\volt}/4096\) from an ideal rail; \SI{-181}{\micro\volt}
      endpoint error at code 3968 \\
    Common-mode output &
      \SI{1.249738}{\volt} &
      mid-code, \SI{-262}{\micro\volt} from \SI{1.25}{\volt} \\
    \addlinespace
    Zero-current output, \(v_{\mathrm{out}} - V_{\mathrm{CM}}\) &
      \SI{-619.64}{\micro\volt} &
      inert electrode, \(N = 0\) \\
    Spread of that zero over gain codes &
      \SI{0.3}{\nano\volt} &
      all 320 codes \\
    A2 input bias current &
      \SI{2.04}{\pico\ampere} &
      \SI{3.91}{\micro\volt} of drift over \(\Delta R_f = \SI{1.914}{\mega\ohm}\) \\
    Uncalibrated current zero &
      \(34.4\,/\,1.57\,/\,\SI{0.323}{\nano\ampere}\) &
      \(N = 0\,/\,63\,/\,319\) \\
    Converter-input zero &
      \SI{1.248499}{\volt} &
      inert electrode, readout slot; \SI{-619.6}{\micro\volt} from the
      transimpedance output, which is the input buffer's own offset \\
    \addlinespace
    Feedback resistance, in situ &
      \SI{37.16}{\kilo\ohm} &
      \(N = 3\), \SI{37}{\celsius}; median, p10--p90
      \SIrange{36.95}{37.44}{\kilo\ohm} \\
    Readout against the in-line probe &
      \(\le \SI{39.4}{\nano\ampere}\) &
      331 treads above \SI{100}{\nano\ampere}, on \SI{7.7}{\micro\ampere}
      peaks, i.e. \SI{0.5}{\percent} \\
    Voltammetric peak potential shift &
      \(0.3\,/\,\SI{2.3}{\milli\volt}\) &
      forward\,/\,reverse, against the ideal-source run of
      \S\ref{ss:analog-electrode-bc} \\
    Sampled peak height &
      \(-11.2\,/\,{-}13.9\,\%\) &
      forward\,/\,reverse, same reference; staircase sampling, not an error \\
    \addlinespace
    Converter agreement with its own bench &
      \(\le \SI{3.3}{LSB}\) &
      extracted converter, \SIrange{0.62}{1.87}{\volt} \\
    Current per converter LSB &
      \(88.5\,/\,\SI{0.825}{\nano\ampere}\) &
      \(N = 0\,/\,319\); \SI{1.593}{\milli\volt} LSB over the measured
      \SI{1.63}{\volt} span, on the nominal ladder. The measured in-loop
      \(R_f\) puts the bottom code at \SI{81.2}{\nano\ampere} \\
    Calibration-slot disturbance &
      \SI{1.75}{\milli\volt} &
      peak on the transimpedance output over a \SI{21}{\micro\second} record
      with the common-mode slot selected; settles in
      \SI{0.54}{\micro\second}, residual below \SI{1}{\nano\volt} \\
    \addlinespace
    Transimpedance phase margin &
      \SI{75.7}{\degree} &
      worst of 24 electrode\(\,\times\,\)gain points, \SI{1.07}{\decibel}
      peaking; both buffer loops \(\ge \SI{78.4}{\degree}\) \\
    Closed-loop bandwidth &
      \(85.8\,/\,\SI{3.4}{\kilo\hertz}\) &
      \(N = 0\,/\,319\), \(C_{\mathrm{dl}} = \SI{100}{\nano\farad}\), with the
      \SI{22.0}{\pico\farad} compensation capacitor \\
    Input-referred noise, one tread &
      \SIrange{1.5}{2.9}{\nano\ampere} &
      \SI{2}{\milli\second} average; \(0.036\)--\(1.78\,\)LSB as the gain code
      rises \\
    Channel power &
      \SI{1.323}{\milli\watt} &
      extracted converter; \SI{487.1}{\micro\watt} of it is the two
      unity-gain buffers \\
    \addlinespace
    Zero-current output, \(\sigma\) &
      \SI{5.488}{\milli\volt} &
      200 Monte Carlo samples, process and mismatch
      (\S\ref{sss:pixtop-mc}); \(\pm\SI{2}{\milli\volt}\) budget \\
    Reference tracking error, \(\sigma\) &
      \SI{5.287}{\milli\volt} &
      same run, mid-code; \SI{5.72}{\milli\volt} on the flat bench
      (\S\ref{sss:pixel-mc}) \\
    Converter-input zero, \(\sigma\) &
      \SI{8.040}{\milli\volt} &
      same run; the transimpedance offset and the buffer's own, uncorrelated.
      It cancels in the zero subtraction (\S\ref{sss:pixtop-mc}) \\
    \bottomrule
  \end{tabular}}
\end{table}

% % HAND-WRITTEN fragment -- NOT generated by maestro2tex.
% Source: ~/vestarv/docs/afe_rev2_pixel_verify.html rev. 2026-08-23,
%   S7.1 and S8.1 (run simruns/pixtop2_dc, dc sweep of vp_code 0..3968 step 128,
%   ncode = 0, vcm_code = 2048, adc_sel = 0, nominal, 37 C):
%     re_err 0.6724 mV over the 0.4-2.1 V window, 1.4194 mV over 0-2.42 V
%     (+0.402 mV at code 0, -1.419 mV at code 3968); v_pattern 104.09 uV at
%     code 0 and 2.421694 V at code 3968 (ideal 2.421875 V, -181 uV);
%     v_cm_dc 1.249738 V; A2 offset -619.64 uV from S8.2.
%   S4 F-4 for the unconnected DAC supply regulators and the 25.8 mV figure.
% The 528.96 uV/LSB and 2.166 V full scale quoted for comparison are the
% reference-DAC section's own numbers (tab_dacr2r12_mc_reference), measured on
% the shunt-clamp regulator that is absent from this cell's supply path.

\subsubsection{Potential Axis}\label{sss:pixtop-potential}

The reference electrode follows the commanded DAC output to
\SI{0.672}{\milli\volt} over the \SIrange{0.4}{2.1}{\volt} window
Section~\ref{ss:pixel} uses, and to \SI{1.419}{\milli\volt} over the full
\SIrange{0}{2.42}{\volt} code range --- \(+\SI{0.402}{\milli\volt}\) at code 0,
\SI{-1.419}{\milli\volt} at code 3968. The ideal-source measurement is
\SI{0.814}{\milli\volt} over the same window. Two terms make it up, both
static: the ladder's own error on a \SI{610.35}{\micro\volt} step, which is
\SI{-181}{\micro\volt} at the top code, and A1's \SI{-620}{\micro\volt} input
offset, which is constant across the range. A two-point fit per channel over
the potential code axis absorbs the pair and leaves a residual well inside the
\SI{2}{\milli\volt} budget of Section~\ref{ss:ampab}; the offset-and-slope
constant of Table~\ref{tab:calibration-constants} is what stores it. The
common-mode DAC sits \SI{262}{\micro\volt} below \SI{1.25}{\volt} at mid-code
and every current referred to that node inherits the error.

The step is \SI{610.35}{\micro\volt}, \(\SI{2.5}{\volt}/4096\), not the
\SI{528.96}{\micro\volt} of Section~\ref{ss:dacr2r12}, because in this cell
both ladders run from the \SI{2.5}{\volt} rail directly: the DAC supply
regulators are instantiated but their outputs reach nothing. The
\SI{25.8}{\milli\volt} of code-dependent reference movement measured in
Section~\ref{sss:dacr2r12-supply} is therefore absent from every potential
number here. It is \(40\times\) the tracking error above and will dominate the
potential axis the moment the regulators are placed in the path --- the reason
to read \SI{0.672}{\milli\volt} as a missing stimulus rather than a result.

% % HAND-WRITTEN fragment -- NOT generated by maestro2tex.
% Source: ~/vestarv/docs/afe_rev2_pixel_verify.html rev. 2026-08-23, S7.3 and
% S8.2 (run simruns/pixtop2_gain, dc sweep ncode 0..319 step 1, inert electrode
% cox = cred = 0, nominal, 37 C):
%   V(WE) - V_CM = -619.6424 uV at N = 0 and -619.6421 uV at N = 319, i.e.
%     0.3 nV of spread over all 320 codes;
%   V_OUT - V_CM = -619.6444 uV (N = 0) ... -623.5525 uV (N = 319), a 3.908 uV
%     drift over dRf = 1.914 MOhm -> 2.042 pA of A2 input current;
%   i_off = -34.4247 nA (N = 0), -1.5651 nA (N = 63), -0.3227 nA (N = 319).
% The 2 nA input-referred noise budget is the in_rma spec used throughout this
% tree (see ss:tiarprog).

\subsubsection{Zero-Current Axis}\label{sss:pixtop-zero}

With the electrode inert the working electrode sits \SI{619.64}{\micro\volt}
below the common mode, and it is the same \SI{619.64}{\micro\volt} at every one
of the 320 gain codes: \SI{0.3}{\nano\volt} of spread. No current flows in
\(R_f\), so the output carries A2's input offset alone and nothing that scales
with the tap. \textbf{One zero measurement therefore calibrates every gain
code} --- the result the per-gain-code current zero of
Table~\ref{tab:calibration-constants} rests on, and the reason the common-mode
slot on the converter multiplexer is worth its silicon.

What survives that calibration is input bias current. The output drifts a
further \SI{3.91}{\micro\volt} between \(N = 0\) and \(N = 319\) while \(R_f\)
grows by \SI{1.914}{\mega\ohm}: \SI{2.04}{\pico\ampere} into A2's inverting
input, \SI{0.5}{\percent} of the \SI{0.41}{\micro\ampere} full scale at maximum
gain and negligible at every code below it.

Referred to current, the uncalibrated zero is \SI{34.4}{\nano\ampere} at
\(N = 0\), \SI{1.57}{\nano\ampere} at \(N = 63\) and
\SI{0.323}{\nano\ampere} at \(N = 319\). The last two are already inside the
\SI{2}{\nano\ampere} input-referred noise budget, so at high gain the
uncalibrated offset is not what limits a reading; at the bottom code it is
\SI{34.4}{\nano\ampere} on a \SI{40.7}{\micro\ampere} full scale and below one
converter code either way (Section~\ref{sss:pixtop-adc}).

% % HAND-WRITTEN fragment -- NOT generated by maestro2tex.
% Source: ~/vestarv/docs/afe_rev2_pixel_verify.html rev. 2026-08-23, S9
% (run simruns/pixtop2_cv, 1.007 s transient, 175 205 stored points, 0 errors):
%   programme  triangle staircase -0.5 -> +0.5 -> -0.5 V vs vcm, 4 mV treads
%              (6-7 DAC codes), 2 ms/tread = 2 V/s, 20 ns bus edges, one cycle,
%              501 treads, codes 1229..2867, ncode = 3;
%   electrode  cdl_we 200 nF, r_u = r_ce = 1 kOhm, e0 0.3 V, k0 0.1, Cox 1e-6,
%              Cred 0, D 5e-6, 37 C -- the same parameter set as the
%              ideal-source run of note_electrode_afe.tex (8.666 uA peak);
%   S9.1  Rf_eff = 37 161.9 ohm median over 124 921 points, p10 36 953.7,
%         p90 37 444.1; nominal law 36 000 ohm; resistor bench 37 620 ohm at 25 C;
%   S9.2  peaks: forward -7.6970 uA @ +0.3436 V (flat bench -8.6658 @ +0.3433),
%         reverse +5.2121 uA @ +0.2551 V (flat bench +6.0553 @ +0.2574);
%         raw (unsampled) peaks are +26 % and +42 % on the same waveform;
%   S9.3  |i_readout - I(IPRB_F)| <= 39.4 nA over 331 treads above 100 nA;
%         Vos trap: the extractor's own pre-scan estimate gives +0.3805 mV.
% Netlist-level in one respect: the 12-bit potential bus is driven by per-bit
% PWL sources so the programme is a real code sequence, and the SAR clock is
% held static.
% 2026-08-24 RE-BASELINE: this 1 s transient was NOT repeated on the edited cell
%   and nothing in simruns/rebase_20260824 moves it.  The two new buffers are
%   outside both loops; deleting the DAC supply regulators left every dc value
%   bit-identical; and stb/stb_table.txt reproduces Zt(1 Hz), peaking and f3dB
%   at every printed digit, so Rf_eff, the peak positions and the readout
%   residual are unaffected.  Re-run it if the loop is ever recompensated.

\subsubsection{Cyclic Voltammetry Through the Complete Chain}\label{sss:pixtop-cv}

Figure~\ref{fig:pixtop-cv} is the electrode of Section~\ref{ss:analog-electrode-bc}
measured through the whole assembled channel: the potential program is a real
12-bit code sequence, the applied potential is servo'd by A1, and the current
is read back through the programmable resistor at gain code~3. The programme is
a triangle staircase, \(\pm\SI{0.5}{\volt}\) about the common mode in
\SI{4}{\milli\volt} treads of 6 or 7 codes, \SI{2}{\milli\second} per tread ---
\SI{2}{\volt\per\second} --- 501 treads over one cycle, sampled once per tread
\SI{50}{\micro\second} before the next code change.

\paragraph{Potential.} The peaks land within \SI{0.3}{\milli\volt} (forward) and
\SI{2.3}{\milli\volt} (reverse) of the same electrode driven by an ideal source.
The potential axis therefore survives \SI{610}{\micro\volt} quantisation, the
ladder's own error, the common-mode DAC's \SI{-262}{\micro\volt} and A1's
\SI{-620}{\micro\volt} offset and still puts the redox features where an ideal
source puts them --- which is what identifies an analyte.

\paragraph{Current.} Sampled peak heights are \SI{11.2}{\percent} and
\SI{13.9}{\percent} below the ideal-source run. A tread sampled at its end has
shed the double-layer charging current that a continuous ramp carries, and the
diffusion layer relaxes over each \SI{2}{\milli\second} hold, so the two
numbers measure different things about the same electrode rather than
disagreeing: kept raw, the same waveform peaks \SI{26}{\percent} and
\SI{42}{\percent} \emph{above} the ideal-source run on the charging tails.

\paragraph{Feedback resistance, measured in the loop.} A current probe in the
feedback branch makes \(R_f\) a measurement rather than an assumption:
\((v_{\mathrm{out}} - v_{\mathrm{we}})/i_f = \SI{37.16}{\kilo\ohm}\) median over
124\,921 points, p10--p90 \SIrange{36.95}{37.44}{\kilo\ohm}, against
\SI{36.0}{\kilo\ohm} from the nominal law and \SI{37.62}{\kilo\ohm} measured
two-terminal at \SI{25}{\celsius} (Section~\ref{ss:tiarprog}) ---
\SI{1.2}{\percent} below the latter at \SI{37}{\celsius}.

\paragraph{Readout.} Over the 331 treads carrying more than
\SI{100}{\nano\ampere} the voltage readout
\((v_{\mathrm{out}} - V_{\mathrm{CM}} - V_{\mathrm{os}})/R_f\) --- the
arithmetic firmware performs --- reproduces the probe current to within
\SI{39.4}{\nano\ampere} on \SI{7.7}{\micro\ampere} peaks, \SI{0.5}{\percent},
inside the \(\pm\SI{0.65}{\percent}\) spread of \(R_f\) itself. No systematic
term is left over. \(V_{\mathrm{os}}\) has to be supplied from a genuine
zero-current measurement (Section~\ref{sss:pixtop-zero}): estimated instead from
the pre-scan quiet period, where this electrode already passes
\SI{27}{\nano\ampere}, it comes out \(+\SI{0.38}{\milli\volt}\) and silently
subtracts a standing current from the whole voltammogram.

One cycle, one gain code, one scan rate, at nominal, with the converter clock
held static.

% %% HAND-WRITTEN fragment -- NOT generated by maestro2tex (XY figure kind; the
%% engine has no signal-vs-signal plot). Precedent: fig_electrode_panel.tex.
%%
%% Data: data/pixtop_cv_{fwd,rev}.dat, written by
%% ~/chips/castalia/skill/pixtop_cv_fig.py from
%% simruns/pixtop2_cv/cv_pixtop.dat (the sampled staircase voltammogram
%% extracted by skill/pixtop_cv_extract.py --vos -619.644u from the full-chain
%% run of 2026-08-23; one row per tread, 503 rows, split at the potential apex).
%% Columns: E = V(RE) - V(V_CM) [V], i = (V_OUT - V_CM - Vos)/Rf_eff [uA] with
%% Rf_eff = 37 161.9 ohm measured in the same run. Sign convention is the
%% measured one of ss:pixel: oxidation reads negative.
%% Re-run the extractor and the script to regenerate.
\providecommand{\MaestroRoot}{include/analog/}
\providecolor{mtxA}{HTML}{2A78D6}
\providecolor{mtxB}{HTML}{EB6834}
\begin{figure}[htbp]
  \centering
  {\pgfplotsset{compat=1.18}%
  \begin{tikzpicture}
    \begin{axis}[
      width=0.72\linewidth, height=6.2cm,
      grid=both,
      major grid style={line width=0.2pt, draw=black!14},
      minor grid style={line width=0.2pt, draw=black!7},
      minor tick num=1,
      axis line style={line width=0.4pt, draw=black!45},
      tick style={line width=0.4pt, draw=black!45},
      tick align=outside,
      label style={font=\small}, tick label style={font=\small},
      legend style={font=\footnotesize, draw=none, fill=white,
                    fill opacity=0.85, text opacity=1,
                    at={(0.03,0.03)}, anchor=south west, row sep=1pt},
      every axis plot/.append style={line join=round, no marks},
      xlabel={\(E = V_{\mathrm{RE}} - V_{\mathrm{CM}}\)~[\si{\volt}]},
      ylabel={\(i = (v_{\mathrm{out}} - V_{\mathrm{CM}})/R_f\)~[\si{\micro\ampere}]},
      enlarge x limits=0.03, enlarge y limits=0.10,
    ]
      \addplot[mtxA, line width=1.0pt] table {\MaestroRoot data/pixtop_cv_fwd.dat};
      \addlegendentry{forward, \SIrange{-0.5}{0.5}{\volt}}
      \addplot[mtxB, line width=1.0pt, dashed] table {\MaestroRoot data/pixtop_cv_rev.dat};
      \addlegendentry{reverse}
      \addplot[black!45, line width=0.4pt, domain=-0.52:0.52, samples=2,
        forget plot] {0};
    \end{axis}
  \end{tikzpicture}}
  \caption[{Staircase voltammogram through the assembled channel.}]{Staircase
  cyclic voltammogram taken through the complete assembled channel --- real
  12-bit potential codes, real programmable feedback resistor at
  \(N = 3\) (\SI{37.16}{\kilo\ohm} measured in the loop) --- of the electrode
  of Section~\ref{ss:analog-electrode-bc} at \SI{2}{\volt\per\second},
  \SI{37}{\celsius}, nominal process point. One point per
  \SI{4}{\milli\volt} tread, sampled \SI{50}{\micro\second} before the next
  code change, so the double-layer charging current has decayed and what
  remains is faradaic; the sign convention is the measured one of
  Section~\ref{ss:pixel}, \(\mathrm{d}v_{\mathrm{out}}/\mathrm{d}i_{\mathrm{we}}
  = +R_f\), so oxidation reads negative. Peak potentials sit within
  \SI{0.3}{\milli\volt} and \SI{2.3}{\milli\volt} of the same electrode driven
  by an ideal source; sampled peak heights are \SI{11.2}{\percent} and
  \SI{13.9}{\percent} lower, which is the staircase sampling and not a channel
  error (Table~\ref{tab:pixtop-summary}).}
  \label{fig:pixtop-cv}
\end{figure}

% % HAND-WRITTEN fragment -- NOT generated by maestro2tex.
% Source: ~/vestarv/docs/afe_rev2_pixel_verify.html rev. 2026-08-23, S7.4/S7.5
% (runs simruns/pixtop2_adcx_*, 21 us transients on the maestro netlist with the
% schematic converter subckt replaced by the flattened av_extracted one from the
% standalone ADC bench; 19 ports matched by name and order, 4451 device lines;
% decoded by skill/pixtop_adc_decode.py):
%   agreement with the standalone bench <= 3.3 LSB over 0.62-1.87 V, three
%     conversions per run within 2 LSB;
%   1 LSB = 1.5934 mV over the measured 1.63 V span -> 88.5 nA at N = 0,
%     42.9 nA at N = 3, 4.02 nA at N = 63, 0.825 nA at N = 319;
%   at N = 0 the zero-current point reads code_msb 507 = 1.2468 V against a true
%     1.249137 V: -1.5 LSB, i.e. -130 nA of apparent current on a -34 nA signal;
%   bit-9 inversion: the raw bus is code = 512*b9 + ... + b0, and every
%     published number (the transfer table of ss:adc, SIM_STATUS 2.4) is
%     code_msb = code + 512 - 1024*b9. The standalone bench's own output set
%     already carried both definitions; no document stated the consequence.
% At schematic level the CDAC subckt has zero devices, so the code is stuck at
% 1022 and RDY never rails (v_rdy_p 0.536 V) -- deliberately not specced.

\subsubsection{Converter Readout}\label{sss:pixtop-adc}

The converter's charge-redistribution array exists only in layout, so at
schematic level it makes no bit decisions and the code the assembled channel
reports is not a conversion. The numbers below come from a netlist in which the
extracted converter replaces the schematic one and nothing else changes.

Against known voltages the channel's transfer matches the standalone
characterisation of Section~\ref{ss:adc} to within \SI{3.3}{LSB} across
\SIrange{0.62}{1.87}{\volt}, and three conversions in one run repeat to
\SI{2}{LSB}. One LSB is \SI{1.593}{\milli\volt} over the converter's measured
\SI{1.63}{\volt} span, so what it resolves in current is set by \(R_f\) alone:
\SI{88.5}{\nano\ampere} at \(N = 0\), \SI{4.02}{\nano\ampere} at \(N = 63\) and
\SI{0.825}{\nano\ampere} at \(N = 319\). The zero calibration consequently
cannot be taken at the bottom gain code, where the whole
\SI{34.4}{\nano\ampere} offset is sub-LSB: measured there it reads
\SI{-1.5}{LSB}, \SI{-130}{\nano\ampere} of apparent current on a
\SI{-34}{\nano\ampere} signal. It must be taken at a high-gain tap, where one
LSB is smaller than the offset being measured --- which the gain-independent
zero of Section~\ref{sss:pixtop-zero} makes free.

\paragraph{Bit 9 is inverted (firmware-critical).} The raw output bus carries
its most significant bit inverted --- the two's-complement form of the
offset-binary code Section~\ref{ss:adc} describes. Every published converter
number, including the transfer of Table~\ref{tab:adc-transfer}, is the
corrected code \(\mathrm{raw} + 512 - 1024\,b_9\), so firmware must invert bit~9
before using the bus. Read raw, the transfer is a plausible straight line that wraps at the
common-mode code --- a silent 512-code error over half the range.

% The multiplexer fabric that stands between the channel and both the converter
% and the test port. Kept inside this section rather than promoted to a block
% section of its own: it is fabric in the channel's signal path, and its two
% consequences (a sibling-voltage rule and a settling wait) are channel-level
% procedures that the readout paragraph above already depends on.
% Source ~/chips/castalia/simruns/muxchar_20260825/README.md.
% % HAND-WRITTEN fragment -- NOT generated by maestro2tex. Input from
% AnalogChapter.tex immediately after note_pixtop_adc.tex.
%
% PLACEMENT DECISION (2026-08-25): kept as a subsubsection of the assembled
% channel rather than promoted to a block subsection of its own.  The fabric is
% in the channel's signal path (two 16:1 trees, ADC_MUX -> BUF_ADC -> converter
% and ATP_MUX -> BUF_ATP -> test port, 20 select bits across the 1.0/2.5 V
% boundary), its two consequences are channel-level procedures, and the
% converter-readout paragraph above already depends on its numbers.  It carries
% its own row in the coverage table, pointing at this label.
%
% Source: ~/chips/castalia/simruns/muxchar_20260825/README.md, sections 3-8a,
% with results_tables.txt, leak_vs_neighbour.txt and spec_summary.txt.
% Benches: castalia/tb_anatop_amux/AmuxTB + maestro (5 tests amux_leak,
% amux_leakv, amux_ac, amux_tran, amux_mc) and castalia/tb_anatop_lvl/LvlTB +
% maestro (2 tests lvl_dc, lvl_tran).  Both new cells; anatop_pixel,
% anatop_amux_* and anatop_lvl_* were not modified.  Schematic level.
% Corners _default(37 C) / T25 / T40 / T85 / ss_hot / ff_cold.
%
% Numbers used below (all from results_tables.txt unless noted):
%  Leakage, pA, at 25 / 37 / 40 / 85 / ss85 / ff-40 C:
%    leak_sib_vcm  0.062 0.061 0.061 0.048 0.047 0.069
%    leak_sib_0    9.17  20.6  24.7  300   73.7  0.167
%    leak_max      19.2  42.1  50.8  582.6 85.7  1.37
%    (leakv_25, slot at 2.5 V with the tree at 0 V: 113 / 237 / 283 / 2792 /
%     391 / 14.3 pA -- the 0-1 supply-node pair of the slot map)
%  Knee at 40 C (leak_vs_neighbour.txt): |dV| 0-1.00 V -> 0.05-0.07 pA,
%    1.125 V -> 0.37 pA, 1.25 V -> 24.7 pA.  At 85 C: 0.05 / 8.5 / 298 pA.
%    Mechanism GIDL: the off pass device's gate sits at a rail, so |Vdg|
%    reaches ~2.3 V exactly as |dV| passes ~1.1 V.  Junction leakage bounded
%    separately at 0.047 pA (leakv_0), not the limit.
%  Sibling-only proof: victim 1.25 V / sibling 0 V / other fourteen at 0 V =
%    24.71 pA; victim 0 V / sibling 1.25 V / other fourteen at 1.25 V =
%    24.86 pA.  Slot symmetry A<12> vs A<7>: 0.087/24.714/282.316 against
%    0.087/24.714/282.314 pA.
%  Ron, ohm: ron_vcm 1680 1739 1753 1971 2588 1041; ron_max 3572 3682 3709
%    4117 5029 2415, peak at 1.80 V (both pass devices w=1u multi 4, so the
%    PMOS is the weaker half); ron_min (at 0 V) 971 1007 1016 1151 1327 696.
%    Spec ron_max < 6 kohm passes at every corner.
%  Aperture: BUF_ADC input capacitance MEASURED 23.4 fF (decks/exp_cin.py, AC
%    current into InP at the real bias point, flat 1 kHz - 1 MHz).  A 1.25 V
%    slot change settles to 0.5 LSB (797 uV) in 2.36 ns; into the test pad
%    (1.2 kohm + 50 pF) 1.42 us.  Charge injection +57.0 / -7.6 mV on the
%    23.4 fF node, gone in 1.2 ns.  Bandwidth 1.74-3.47 GHz into that load.
%    Channel-level cross-check: rebase_20260824 S2.1's 0.54 us slot-1
%    disturbance is the TIA and buffer recovering, not the mux.
%  Make-before-break (S7), 1 pF probe on slot 0, 0.94 V across the pair, 37 C:
%    6->9 (all four bits) -119.4 fC; 6->7 (Sel<0> only) -118.5 fC;
%    6->14 (Sel<3> only) 0.000 fC -- the control that proves it is per-bit;
%    6->9 with 20 ns select edges -829.6 fC (7x); through the real shifters
%    -11.7 fC.  Corners -91.7 to -137.4 fC (the -17.4 fC at tt/85 C is the
%    10 GOhm probe leaking, not the mux).  ~0.13 fC per mV of pair difference
%    per Sel<0> toggle at 1 ns edges.
%    Skew: 5 ns between Sel<0> and Sel<3:1> swings Y the full 1.25 V to a third
%    slot and stretches settling 2.4 ns -> 6.3 ns; 50 ns of skew -> 51 ns.
%    Level-1 sibling pairs, read from the netlist: (0,1)(2,3)(4,5)(6,7)(8,9)
%    (10,11)(12,13)(14,15).
%  Level shifters (S5): non-inverting, VOH 2.5 V, VOL 3.3 nV - 534 nV,
%    threshold 0.516-0.551 of the 1.0 V rail, t_pLH 1.22-1.99 ns, t_pHL
%    1.96-2.48 ns into 20 fF, static I25 < 2.8 nA per 6 channels.  Crowbar at
%    threshold 15.3 uA per 6 channels = 2.5 uA/channel; 20 select bits = 51 uA
%    on the ANALOG rail.  Power-up with the 1.0 V domain dark: no select line
%    exceeds 283 mV against a ~550 mV pass threshold, parked slot moves
%    <= 517 uV when the 1.0 V rail rises.  _x6 is 3 x the dual cell and _x14 is
%    7 x, so characterising the dual characterises all three.
%  Crosstalk (S6): 114 nV pp from the nearest unselected slot for a 2.5 V step,
%    -174 dB at 10 kHz, off-state capacitance 5.3 fF measured.  MODEL FLOOR --
%    BSIM4 gives a cut-off device no source-drain through-capacitance and the
%    off gate acts as a grounded shield, so the physical number is a layout
%    question and needs an extracted view.  Deliberately not published as a
%    silicon figure.
%  Ron matching MC (S8): 50 mismatch samples, four instances at select codes
%    0/5/10/15, sigma = 0 exactly, slot-to-slot difference 21.6 uOhm.  All 90
%    pass devices per 16:1 are on raw nch_25/pch_25 cards, which carry no
%    per-device mismatch.  Bound: at 1 uA Ron drops 1.7 mV, so at the fA-to-nA
%    currents this fabric carries a 20 % Ron mismatch is < 1 uV.
%  Floating inputs (S8a), from simruns/rebase_20260824/canon/input.scs:
%    ADC_MUX has 11 of 16 inputs floating (A<3>, A<6..15>), ATP_MUX all 16.
%    A floating input settles at a leakage equilibrium (0.9376 V at 37 C
%    through 10 GOhm), and A<3>'s level-1 sibling is A<2> = V_PATTERN, the
%    potential-program DAC output.  Fix is one wire per input, to V_CM and not
%    to local ground (S3's sibling rule).  note_pixtop_open already carries the
%    floating-input finding qualitatively; this is its cost.
%  Netlisting blocker, unrelated: anatop_lvl_d10a25_x6 needs a GUI
%    Check-and-Save (OSSHNL-109) before it can be netlisted; _x14 untested.
%    Not a TRM statement.
% Library cell names are kept out of the rendered prose.

\subsubsection{Input Multiplexer}\label{sss:pixtop-mux}

Both the converter and the test port see the channel through a 16:1
transmission-gate tree, and twenty select bits reach those trees through
\SI{1.0}{\volt}-to-\SI{2.5}{\volt} level shifters. The fabric is in the signal
path, so its leakage, its on-resistance and its switching behaviour are channel
specifications rather than debug conveniences.

\paragraph{Off-state leakage meets its budget conditionally.} A slot's leakage
is set by the voltage across the one off switch it shares with its level-1
sibling, and the dependence has a knee rather than a slope: at \SI{40}{\celsius}
it is \SI{0.05}{\pico\ampere} anywhere within \SI{1.00}{\volt} of the sibling,
\SI{0.37}{\pico\ampere} at \SI{1.125}{\volt} and \SI{24.7}{\pico\ampere} at
\SI{1.25}{\volt} (Table~\ref{tab:pixtop-mux}). At \SI{85}{\celsius} the same
three points are \num{0.05}, \num{8.5} and \SI{298}{\pico\ampere}. The
mechanism is gate-induced drain leakage in the off pass device, whose gate sits
at a supply rail, so the drain--gate field reaches its onset exactly as the
slot-to-sibling difference passes \SI{1.1}{\volt}; junction leakage is bounded
separately at \SI{0.047}{\pico\ampere} and is not the limit. Against a
\SI{20}{\pico\ampere} budget the fabric therefore passes with \(280\times\)
margin while neighbouring slots sit near one another, and misses by
\(1.2\times\) at \SI{40}{\celsius} and \(15\times\) at \SI{85}{\celsius} when a
slot's sibling sits at a rail.

\paragraph{The tree earns its place, measurably.} Only the level-1 sibling
appears in the answer: holding the other fourteen inputs at ground and at
\SI{1.25}{\volt} gives \num{24.71} and \SI{24.86}{\pico\ampere} for the same
slot, and every slot crosses the same four gates, so no slot is privileged. A
flat sixteen-way star has no such structure --- every off switch faces the
shared node, so a sense slot's leakage would be set by whichever slot is
currently selected, a quantity the measurement itself changes. The costs are
four gates in series and the switching behaviour below.

\paragraph{On-resistance peaks high, not mid-rail.} It is \SI{3.68}{\kilo\ohm}
at \SI{1.80}{\volt} and \SI{37}{\celsius}, worst \SI{5.03}{\kilo\ohm} at the
slow hot corner, against a \SI{6}{\kilo\ohm} limit. The hump sits above mid-rail
because the two pass devices are the same width, which makes the PMOS the weaker
half. It does not threaten the sampling aperture by four orders of magnitude:
the multiplexer drives a gate, whose input capacitance measures
\SI{23.4}{\femto\farad}, and a \SI{1.25}{\volt} slot change settles to half a
converter code in \SI{2.36}{\nano\second}. Read at the test pad instead, into
\SI{1.2}{\kilo\ohm} and \SI{50}{\pico\farad}, the same change takes
\SI{1.42}{\micro\second} --- irrelevant for a meter, decisive for a procedure
that assumes a settled node. Coupling from unselected slots simulates at
\SI{114}{\nano\volt} peak-to-peak for a full-scale aggressor step, but that is a
model floor: a cut-off device is given no source-to-drain path, so the physical
number is set by routing and needs an extracted view.

\paragraph{Make-before-break is real.} Each two-way cell drives its pass gates
from a select bit and its locally inverted copy, so during any transition both
conduct. A toggle of the lowest select bit does this in all eight level-1 cells
at once, shorting all eight sibling pairs; the higher bits short the group
outputs. Measured on a parked slot with \SI{0.94}{\volt} across its pair,
\SI{119}{\femto\coulomb} is transferred per toggle at \SI{1}{\nano\second}
select edges, \(\approx \SI{0.13}{\femto\coulomb}\) per millivolt of pair
difference, rising to \SI{830}{\femto\coulomb} at \SI{20}{\nano\second} edges. A
change that toggles only the top select bit moves the same probe by
\SI{0}{\femto\coulomb}, which is what makes the effect per-bit and predictable
rather than diffuse. Select-bit skew is the sharper hazard: \SI{5}{\nano\second}
of skew between select bits swings the output the full \SI{1.25}{\volt} through
a third, unintended slot and stretches settling from \num{2.4} to
\SI{6.3}{\nano\second}; \SI{50}{\nano\second} of skew stretches it to
\SI{51}{\nano\second}. A reading taken too soon after a slot change is not a
noisy version of the right slot --- it is another slot's voltage. Removing this
would need a decoder and a non-overlap generator, which is the device count the
tree exists to avoid, so it is handled by procedure:

\begin{enumerate}[nosep]
  \item \textbf{Keep every unselected slot near the common mode}, and in no case
    more than \SI{1}{\volt} from its level-1 sibling. Tying spares to local
    ground satisfies neither. The level-1 pairs are fixed by the tree ---
    \((0,1)\), \((2,3)\) and so on --- and the channel as drawn violates the
    rule: the test-port tree has all sixteen inputs undriven and the converter
    tree eleven of sixteen, and one of those undriven slots is the level-1
    sibling of the commanded-potential output, into which it dumps its
    accumulated charge on every toggle of the lowest select bit.
  \item \textbf{Wait after any slot change before sampling}, for the settling
    time plus the worst-case skew between the four select bits of that tree, and
    drive those four bits from one shifter instance so the skew is bounded.
\end{enumerate}

\paragraph{Level shifters.} Non-inverting, output high at \SI{2.5}{\volt} and
low below \SI{0.6}{\micro\volt}, threshold \numrange{0.52}{0.55} of the
\SI{1.0}{\volt} rail, propagation \SIrange{1.2}{2.5}{\nano\second}. The figure
to carry into the power budget is the crowbar: a select input parked at its own
threshold --- a floating select bus, a half-driven pin, a slow ramp --- draws
\SI{2.5}{\micro\ampere} per channel, and twenty select bits is
\SI{51}{\micro\ampere}, on the analog rail rather than the digital one. With the
\SI{1.0}{\volt} domain unpowered and the analog domain live, no select line
exceeds \SI{283}{\milli\volt} against a \(\approx\SI{550}{\milli\volt}\) pass
threshold and no switch closes that would not have closed anyway; a parked slot
moves by at most \SI{0.52}{\milli\volt} when the \SI{1.0}{\volt} rail comes up.

\paragraph{Slot-to-slot on-resistance matching is not measurable here.} Every
path crosses the same four gates, and nominal on-resistance is identical across
slots to seven digits; a 50-sample mismatch run returns \(\sigma = 0\) exactly,
because all ninety pass devices per tree sit on model cards that carry no
per-device mismatch. The concern is bounded rather than measured: at
\SI{1}{\micro\ampere} the on-resistance drops \SI{1.7}{\milli\volt}, so at the
femtoampere-to-nanoampere currents this fabric carries a \SI{20}{\percent}
mismatch is under \SI{1}{\micro\volt}. Slot-dependent on-resistance is not a
measurement-error mechanism for this path.

\begin{table}[htbp]
  \centering
  \caption[{Off-state leakage and on-resistance of one 16:1 input multiplexer tree.}]{Off-state leakage and on-resistance of one 16:1 input multiplexer tree, schematic level. Leakage is the current out of one slot with the tree selecting elsewhere, and is quoted against the position of that slot's level-1 sibling; the worst-neighbour row sweeps the sibling over the whole rail. On-resistance is measured with \SI{1}{\micro\ampere} forced out of the tree output, and its peak sits at \SI{1.80}{\volt} rather than mid-rail. The leakage budget is \SI{20}{\pico\ampere} and the on-resistance limit \SI{6}{\kilo\ohm}.}
  \label{tab:pixtop-mux}
  \begin{tabular}{lrrrrrr}
    \toprule
    Parameter & \SI{25}{\celsius} & \SI{37}{\celsius} & \SI{40}{\celsius} & \SI{85}{\celsius} & ss \SI{85}{\celsius} & ff \SI{-40}{\celsius} \\
    \midrule
    Leakage, sibling at $V_{\mathrm{CM}}$ [\si{\pico\ampere}] & 0.062 & 0.061 & 0.061 & 0.048 & 0.047 & 0.069 \\
    Leakage, sibling at ground [\si{\pico\ampere}]            & 9.17  & 20.6  & 24.7  & 300   & 73.7  & 0.167 \\
    Leakage, worst neighbour [\si{\pico\ampere}]              & 19.2  & 42.1  & 50.8  & 582.6 & 85.7  & 1.37 \\
    \midrule
    $R_{\mathrm{on}}$ at $V_{\mathrm{CM}}$ [\si{\ohm}] & 1680 & 1739 & 1753 & 1971 & 2588 & 1041 \\
    $R_{\mathrm{on}}$ peak [\si{\ohm}]                 & 3572 & 3682 & 3709 & 4117 & 5029 & 2415 \\
    \bottomrule
  \end{tabular}
\end{table}

% Small-signal, noise and power: simruns/rebase_20260824/{stb,stb_buf,noise,
% noise_adc,pwr,corner} on the 2026-08-24 OA, with simruns/trmnoise_gain (the
% 16-code noise sweep behind fig_pixtop_noise, whose numbers the re-baseline
% reproduces to four significant figures) and simruns/postfix_decim (the
% transient-noise decimation records).
% % HAND-WRITTEN fragment -- NOT generated by maestro2tex. Input from
% AnalogChapter.tex after note_pixtop_adc.tex.
%
% Provenance of every number below (all measured 2026-08-23 on the post-fix
% deck simruns/postfix_canon/input.scs, netlisted from the OA by
% skill/eis_headless_run.sh; engine skill/pixchar_engine.py --mode smallsig,
% margins from spectre's own stbF/stbC.margin.stb):
%   C_F_TIA = tsmcN65/mimcap_sin lt=wt=52.09u mimflag=3 mf=4 = 22.0052 pF,
%     FB -> V_OUT inside castalia/anatop_tia; area 4 x 52.09^2 = 10 853 um^2.
%     All three mimcap_sin in that cell now carry mismatchflag=1 (owner's edit,
%     2026-08-23) -- this is the second of the two Monte Carlo blockers of
%     sss:pixtop-open, now cleared.
%   24-point grid, PM/GM/peaking/f3dB, pre- and post-fix side by side:
%     simruns/postfix_stb/stb_table.txt (post), simruns/pixchar_stb/stb_table.txt
%     (pre; 6 of 24 points FAIL, worst 35.95 deg / 4.542 dB at Cdl 100n Rct 1M
%     N = 191).
%   5 corners at N = 191 and N = 63: simruns/postfix_corner_stb/corner_stb_table.txt
%     (pre-Cf worst 30.00 deg / 5.957 dB at ss, same cell).
%   Control-loop shift 21 -> 36-41 dB GM and 18 -> 6 MHz crossover: same two
%     tables, PM_ctrl/GM_ctrl/fx_ctrl columns.
%   f3dB pre -> post at Cdl 100 nF, Rct 1 M: 116.6k -> 85.77k (N=0),
%     13.59k -> 10.00k (63), 7356 -> 4642 (191), 5412 -> 3415 (319).
% Library cell names are kept out of the rendered prose (owner request).

\subsubsection{Compensation and Stability}\label{sss:pixtop-stability}

The transimpedance stage carries a \SI{22.0}{\pico\farad} capacitor across
\(R_f\), a four-unit metal--insulator--metal array occupying
\SI{10853}{\micro\meter\squared} per channel. All three capacitors in that
stage now carry mismatch statistics.

It exists because the assembled loop did not close without it. With the
programmable ladder and the real reference in the loop --- neither of which
Section~\ref{sss:pixel-grid} had --- 6 of 24 electrode-by-gain grid points
failed, falling to \SI{35.95}{\degree} of phase margin with \SI{4.54}{\decibel}
of peaking at the nominal corner and to \SI{30.00}{\degree} with
\SI{5.96}{\decibel} at the slow one. With the capacitor all 24 points pass:
the worst cell is \(C_{\mathrm{dl}} = \SI{100}{\nano\farad}\),
\(R_{\mathrm{ct}} = \SI{1}{\mega\ohm}\), \(N = 63\) at \SI{75.71}{\degree} and
\SI{1.068}{\decibel}, and it holds \SI{72.37}{\degree} and
\SI{1.379}{\decibel} at the slow corner. Gain margin is at least
\SI{21.15}{\decibel} everywhere. The grid spans
\(C_{\mathrm{dl}} = \SI{100}{\nano\farad}\,/\,\SI{1}{\micro\farad}\,/\,\SI{10}{\micro\farad}\),
\(R_{\mathrm{ct}} = \SI{10}{\kilo\ohm}\,/\,\SI{1}{\mega\ohm}\) and
\(N = 0\,/\,63\,/\,191\,/\,319\), and the two worst cells were re-run at five
corners.

The capacitor also lands in the control loop, because it ties the
transimpedance output to the working electrode and therefore sits inside A1's
path through the electrode: A1's gain margin rises from \SI{21}{\decibel} to
\SIrange{35.5}{40.7}{\decibel}, its crossover falls from \SI{18}{\mega\hertz}
to \SI{6.0}{\mega\hertz}, and its phase margin from
\SIrange{68.1}{71.4}{\degree} to \SIrange{79.0}{81.8}{\degree}.

The cost is closed-loop bandwidth, which falls by a third to a half at every
code: \SIrange{116.6}{85.8}{\kilo\hertz} at \(N = 0\),
\SIrange{13.6}{10.0}{\kilo\hertz} at \(N = 63\) and
\SIrange{5.4}{3.4}{\kilo\hertz} at \(N = 319\). Against the
\SI{2}{\milli\second} staircase tread of Section~\ref{sss:pixtop-cv} even the
slowest of those settles two decades early, so no measurement in this section
is bandwidth-limited.

The five-corner rows keep the resistor and capacitor models typical, so neither
\(R_f\) nor the compensation capacitor corners; what they exercise is the
transistors alone.

% % HAND-WRITTEN fragment -- NOT generated by maestro2tex. Input from
% AnalogChapter.tex after note_pixtop_comp.tex; fig_pixtop_noise.tex follows it.
%
% Provenance (all 2026-08-23, post-C_F deck simruns/postfix_canon/input.scs,
% nominal section castalia_typ_25 at 37 C, SAR frozen, Rct = 1 M):
%   ac noise, 4 gain codes x 3 Cdl, 0.1 Hz - 10 MHz:
%     simruns/postfix_noise/{noise_table.txt,boxcar_table.txt,antialias_table.txt}
%   ac noise, 16 gain codes at Cdl = 100 nF (the figure):
%     simruns/trmnoise_gain/noise_vs_gain_table.txt, driver skill/trmnoise_run.py,
%     reduction + .dat by skill/pixtop_noise_fig.py
%   transient noise, 3 codes x 6 seeds x 20 ms, strobed at 1.150748 us:
%     simruns/postfix_decim/{decim_table.txt,correlation_table.txt}
%   the spec rewrite itself, and the 1 uF / 10 uF electrode rows:
%     simruns/postfix_specs.txt section 2
% Key numbers: retired px_noise.in_rma 4.588/4.378/4.377/4.377 nA at
%   N = 0/63/191/319; sigma_tread 2.889/1.707/1.559/1.475 nA = 0.036/0.423/1.134
%   /1.781 LSB; I_LSB 81.17/4.036/1.374/0.8282 nA on the MEASURED Zt(dc)
%   (19.63k/394.8k/1.159M/1.924M), which is 8.4 % above the nominal ladder at
%   the bottom code -- hence 81.2 nA here against the 88.5 nA of sss:pixtop-adc,
%   which is quoted on the nominal 18 kOhm.
%   sqrt(N) error 6.674x at N = 100, code 319 (optimistic) and 0.350x at
%   N = 1738, code 0 (pessimistic); rho(1) = 0.988, rho(10) = 0.830 at code 63+,
%   rho(10) = -0.077 at code 0.  Dirichlet prediction vs measurement 0.4-1.8 %
%   except the two largest N at code 0, where a 20 ms record realises nothing
%   below 50 Hz and the prediction is the upper bound.
%   Anti-alias: brickwall at the 434.5 kHz Nyquist multiplies sigma_1 by 1.054
%   (N=0), 1.003 (63), 1.001 (191, 319); alias share of variance 10.0 % at N = 0.
% NOTE the `in 0.1-10M` column of noise_table.txt is spectre's in(f) integral and
% is NOT a per-sample error -- it divides by a rolled-off transimpedance and comes
% out gain-independent at ~1170 nA.  Do not publish it.
% Library cell names are kept out of the rendered prose.

\subsubsection{Noise and the Sampled Readout}\label{sss:pixtop-noise}

A flat current-noise budget cannot be written for this channel. One converter
LSB is \SI{81.2}{\nano\ampere} at \(N = 0\) and \SI{0.828}{\nano\ampere} at
\(N = 319\), so the \SI{2}{\nano\ampere} rms budget of
Sections~\ref{ss:ampab} and~\ref{ss:pixel} is \SI{0.025}{LSB} at one end of the
ladder and \SI{2.4}{LSB} at the other. The budget is therefore stated against the code in use: the rms of the
readout current, after averaging every converter sample taken inside one
staircase tread, must stay below one LSB at that code. A tread is
\SI{2.00}{\milli\second} and holds 1738 samples at the converter's
\SI{869}{\kilo\hertz} sample rate.

Measured that way, the channel passes at \(N = 0\) (\SI{0.036}{LSB}) and
\(N = 63\) (\SI{0.423}{LSB}) and misses at \(N = 191\) (\SI{1.13}{LSB}) and
\(N = 319\) (\SI{1.78}{LSB}). Figure~\ref{fig:pixtop-noise} shows why: the
averaged floor is \SIrange{1.5}{1.9}{\nano\ampere} at every code above the
bottom one and \SI{2.9}{\nano\ampere} at \(N = 0\), while one LSB falls
\(98\times\) across the ladder. What survives a \SI{2}{\milli\second} average
is the low-frequency input-referred noise of A2, which is a property of the
amplifier and not of \(R_f\), so the top gain codes are noise-limited and not
quantisation-limited --- which is what the ladder was sized for. The
compensation capacitor of Section~\ref{sss:pixtop-stability} does not change
any of this: it divides output noise and transimpedance by the same factor, and
the input-referred figures are unchanged by it to four digits.

The electrode sets the floor at larger area. At
\(C_{\mathrm{dl}} = \SI{1}{\micro\farad}\) the tread-averaged noise is
\SIrange{11.8}{13.7}{\nano\ampere} and at \SI{10}{\micro\farad}
\SIrange{48}{76}{\nano\ampere}; a \SI{10}{\micro\farad} electrode cannot be
read to one LSB in one tread at any gain code, and the tread must lengthen or
the electrode shrink.

\paragraph{Averaging is not \(\sqrt{N}\).} Eighteen transient-noise records ---
three gain codes, six seeds, \SI{20}{\milli\second} each, strobed at the
converter's sample period --- put \(\sqrt{N}\) wrong by up to \(6.7\times\) and
in both directions: at high gain the noise sits three decades below Nyquist and
consecutive samples correlate at \(\rho(1) = 0.988\), so averaging ten of them
buys \SI{2.5}{\percent} rather than \(3.16\times\), while at \(N = 0\)
adjacent samples anticorrelate and the average beats \(\sqrt{N}\). The exact
Dirichlet-kernel integral of the measured spectrum reproduces the measurement
to \SIrange{0.4}{1.8}{\percent} and is the predictor used throughout this
section.

\paragraph{No anti-alias filter is required.} Replacing the sampler with an
ideal brickwall at the \SI{434.5}{\kilo\hertz} Nyquist frequency multiplies the
per-sample rms by \(1.054\) at \(N = 0\) and by \(1.001\)--\(1.003\) at every
other code. At most \SI{10.5}{\percent} of the per-sample variance folds in
from above Nyquist, and only at minimum gain, where it moves the tread-averaged
figure from \SI{0.036}{LSB} to \SI{0.027}{LSB}. An analog anti-alias filter
would buy nothing measurable; a post-transimpedance pole placed in the signal
band is a different question, and the tread average already answers it.

% %% HAND-WRITTEN fragment -- NOT generated by maestro2tex (the engine has no
%% quantity-vs-code figure kind and these are raw Spectre runs, not a Maestro
%% results database). Precedent: fig_electrode_validation.tex (two panels as two
%% tikzpictures sharing a pgfplotsset style), fig_pixtop_cv.tex (colours,
%% grid/legend styling, \MaestroRoot).
%%
%% Data: data/pixtop_noise_{sig1,tread,lsb,sig1_lsb,tread_lsb}.dat, written by
%% ~/chips/castalia/skill/pixtop_noise_fig.py from the 16 noise runs of
%% ~/chips/castalia/simruns/trmnoise_gain/ (driver skill/trmnoise_run.py).
%% Each run: pixchar_engine --mode smallsig --inj --freeze-sar --no-stb
%% --with-noise, `noise ( DUT.V_OUT DUT.V_CM ) iprobe=I_NIN` 0.1 Hz - 10 MHz at
%% 15 pts/dec, --rct 1M --param cdl_we=100n, nominal section castalia_typ_25 at
%% 37 C, on the post-C_F deck simruns/postfix_canon/input.scs. 16/16 completed,
%% 0 spectre errors, ~14 s each.
%% Codes 0/8/16/32/63/64/96/127/128/160/191/192/255/256/288/319 -- the four
%% ladder seams (63/64, 127/128, 191/192, 255/256) are straddled deliberately.
%%
%% Columns:
%%   sigma_1     = total OUTPUT noise rms / Zt(dc), an ideal instantaneous
%%                 sampler with no anti-alias filter. NOT spectre's
%%                 input-referred in(f) integral to 10 MHz, which divides by a
%%                 rolled-off transimpedance and is gain-independent (~1170 nA)
%%                 -- see simruns/postfix_noise/noise_table.txt footnote.
%%   sigma_tread = the same PSD through the exact Dirichlet-kernel integral of
%%                 an N = 1738 boxcar at Ts = 1/869 kHz (skill/postfix_boxcar.py).
%%                 Validated against 18 spectre transient-noise records to
%%                 0.4-1.8 % (simruns/postfix_decim/decim_table.txt); sqrt(N) is
%%                 wrong there by up to 6.7x, in both directions.
%%   I_LSB       = 1.5934 mV / Zt(dc). 1.5934 mV is the converter's measured LSB
%%                 over its 1.63 V span (sss:pixtop-adc). Zt(dc) is the MEASURED
%%                 in-loop transimpedance, 8.4 % above the nominal 18k+6k*N at
%%                 the bottom code because of the closed-switch stack, so
%%                 I_LSB(N=0) is 81.2 nA and not the nominal-ladder 88.5 nA.
%% Cross-check: at N = 0/63/191/319 these reproduce
%% simruns/postfix_noise/boxcar_table.txt to every printed digit.
%% Two panels, not a twin axis: the normalisation divides by I_LSB, which itself
%% spans 98x over the codes, so (b) is not a rescaling of (a).
\providecommand{\MaestroRoot}{include/analog/}
\providecolor{mtxA}{HTML}{2A78D6}
\providecolor{mtxB}{HTML}{EB6834}
\providecolor{mtxC}{HTML}{1BAF7A}
\begin{figure}[htbp]
  \centering
  {\pgfplotsset{compat=1.18}%
  \pgfplotsset{
    nzax/.style={
      width=0.42\linewidth, height=5.6cm,
      grid=both,
      major grid style={line width=0.2pt, draw=black!14},
      minor grid style={line width=0.2pt, draw=black!7},
      minor tick num=1,
      axis line style={line width=0.4pt, draw=black!45},
      tick style={line width=0.4pt, draw=black!45},
      tick align=outside,
      label style={font=\small}, tick label style={font=\small},
      legend style={font=\footnotesize, draw=none, fill=white,
                    fill opacity=0.85, text opacity=1, row sep=1pt},
      every axis plot/.append style={line join=round, line width=0.9pt,
                                     mark size=1.5pt},
      xlabel={Gain code \(N\)},
      xmin=-12, xmax=331, xtick={0,64,128,192,256,320},
    }}%
  \begin{tikzpicture}
    \begin{axis}[nzax, ymode=log,
      ylabel={Noise current~[\si{\nano\ampere} rms]},
      ymin=0.5, ymax=600,
      legend style={at={(0.97,0.97)}, anchor=north east},
    ]
      \addplot[mtxA, mark=*] table {\MaestroRoot data/pixtop_noise_sig1.dat};
      \addlegendentry{\(\sigma_1\), per sample}
      \addplot[mtxC, mark=square*, dashed]
        table {\MaestroRoot data/pixtop_noise_lsb.dat};
      \addlegendentry{one converter LSB}
      \addplot[mtxB, mark=triangle*]
        table {\MaestroRoot data/pixtop_noise_tread.dat};
      \addlegendentry{\(\sigma_{\mathrm{tread}}\), \SI{2}{\milli\second} average}
    \end{axis}
  \end{tikzpicture}\hfill
  \begin{tikzpicture}
    \begin{axis}[nzax, ymode=log,
      ylabel={Noise~[LSB at that code]},
      ymin=0.02, ymax=60,
      legend style={at={(0.97,0.03)}, anchor=south east},
    ]
      \addplot[mtxA, mark=*] table {\MaestroRoot data/pixtop_noise_sig1_lsb.dat};
      \addlegendentry{\(\sigma_1\), per sample}
      \addplot[mtxB, mark=triangle*]
        table {\MaestroRoot data/pixtop_noise_tread_lsb.dat};
      \addlegendentry{\(\sigma_{\mathrm{tread}}\), \SI{2}{\milli\second} average}
      \addplot[black!55, line width=0.7pt, densely dotted, no marks,
        domain=-12:331, samples=2] {1};
      \addlegendentry{one LSB}
    \end{axis}
  \end{tikzpicture}}
  \caption[{Input-referred noise against gain code, in current and in
  converter LSBs.}]{Input-referred current noise of the assembled channel
  against gain code, at the nominal electrode
  (\(C_{\mathrm{dl}} = \SI{100}{\nano\farad}\), \(R_{\mathrm{ct}} =
  \SI{1}{\mega\ohm}\)), nominal process point, \SI{37}{\celsius}, integrated
  from \SI{0.1}{\hertz} to \SI{10}{\mega\hertz} and referred by the measured
  in-loop transimpedance. \(\sigma_1\) is one converter sample; \(\sigma_{\mathrm{tread}}\)
  is the average of the 1738 samples taken inside one \SI{2}{\milli\second}
  staircase tread, by the exact Dirichlet integral of the same spectrum ---
  \emph{not} \(\sigma_1/\sqrt{N}\), which the transient-noise measurement of
  Section~\ref{sss:pixtop-noise} finds wrong by up to \(6.7\times\). Left, in
  current: the averaged floor is \SIrange{1.5}{2.9}{\nano\ampere} and barely
  moves with gain, while one LSB falls \(98\times\), so the two cross at
  \(N \approx 165\) and every code above it is noise- rather than
  quantisation-limited. Right, the same referred to the LSB at each code; a
  single sample is never better than \SI{3.4}{LSB}, which is why the readout
  average is part of the specification and not an option. The feedback
  capacitor does not appear here: it divides output noise and transimpedance by
  the same factor and leaves the input-referred quantity unchanged to four
  digits.}
  \label{fig:pixtop-noise}
\end{figure}

% % HAND-WRITTEN fragment -- NOT generated by maestro2tex. Input from
% AnalogChapter.tex after fig_pixtop_noise.tex, before note_pixtop_open.tex.
%
% Provenance: simruns/rebase_20260824/pwr/power_table.txt and
% corner/corner_table_raw.txt (2026-08-24 OA, deck rebase_20260824/canon/input.scs;
% dcOp, SAR frozen, pwr=subckts, tt @37 C).  OLD = simruns/postfix_pwr, the
% pre-edit cell.
%   tt, extracted converter: 528.7005 uA on 2.5 V + 1.0391 uA on 1.0 V
%     = 1.3228 mW (OLD 656.5548 + 1.0391 = 1.6424 mW).  Schematic converter:
%     528.7669 + 365.2395 uA, whose 1.0 V column is the empty-CDAC crowbar
%     artefact -- NOT a power number, do not publish it.
%   corner 2.5 V rail, extracted: 913.5534 (ff) 579.2546 (sf) 528.7005 (tt)
%     483.5714 (fs) 431.9966 (ss) uA.  OLD 1075.1642 ... 517.1019.
%   per block, extracted, pwr=subckts: TIA 243.7677; CONTROLAMP 243.7675;
%     BUF_ADC 243.7761; BUF_ATP 243.3469; BIASGEN 59.2659; DAC_VP 44.4113 +
%     DAC_VCM 44.4113; ADC 12.3496; the two muxes 0.0007.  Rollup 1135.0970 uW
%     against 1322.7904 delivered by the rails -- the gap is the rollup not
%     closing across the DUT's inherited vdd!/vss! ports, so the rail figure is
%     the authoritative one.
%   TRAP: pwr=subckts vs pwr=total moved the SCHEMATIC-converter 2.5 V rail by
%     5.7 %; with the extracted converter they agree to every digit.  All
%     figures above are the extracted-converter ones.
%   ATTRIBUTION, closes to 0.05 uW: -806.758 uW (DAC_VCM_PSU + DAC_VP_PSU,
%     403.3791 each, DELETED from the cell) +243.776 (BUF_ADC) +243.347
%     (BUF_ATP) = -319.635 uW against a measured 1642.426 -> 1322.790 uW.
%     The regulators were dead load (old F-4: instantiated, drove nothing --
%     DAC_VCM and DAC_VP draw 44.4113 uW each before AND after, to every printed
%     digit), so their 806.8 uW was pure waste.
%   corners, 20 dc-reachable specs: corner/corner_table_raw.txt, 20/20 pass at
%     tt/ff/fs/sf/ss as before.  The only entry that moved is v_adcin_p:
%     1.24850 (tt) against 1.24912 before, margin to the 1.245 V limit 3.50 mV
%     where it was 4.12 mV -- exactly BUF_ADC's own -619.6 uV offset.  Those
%     corner rows keep tt_res and tt_mim.
% 2026-08-24, OWNER INSTRUCTION "do not talk about pixel level yield yet": the
%   closing sentence quoted an uncalibrated per-channel yield; replaced with the
%   distributional statement.  yield_table.txt in the archive still holds it.
% Library cell names are kept out of the rendered prose.

\subsubsection{Power and Corners}\label{sss:pixtop-power}

One channel draws \SI{1.323}{\milli\watt} at the nominal process point,
\SI{528.70}{\micro\ampere} from the \SI{2.5}{\volt} analog rail and
\SI{1.04}{\micro\ampere} from the \SI{1.0}{\volt} converter core, against a
\SI{1}{\milli\watt} target. Over the four process corners the analog rail spans
\SIrange{432}{914}{\micro\ampere}, so the fast corner misses the
\SI{400}{\micro\ampere} rail budget by \(2.3\times\).

Four amplifiers carry it: A1, A2 and the two multiplexer buffers, at
\SI{243.8}{\micro\watt} each and \SI{86}{\percent} of the
\SI{1.135}{\milli\watt} that rolls up per block. The rest is the local bias
generator at \SI{59.3}{\micro\watt}, the two ladders at
\SI{88.8}{\micro\watt} and the converter at \SI{12.3}{\micro\watt}. The rollup
does not close against the \SI{1.323}{\milli\watt} the rails deliver, because it
does not follow the inherited supply ports; the rail figure is the authoritative
one.

The channel is \SI{319.6}{\micro\watt} cheaper than the revision before it, and
the two effects cancel almost exactly: deleting the two reference-DAC supply
regulators returned \SI{806.8}{\micro\watt} that drove nothing, and the two
buffers cost \SI{487.1}{\micro\watt}. The regulators were dead load --- both
ladders draw \SI{44.4113}{\micro\watt} before and after, to every printed digit
--- so the disturbance fix of Section~\ref{sss:pixtop-adc} was bought out of
waste rather than out of budget.

All twenty dc-reachable specs --- reference tracking, the potential range, the
common mode, the zero-current offsets and the gain-code sweep --- pass at all
five corners. Only one moved: the converter-input voltage now sits
\SI{620}{\micro\volt} lower at every corner, the input buffer's own offset,
leaving \SI{3.50}{\milli\volt} of margin to its window where there was
\SI{4.12}{\milli\volt}. That is weaker evidence than it looks: the corner rows
keep the resistor and capacitor models typical, so for a block whose accuracy is
set by matching they exercise almost nothing. Monte Carlo is the real test, and
Section~\ref{sss:pixtop-mc} spreads several of these same quantities by more
than their budget.

% Monte Carlo on the assembled cell, 2026-08-24:
% ~/chips/castalia/simruns/rebase_20260824/{mc200.json,mc_sigma_table.txt} (200
% samples, process + mismatch, section castalia_pm_25, nominal corner, 37 C,
% 1264 random variables) and mc100_bufcm (the buffer's CMRR under mismatch).
% Figure data reduced as ~/chips/castalia/skill/pixmc_trm_fig.py panel_a does.
% % HAND-WRITTEN fragment -- NOT generated by maestro2tex. Input from
% AnalogChapter.tex after note_pixtop_power.tex, before note_pixtop_open.tex.
%
% Provenance: ~/chips/castalia/simruns/rebase_20260824/ (2026-08-24), the
% re-baseline of the edited cell.  200 samples, process and mismatch, model
% section castalia_pm_25, nominal corner, 37 C, savefamilyplots=yes,
% save=selected, subcktprobelvl=0, 25 chunks; 1264 random variables (was 1070 --
% +32 from the four _mac devices x 4 instantiations of anatop_tia_amp_ab, minus
% the deleted regulators' 194 draws), 200/200 converged, no 1.11e36 sentinel.
% Every deck is a textual splice into rebase_20260824/canon/input.scs, netlisted
% from the setup's own Maestro state opened ?mode "r".
%   distributions       mc_sigma_table.txt, mc200.json
%   gain-independence   tia_zero_n0/n63/n319, mean and sigma agree to 5 figures
%                       (-1.00394 / -1.00408 / -1.00788 mV, sigma 5.48789 /
%                       5.48791 / 5.48820 mV); worst per-sample n0-vs-n319
%                       difference 5.40 uV, mean 3.94 uV
%   tracking vs code    mc200.json vec_pat/vec_err, via the same reduction
%                       skill/pixmc_trm_fig.py panel_a performs
%   buffer offset       BUF_ADC's own, mean -0.605 mV sigma 5.325 mV,
%                       r = 0.055 against A2's -- hence v_adcin's 8.040 mV,
%                       close to sqrt(5.616^2 + 5.325^2) = 7.74 mV
%   CMRR residual       mc100_bufcm/buf_cmrr_table.txt (100 samples)
% CAUTION recorded in the archive: this run and pixmc_20260823 are UNPAIRED --
%   deleting the regulator subckt changes the random-variable count and so the
%   whole draw sequence.  A few-per-cent sigma difference on a quantity the edit
%   cannot physically reach (v_cm_dc +12 %, dac_step +8 %) is sampling noise; the
%   95 % interval on sigma at n = 200 is about +-10 %.
%   simruns/macswap_20260824/mc_sigma_table.txt is the matched pair on this same
%   netlist and is the authority on what each change did.
% NOT RE-MEASURED on this OA, and flagged as such in the prose below:
%   the process-only / mismatch-only split and the six single-structure
%   ablations (pixmc_20260823/attribution_table.txt, n = 200 and n = 100), and
%   the separate 200-sample stability run at Cdl = 100 nF, Rct = 1 M, N = 63.
%   The nominal grid moved +0.36 deg / -0.5 dB there, so those distributions are
%   the previous revision's; the mechanism (A2 sets the current axis, A1 the
%   potential axis) is structural and the re-baseline contradicts none of it.
%   Process share of sigma on tia_offset was 0.039263/5.7541 = 0.682 %, i.e.
%   0.0047 % of the variance; sqrt(sigma_proc^2 + sigma_mis^2) reproduced
%   sigma_all to 0.1-0.4 % on every quantity.
% Calibrated worst tracking error, recomputed from mc200.json with
%   pixmc_yield.calibrate()'s model and seed 7 (0.65 mV rms readout floor):
%   0.885 mV over the 0.4-2.1 V window, 5.781 mV over the full sweep.  The same
%   computation reproduces the previously published 0.873 / 5.841 mV exactly on
%   pixmc_20260823, so the method is verified, not re-derived.
%
% 2026-08-24, OWNER INSTRUCTION "do not talk about pixel level yield yet":
%   every per-channel and four-channel yield figure REMOVED -- from this
%   fragment, from tab_pixtop_mc.tex (the Yield column), from
%   tab_pixtop_summary.tex (the yield row), from note_pixtop_power.tex and
%   note_pixtop_open.tex -- and the right-hand yield-vs-window panel deleted
%   from fig_pixtop_mc.tex along with data/pixmc_yield_{track,chan,4ch}.dat.
%   Reason: the figures are upper bounds (four amplifier devices carried no
%   mismatch), the calibration they assume is not executable until the
%   common-mode tap is buffered, and the cell has been edited since the run --
%   the whole characterisation is being re-measured.  yield_table.txt in the
%   archive still holds them; nothing here scores samples against a spec.
%   The subsection was retitled "Monte Carlo" (was "Monte Carlo and Yield").
%   NOTE for whoever restores them: skill/pixmc_trm_fig.py's panel_b() and its
%   three pixmc_yield_*.dat writes are now the only producers of the deleted
%   data files; the script still runs, it just writes files nothing inputs.
% Library cell names are kept out of the rendered prose (owner request).

\subsubsection{Monte Carlo}\label{sss:pixtop-mc}

200 samples of process and mismatch together, at the nominal process point and
\SI{37}{\celsius}, on the tapeout cell itself: the statistical model section
resolves \num{1264} random variables and all 200 samples converge.
Table~\ref{tab:pixtop-mc} is the result. The corner runs of
Section~\ref{sss:pixtop-power} cannot show it: they hold the resistor and
capacitor models typical, and every quantity here is set by matching.

Mismatch is the whole of it. A process-only run leaves \SI{0.68}{\percent} of
\(\sigma\) on the zero-current output --- \SI{0.005}{\percent} of the variance
--- and under \SI{1.5}{\percent} on every other quantity, and the two runs add
in quadrature to the combined \(\sigma\) within \SI{0.4}{\percent}. Six further
mismatch-only runs, each suppressing one matched structure, locate it: the
transimpedance amplifier A2 alone sets the current axis --- suppressing it
collapses \(\sigma\) of the zero-current output to
\SI{0.034}{\milli\volt} --- and the control amplifier A1 alone sets the
potential axis, \SI{0.042}{\milli\volt}. The two R-2R ladders move only the
common mode and the DAC gain; the 320-code feedback ladder and the local bias
generator move no dc quantity by more than the sampling error of a 100-sample
run. The feedback ladder cannot enter: the readout is
\(V_{\mathrm{CM}} - I R_f\), so \(R_f\) is a gain a constant divides out while
the two amplifiers are a raw additive input offset. That offset is what the flat
bench of Section~\ref{ss:pixel} measured mismatch-only at
\SI{5.72}{\milli\volt}; the assembled channel with process \emph{and} mismatch
reproduces it at \SIrange{5.29}{5.49}{\milli\volt}, so assembly and process add
nothing. The die-common term being that small, the four channels on a die are
very nearly independent. The attribution and ablation runs behind this
paragraph were taken on the previous revision of the cell; the re-baseline moves
both amplifier-set spreads by under \SI{5}{\percent} and contradicts none of it.

The converter input is the one node with a wider spread than the transimpedance
output: \(\sigma = \SI{8.040}{\milli\volt}\) against \SI{5.488}{\milli\volt},
because the input buffer contributes its own offset,
\(\sigma = \SI{5.325}{\milli\volt}\), uncorrelated with A2's at \(r = 0.055\).
It does not reach a calibrated reading. Both the signal and the zero are read
through that one buffer, so its offset subtracts out and only its variation with
input level survives: \SI{0.62}{\milli\volt} mean and \SI{2.07}{\milli\volt}
worst of 100 samples over the full readout swing, \SI{0.031}{\milli\volt} and
\SI{0.11}{\milli\volt} over \(\pm\SI{150}{\milli\volt}\)
(Section~\ref{sss:pixtop-adc}).

The zero-current output is gain-independent \emph{per sample}. Its \(\sigma\)
agrees to five significant figures at \(N = 0\), 63 and 319 and its mean to
\SI{4}{\micro\volt}, and within one sample the three taps agree to at worst
\SI{5.4}{\micro\volt}. One stored constant
therefore calibrates all 320 gain codes on a real channel, not only at the
typical point where Section~\ref{sss:pixtop-zero} measured it --- the premise
the whole current-axis calibration rests on.

\paragraph{Calibration.} Every spread above is a per-channel additive constant,
so the two constants of Section~\ref{ss:calibration} remove it, each measured
through the channel's own converter and so carrying that path's
\SI{0.65}{\milli\volt} rms floor. The reference-electrode floor of
Table~\ref{tab:pixtop-mc} is that same offset seen from the other side: the
electrode cannot be driven below ground, so the floor is A1's offset rectified
at zero commanded potential.

\paragraph{The tracking specification is the working window.} The
\SI{2}{\milli\volt} reference-tracking limit is specified over the
\SIrange{0.4}{2.1}{\volt} commanded range, not over the full
\SIrange{0}{2.42}{\volt} code sweep. Inside that window the tracking error is a
pure constant per channel: after the mid-code value is removed the residual is
\SI{0.36}{\milli\volt} at \SI{0.63}{\volt} and falls further toward mid-code
(Figure~\ref{fig:pixtop-mc}). Outside it A1's input pair is running out of
common-mode range, the error stops being a constant, and the same residual is
\SI{3.4}{\milli\volt} at \SI{0.31}{\volt} --- which no stored constant reaches.
On the same samples and the same calibration the worst tracking error averages
\SI{0.885}{\milli\volt} over the window and \SI{5.78}{\milli\volt} over the
full sweep. The difference is a choice of range, not a circuit difference, and
the electrochemistry does not use the outermost \SI{0.4}{\volt} at the bottom or
\SI{0.32}{\volt} at the top.

Stability is the one axis that needs nothing. A separate 200-sample run at the
worst electrode-by-gain point of Section~\ref{sss:pixtop-stability} gives
\SI{74.11}{\degree} of transimpedance phase margin,
\(\sigma = \SI{3.07}{\degree}\) and a worst sample of \SI{64.35}{\degree} ---
\num{9.5} standard deviations above the \SI{45}{\degree} limit. Gain margins are
\(21.57 \pm \SI{0.88}{\decibel}\) and \(38.17 \pm \SI{0.64}{\decibel}\). The
compensation capacitor is not a marginal fix that mismatch undoes. That run is
the previous revision's and has not been repeated; at the same grid point the
nominal phase margin did not move and the nominal gain margin fell
\SI{0.5}{\decibel}.

% %% HAND-WRITTEN table fragment (maestro2tex house style; nothing here is
%% generated). Every value transcribed from
%% ~/chips/castalia/simruns/pixmc_20260823/ -- 200 samples, process and
%% mismatch, model section castalia_pm_25, nominal corner, 37 C, 1070 random
%% variables, 200/200 converged, no 1.11e36 sentinel:
%%   rows 1-10  mc_distributions.txt + yield_table.txt (UNCALIBRATED block)
%%   rows 11-12 stb_table.txt (a separate 200-sample run at the worst point of
%%              the postfix_stb grid: Cdl = 100 nF, Rct = 1 M, N = 63)
%% Means and sigmas re-derived from samples/main200.json and cross-checked
%% against mc_distributions.txt; dac_step 610.2484 uV (the archive README's
%% headline table rounds it to 610.22, the extracted data does not) and
%% v_cm_dc 1.249697 V (README quotes 1.24966).
%% Gain margins, not tabulated: 21.57 +- 0.88 dB (transimpedance loop, min
%% 19.37) and 38.17 +- 0.64 dB (control loop, min 36.07), both 100 %.
%% The yield column is the UNCALIBRATED fraction; the calibrated scenarios are
%% in the prose of sss:pixtop-mc and in yield_table.txt.
%% Library cell names are kept out of the rendered prose.
\begin{table}[!htb]
  \centering
  \caption[{Monte Carlo distributions on the assembled channel.}]{Monte Carlo
  distributions on the assembled channel: 200 samples, process and mismatch
  together, nominal process point, \SI{37}{\celsius}, \SI{2.5}{\volt} supply,
  \num{1070} random variables, all 200 converged. Specifications and yields are
  those the channel's own bench carries, scored \emph{uncalibrated}; the
  calibrated figures are in the text. The two tracking rows are the same
  measurement over two ranges --- the \SIrange{0.4}{2.1}{\volt} row is the
  specification (Section~\ref{sss:pixtop-potential}), the full-sweep row is
  reported because the bench writes it that way. Phase margins come from a
  separate 200-sample run at the worst electrode-by-gain point of
  Section~\ref{sss:pixtop-stability}. \textbf{Every yield here is an upper
  bound}: four devices in each amplifier instantiate raw model cards instead of
  the mismatch-enabled wrappers and therefore contributed no variance
  (Section~\ref{sss:pixtop-open}).}
  \label{tab:pixtop-mc}
  {\footnotesize
  \setlength{\tabcolsep}{5pt}
  \begin{tabular}{@{}%
    >{\raggedright\arraybackslash}p{0.335\textwidth}%
    >{\raggedleft\arraybackslash}p{0.135\textwidth}%
    >{\raggedleft\arraybackslash}p{0.135\textwidth}%
    >{\raggedleft\arraybackslash}p{0.155\textwidth}%
    >{\raggedleft\arraybackslash}p{0.105\textwidth}@{}}
    \toprule
    Quantity & Mean & \(\sigma\) & Specification & Yield \\
    \midrule
    Zero-current output, \(v_{\mathrm{out}} - V_{\mathrm{CM}}\) &
      \SI{-0.425}{\milli\volt} & \SI{5.754}{\milli\volt} &
      \(\pm\SI{2}{\milli\volt}\) & \SI{22.0}{\percent} \\
    Reference tracking error, mid-code &
      \SI{-0.279}{\milli\volt} & \SI{5.552}{\milli\volt} &
      \(\pm\SI{2}{\milli\volt}\) & \SI{31.5}{\percent} \\
    Reference tracking error, worst over \SIrange{0.4}{2.1}{\volt} &
      \SI{4.850}{\milli\volt} & \SI{3.800}{\milli\volt} &
      \(<\SI{2}{\milli\volt}\) & \SI{25.5}{\percent} \\
    Reference tracking error, worst over \SIrange{0}{2.42}{\volt} &
      \SI{9.751}{\milli\volt} & \SI{6.303}{\milli\volt} &
      \(<\SI{2}{\milli\volt}\) & \SI{5.0}{\percent} \\
    \addlinespace
    Zero-current offset referred to current, \(N = 0\) &
      \SI{-23.6}{\nano\ampere} & \SI{319.7}{\nano\ampere} &
      \(\pm\SI{100}{\nano\ampere}\) & \SI{20.5}{\percent} \\
    Same at maximum gain, \(N = 319\) &
      \SI{2.470}{\nano\ampere} & \SI{1.671}{\nano\ampere} &
      \(<\SI{2}{\nano\ampere}\) & \SI{46.5}{\percent} \\
    \addlinespace
    Reference-electrode floor &
      \SI{4.428}{\milli\volt} & \SI{5.913}{\milli\volt} &
      \(<\SI{5}{\milli\volt}\) & \SI{65.5}{\percent} \\
    Reference-electrode ceiling &
      \SI{2.4211}{\volt} & \SI{10.7}{\milli\volt} &
      \(>\SI{2.4}{\volt}\) & \SI{97.0}{\percent} \\
    Common-mode output &
      \SI{1.24970}{\volt} & \SI{0.684}{\milli\volt} &
      \SIrange{1.245}{1.255}{\volt} & \SI{100}{\percent} \\
    Reference DAC step &
      \SI{610.25}{\micro\volt} & \SI{0.171}{\micro\volt} &
      --- & --- \\
    \addlinespace
    Transimpedance phase margin &
      \SI{74.11}{\degree} & \SI{3.07}{\degree} &
      \(>\SI{45}{\degree}\) & \SI{100}{\percent} \\
    Control-loop phase margin &
      \SI{79.94}{\degree} & \SI{1.06}{\degree} &
      \(>\SI{45}{\degree}\) & \SI{100}{\percent} \\
    \bottomrule
  \end{tabular}}
\end{table}

% %% HAND-WRITTEN fragment -- NOT generated by maestro2tex (the engine has no
%% yield-vs-window figure kind, and this is a raw Spectre Monte Carlo spliced
%% into the setup's own netlist, not a Maestro results database). Precedent:
%% fig_pixtop_noise.tex -- two tikzpictures sharing one pgfplotsset style,
%% \MaestroRoot, the mtxA/mtxB/mtxC colours.
%%
%% Data: data/pixmc_track_{sigma,res}.dat and data/pixmc_yield_{track,chan,4ch}.dat,
%% written by ~/chips/castalia/skill/pixmc_trm_fig.py from
%% ~/chips/castalia/simruns/pixmc_20260823/samples/main200.json -- 200 samples,
%% process and mismatch, model section castalia_pm_25, nominal corner, 37 C,
%% 1070 random variables, 200/200 converged.
%%
%% Panel (a) reproduces that run's tracking_vs_code.txt to the printed digits.
%% sigma_res is zero at the mid-code point by construction (it is where the
%% constant is measured) and falls off the bottom of the log axis there.
%%
%% Panel (b) re-uses skill/pixmc_yield.py's calibrate()/passes()/SPECS verbatim
%% at its default seed 7, so its x = 0 point is yield_table.txt's "calibrated
%% (real)" row (9.0 %) and its grid-point-6 value is the "calibrated + window"
%% row (61.5 % per channel, 14.31 % over four, 96.5 % on the tracking spec).
%% x is the commanded potential trimmed off EACH end on the sweep's own
%% 128-code grid; the 0.4-2.1 V working window selects grid points 6..25, i.e.
%% 0.469-1.952 V, which is exactly the set pixmc_yield.py scores as re_err_trm.
%% Every yield here is an UPPER bound -- four devices in the amplifier
%% instantiate raw model cards and carry no mismatch (sss:pixtop-open).
%% Library cell names are kept out of the rendered prose.
\providecommand{\MaestroRoot}{include/analog/}
\providecolor{mtxA}{HTML}{2A78D6}
\providecolor{mtxB}{HTML}{EB6834}
\providecolor{mtxC}{HTML}{1BAF7A}
\begin{figure}[htbp]
  \centering
  {\pgfplotsset{compat=1.18}%
  \pgfplotsset{
    mcax/.style={
      width=0.42\linewidth, height=5.6cm,
      grid=both,
      major grid style={line width=0.2pt, draw=black!14},
      minor grid style={line width=0.2pt, draw=black!7},
      minor tick num=1,
      axis line style={line width=0.4pt, draw=black!45},
      tick style={line width=0.4pt, draw=black!45},
      tick align=outside,
      label style={font=\small}, tick label style={font=\small},
      legend style={font=\footnotesize, draw=none, fill=white,
                    fill opacity=0.85, text opacity=1, row sep=1pt},
      every axis plot/.append style={line join=round, line width=0.9pt,
                                     mark size=1.5pt},
    }}%
  \begin{tikzpicture}
    \begin{axis}[mcax, ymode=log,
      xlabel={Commanded potential~[\si{\volt}]},
      ylabel={\(\sigma\) of tracking error~[\si{\milli\volt}]},
      xmin=-0.1, xmax=2.52, xtick={0,0.5,1,1.5,2,2.5},
      ymin=0.03, ymax=400,
      legend style={at={(0.5,0.97)}, anchor=north},
    ]
      \addplot[mtxA, mark=*] table {\MaestroRoot data/pixmc_track_sigma.dat};
      \addlegendentry{as measured}
      \addplot[mtxB, mark=triangle*]
        table {\MaestroRoot data/pixmc_track_res.dat};
      \addlegendentry{after one stored constant}
      \draw[black!45, densely dashed, line width=0.6pt]
        (axis cs:0.4,0.03) -- (axis cs:0.4,400);
      \draw[black!45, densely dashed, line width=0.6pt]
        (axis cs:2.1,0.03) -- (axis cs:2.1,400);
    \end{axis}
  \end{tikzpicture}\hfill
  \begin{tikzpicture}
    \begin{axis}[mcax,
      xlabel={Potential excluded at each end~[\si{\volt}]},
      ylabel={Yield~[\si{\percent}]},
      xmin=-0.03, xmax=0.66, xtick={0,0.2,0.4,0.6},
      ymin=0, ymax=105, ytick={0,20,40,60,80,100},
      legend style={at={(0.03,0.97)}, anchor=north west},
    ]
      \addplot[mtxA, mark=*]
        table {\MaestroRoot data/pixmc_yield_track.dat};
      \addlegendentry{tracking specification}
      \addplot[mtxB, mark=triangle*]
        table {\MaestroRoot data/pixmc_yield_chan.dat};
      \addlegendentry{one channel, all 20 specs}
      \addplot[mtxC, mark=square*, dashed]
        table {\MaestroRoot data/pixmc_yield_4ch.dat};
      \addlegendentry{four channels}
      \draw[black!45, densely dashed, line width=0.6pt]
        (axis cs:0.469,0) -- (axis cs:0.469,105);
    \end{axis}
  \end{tikzpicture}}
  \caption[{Monte Carlo tracking error against commanded potential, and yield
  against the specification window.}]{Reference tracking on the assembled
  channel, 200 samples, process and mismatch, nominal process point,
  \SI{37}{\celsius}. Left: the spread of \(V_{\mathrm{RE}} -
  V_{\mathrm{PATTERN}}\) against the commanded potential, as measured and after
  the same sample's mid-code value is subtracted --- the residual one stored
  per-channel constant leaves behind, zero at mid-code by construction and off
  the axis there. Dashed lines mark the specified window;
  interior to it the error is a
  pure constant to \SI{0.36}{\milli\volt}, while outside it A1's input pair is
  leaving its common-mode range and no constant reaches the remainder. Right: yield
  against how much commanded potential the \SI{2}{\milli\volt} tracking
  specification excludes at each end, after that calibration, on the sweep's
  128-code grid; the \SIrange{0.4}{2.1}{\volt} window selects the points from
  \num{0.469} to \SI{1.952}{\volt} and is the dashed line. Every yield is an
  upper bound
  (Section~\ref{sss:pixtop-open}), and the per-channel curve saturates at
  \SI{62}{\percent} because the reference-electrode floor of
  Table~\ref{tab:pixtop-mc}, not tracking, then binds.}
  \label{fig:pixtop-mc}
\end{figure}

% % HAND-WRITTEN fragment -- NOT generated by maestro2tex.
% Source: ~/vestarv/docs/afe_rev2_pixel_verify.html rev. 2026-08-23, S11.2/S11.3
% (not verified, and the ordered next steps), S4 (F-3, F-4, F-5b, F-5c) and the
% vcm-tap finding at the end of S7.5:
%   with the common-mode slot selected, V_CM dips to 0.731 V and V_OUT is thrown
%   from 1.25 V to 16 mV, recovering over more than 4 us, at every conversion;
%   with the output slot selected V_CM moves 0.74 mV and V_OUT 3.5 mV over the
%   whole 21 us run.
%
% REVISED 2026-08-23 after the C_F fix and the post-fix characterisation.
% REMOVED: the stability bullet -- measured, sss:pixtop-stability, 24 grid points
%   and 5 corners (simruns/postfix_stb, simruns/postfix_corner_stb).
% REMOVED: "Noise, power, corners. None." -- all three now measured
%   (sss:pixtop-noise, sss:pixtop-power; simruns/postfix_{noise,decim,pwr,corner}).
% REVISED AGAIN 2026-08-24 after simruns/pixmc_20260823 (the first Monte Carlo
%   on this cell, 200 samples, process + mismatch, section castalia_pm_25).
% REMOVED: the Monte Carlo bullet -- run, sss:pixtop-mc, tab:pixtop-mc,
%   fig:pixtop-mc.  Both of its blockers are gone: the three mimcap_sin in
%   anatop_tia carry mismatchflag=1 (owner's edit 2026-08-23) and the deck was
%   spliced onto castalia_pm_25 rather than the typical-only castalia_typ_25
%   the setup names.
% ADDED: the four raw-model devices in the amplifier (anatop_tia_amp_ab M48/M49
%   nch_25, M57/M58 pch_25 -- present in BOTH A1 and A2), which drew no
%   mismatch, so the sss:pixtop-mc spreads are narrower than the built
%   channel's.  They cannot be fixed from a splice: they are schematic
%   instances.  README.md "Owner-fixes" of simruns/pixmc_20260823 lists the rest
%   of the census (biasgen_local, pixel_dac_psu, amux_2, the sample switch and
%   level shifters), all of them in blocks the ablation shows contribute
%   nothing.
% CHANGED: the Wiring bullet's common-mode item now states what it costs --
%   it is the tap the calibration of sss:pixtop-mc is measured through, so no
%   calibrated figure there can be realised until it is buffered.
%
% 2026-08-24, OWNER INSTRUCTION "do not talk about pixel level yield yet":
%   the two yield sentences here were reduced to the facts they encoded, with
%   no per-channel or four-channel percentage.  The capacitor clause added to
%   the first bullet is CONFIRMED from the run's own netlist,
%   simruns/pixmc_20260823/decks/input.scs.mc_dc: FOUR mimcap_sin carry an
%   explicit mismatchflag=0 -- CBP, CBPC, CBNC and CBN in anatop_biasgen_local,
%   the bias-rail decoupling capacitors.  There are not six, and they are not
%   in the pixel signal path; the three mimcap_sin in anatop_tia (C_F_TIA,
%   C_CMP1_TIA, C_CMP2_TIA) are at mismatchflag=1 after the owner's 2026-08-23
%   edit.  mismatch_census.txt agrees: "mismatchflag=0 (explicit) : 4 devices".
% ADDED: the readout path (no run has converted a real electrode current through
%   the on-chip converter -- sss:pixtop-cv reads currents from an in-line probe
%   and from Zt, not from a code), and the layout parasitics on the new
%   capacitor's node (4 x 52.09 um mimcap_sin on the TIA summing node, 10 853
%   um^2; no extracted view of anatop_tia exists).
% F-4 (the two DAC supply regulators still placed) is now quantified in
%   sss:pixtop-power and stays in the Wiring bullet as an open item.
% Library cell names are kept out of the rendered prose.

\subsubsection{Not Yet Verified}\label{sss:pixtop-open}

\begin{itemize}[nosep]
  \item \textbf{Devices outside the statistics.} Two n-channel and two
    p-channel devices in the amplifier's enable and class-AB bias network
    instantiate raw model cards rather than the mismatch-enabled wrappers, so
    the Monte Carlo of Section~\ref{sss:pixtop-mc} drew no mismatch on them.
    They sit in both A1 and A2, the blocks that set every distribution there,
    and four more mismatching devices per amplifier can only widen a spread ---
    so every \(\sigma\) in Table~\ref{tab:pixtop-mc} is narrower than the built
    channel's. Those four devices have since been given the mismatch-enabled
    wrappers, but the run tabulated here predates that edit. Six capacitors
    added to the channel in a later edit stand outside the statistics too, and
    four more on the local bias generator's rails are flagged out by intent:
    they decouple, they do not match. The run must be repeated against the
    channel as it now stands.
  \item \textbf{The converter at schematic level.} Its capacitor array exists
    only in the extracted view (Section~\ref{sss:pixtop-adc}), so the code the
    schematic tests report is stuck and is deliberately left unspecced.
  \item \textbf{Current read as a code.} No run has taken a real electrode
    current through the on-chip converter. The voltammogram of
    Section~\ref{sss:pixtop-cv} reads current from an in-line probe and from
    the measured transimpedance; the converter figures of
    Section~\ref{sss:pixtop-adc} come from known voltages applied at the
    multiplexer. The two have not been closed into one measurement.
  \item \textbf{Layout parasitics on the compensation node.} The
    \SI{10853}{\micro\meter\squared} of capacitor added by
    Section~\ref{sss:pixtop-stability} sits on the transimpedance summing node,
    the one node in the channel where stray capacitance costs phase margin. No
    stability or noise figure here carries its wiring parasitics: the stage has
    no extracted view.
  \item \textbf{Wiring.} Three items are drawn but not connected. The sixteen
    test-port multiplexer inputs are undriven, so the debug path selects
    between floating nodes on demand. The two DAC supply blocks' outputs reach
    nothing, which is why the supply-modulation defect of
    Section~\ref{sss:dacr2r12-supply} is not exercised in this cell and why
    \SI{806.8}{\micro\watt} of Section~\ref{sss:pixtop-power} buys nothing. And
    there is no buffer between the common-mode rail and the converter
    multiplexer: selecting that slot loads the bare ladder output --- which is
    also A2's reference --- pulling the common mode to \SI{0.731}{\volt} and
    throwing the transimpedance output from \SI{1.25}{\volt} to
    \SI{16}{\milli\volt} for more than \SI{4}{\micro\second}, at every
    conversion. It collapses the one measurement the gain-independent zero of
    Section~\ref{sss:pixtop-zero} depends on, and it is the tap the calibration
    of Section~\ref{sss:pixtop-mc} is measured through: until it is buffered,
    that calibration cannot be run on this cell. The output slot is quiet by
    comparison: \SI{0.74}{\milli\volt} on the common mode and
    \SI{3.5}{\milli\volt} on the output over a whole \SI{21}{\micro\second}
    run.
\end{itemize}

% % HAND-WRITTEN fragment -- NOT generated by maestro2tex. Input from
% AnalogChapter.tex immediately after the assembled-channel fragments.
%
% Unlike the other note_meas_* fragments, part of this fabric EXISTS in the
% schematic: the converter's 16:1 input multiplexer and its 4-bit select pin
% (slot 0 = transimpedance output, 1 = common mode, 2 = commanded potential,
% 3 = reserved for a buffered RE tap and floating, 4 = analog test point,
% 5 = ground). What does not exist: the RE buffer (F-5c), the common-mode
% buffer (the S7.5 finding), the test-port mux inputs (F-3), the electrode
% isolation switches and the on-chip gain reference resistor.
% Numbers: bit-9 inversion and 88.5 nA/LSB at N = 0 (sss:pixtop-adc); 0.3 nV
% zero spread over 320 codes and 0.825 nA/LSB at N = 319 (sss:pixtop-zero);
% 0.731 V / 16 mV collapse of the common-mode slot (sss:pixtop-open).
% Library cell names are kept out of the rendered prose.

\subsubsection{Measuring this block on chip}

The channel's own converter multiplexer carries this sequence; only the last
step needs the tester.
\begin{enumerate}[nosep]
\item Invert bit~9 of the raw converter bus before anything else reads it. Every
  step below is that number.
\item Select the transimpedance output (slot 0) with the working electrode
  isolated and take the current zero at a \emph{high-gain} code. The zero is
  gain-independent to \SI{0.3}{\nano\volt} over all 320 taps, so one
  measurement calibrates every code --- but at the bottom code one LSB is
  \SI{88.5}{\nano\ampere} against \SI{34}{\nano\ampere} of offset and the
  measurement is quantisation noise.
\item Do not take that zero through the common-mode slot until that tap is
  buffered: sampling it pulls the common mode to \SI{0.731}{\volt} and the
  output to \SI{16}{\milli\volt} at every conversion, and the channel never
  settles.
\item Read the commanded potential at slot 2 to separate a DAC fault from an
  amplifier fault. A1's own contribution is not observable: slot 3 is reserved
  for a buffered reference-electrode tap that does not exist, and the raw
  electrode must not be sampled.
\item Absolute potential scale and transimpedance gain still need the
  production tester at the reference-DAC pad and the working-electrode pad
  (Sections~\ref{ss:dacr2r12} and~\ref{ss:tiarprog}). The in-loop feedback
  resistance of Section~\ref{sss:pixtop-cv} has no on-die counterpart --- it is
  read from a simulation probe.
\end{enumerate}

% \FloatBarrier
% End of the commented-out assembled-channel section.
\FloatBarrier

% Impedance spectroscopy, hand-written throughout from
% ~/chips/castalia/simruns/eisdemo_20260825/README.md and the run data beside it (30 AC
% runs, 24 fixed-ratio runs, 17 real-DAC transients, an extracted-converter probe, and
% the bench cell castalia/tb_anatop_eisdemo, whose four Maestro tests verify headlessly
% at 13/13 specs), plus eiscorner_20260825/README.md (the same bench and extractor with
% a corner and temperature axis: 210 corner AC runs, 42 corner band-edge runs, 24 runs
% with the poly and mimcap cornered too, eight 200-sample Monte Carlo runs, and eis_ac
% over seven corners at 35/35 specs).
% It is a subsection rather than a block section because it is a measurement mode of
% the assembled channel and adds no hardware, while spanning the reference DAC, both
% amplifiers, the programmable feedback resistor, the converter and the firmware
% demodulator.
% The section is two prose fragments and the inline table tab_eis_codelevel.
% note_eis_{coverage,excitation,open,rules}.tex and
% fig_eis_{error,settling,mccal,nyquist,codelevel}.tex stay on disk un-input, each with
% its own provenance header; restoring one is putting an \input back. No \ref
% elsewhere in the chapter reaches them. The coverage table's Monte Carlo tick for the
% impedance row stands on the sigma = 11.1 % of note_eis_intro's pinned-passive
% paragraph, the section's only wafer number. Campaigns:
% ~/chips/castalia/simruns/{eisdemo,eiscorner}_20260825, eisdeep_20260826 and
% eisfigs2_20260826; narrative ~/vestarv/docs/afe_rev2_eis_lite.html.
% VestaRV: hand-written TRM fragment, not maestro2tex output. Input from
% AnalogChapter.tex, first of the impedance section.
%
% Sources for every number here. ~/chips/castalia/simruns/eisdemo_20260825/README.md:
% S3 the measurement, S5 coverage; engine skill/eisdemo_engine.py, extractor
% skill/eisdemo_extract.py, every run a textual splice of the pixel's own
% Maestro-netlisted deck simruns/mcflags_20260825/canon/input.scs.
% eiscorner_20260825/README.md S2.5 (ac_resmim/resmim_table.txt) is the -33.4 % /
% +35.5 % on R_f.  eisdeep_20260826/README.md: S3 the grid-density argument (10
% points/decade puts the chord polygon 0.651 % of the arc radius off the true arc
% against a 0.30 % instrument error; 20 gives 0.165 %, 30 gives 0.074 %), S5 the
% cornered-poly band edge, S7 the Monte Carlo (sigma = 11.1 % on R_f, n = 654, exact
% chi2 95 % interval).  Monte Carlo throughout: castalia_pm_25, nominal corner, 37 C,
% seed 12345.  The corner spread of the reported current, -33.26 % (ff_all) to
% +35.71 % (ss_all) on the nominal R_f law against a true current constant to 0.12 %
% and a settled impedance constant to 0.025 %, with both MOS-skew corners showing
% none of it (+0.46 / +0.56 %), is ~/chips/castalia/simarchive/tb_anatop_eissq/
% 20260905_sq_step/typ_10k_corners (five executed netlists, each with its own
% section=castalia_*).
%
% The small-signal runs of this section carried no 22 pF feedback capacitor, and the
% conditions paragraph says so: without that sentence the phase numbers read as
% channel measurements, which they are not.
% The pinned-passive mechanism is stated once in AnalogChapter.tex (label
% p:pinned-passives); keep only the pointer and this section's own numbers here.
% Library cell names stay out of the rendered prose.
%
% The section is two prose fragments (this one and note_eis_accuracy) and the inline
% table tab_eis_codelevel.  note_eis_rules.tex, note_eis_coverage.tex,
% note_eis_excitation.tex, note_eis_open.tex, fig_eis_nyquist.tex, fig_eis_error.tex,
% fig_eis_mccal.tex and fig_eis_settling.tex stay on disk un-input, with their own
% headers and source lists intact; restoring them is putting their \input lines back.
% Campaign narrative: ~/vestarv/docs/afe_rev2_eis_lite.html.

\subsection{Impedance Spectroscopy}\label{ss:eis}

The channel measures impedance with no impedance hardware. A small excitation
rides on the reference-DAC setpoint of Section~\ref{ss:dacr2r12}, A1 applies it
at the reference electrode, A2 and the converter read the resulting current
through the same programmable feedback resistor a dc measurement uses
(Section~\ref{ss:tiarprog}), and firmware demodulates the two records in
quadrature, an I/Q multiply--accumulate over whole excitation
periods. The excitation is a toggle between two DAC codes, so the
mode adds a clock divider and an accumulate to the channel and nothing else.

Everything below runs on the netlist of that complete channel at a
\SI{2.5}{\volt} analog supply, and is schematic-level except where the converter
appears, which it does as its extracted view. The nominal point is the typical
process corner at \SI{37}{\celsius}; only the small-signal measurement was
repeated over the chapter's four process corners at \SI{25}{\celsius}, over
\SIrange{0}{85}{\celsius}, and over Monte Carlo. The electrodes are the linear
cell model of Section~\ref{ss:analog-models}, \(R_u = \SI{1}{\kilo\ohm}\)
throughout, with \(R_{\mathrm{ct}} \parallel C_{\mathrm{dl}}\) of
\SI{1}{\mega\ohm} \(\parallel\) \SI{100}{\nano\farad},
\SI{100}{\kilo\ohm} \(\parallel\) \SI{1}{\micro\farad} and
\SI{10}{\kilo\ohm} \(\parallel\) \SI{10}{\micro\farad}. Sweeps carry \num{20}
points per decade. \textbf{The small-signal measurements were taken on a
replica that carries no feedback capacitor}, so their phase numbers and the
band edge that follows from them are optimistic; the transient campaign runs
through the shipping channel with the \SI{22}{\pico\farad} present, and the
accuracy note states both.

\textbf{The pinned-passive caveat of page~\pageref{p:pinned-passives} governs
every corner number in this section.} Neither the ladder that is \(R_f\) nor the
capacitor that sets the feedback pole corners here. The Monte
Carlo does vary the poly, and measures \(\sigma = \SI{11.1}{\percent}\) on
\(R_f\) with process and mismatch together at the nominal corner. Across the
five-corner set the current the channel reports through the nominal feedback
law spans \SIrange{-33.3}{+35.7}{\percent} while the current that actually
flows is constant to \SI{0.12}{\percent} and the impedance derived from it to
\SI{0.025}{\percent}, all of it the poly sheet: the ratiometric calibration is
not a refinement of this mode but a condition of it.

% HAND-WRITTEN fragment -- NOT generated by maestro2tex. Input from
% AnalogChapter.tex after note_eis_intro.tex.
%
% ======================================================================
% CUT 2026-08-26, OWNER'S EXPLICIT INSTRUCTION (see note_eis_intro's header for
% the scope of the cut).  This fragment absorbed what survives of the dropped
% note_eis_coverage.tex and carries the section's headline, code-level and
% calibration results in three paragraphs.  It lost its \subsubsection and its
% label sss:eis-accuracy.
%
% WHAT WAS CUT FROM THIS FRAGMENT, AND WHERE IT IS:
%   - the per-load / per-gain-code error table in prose (fresh, typical, fouled
%     against six codes at six frequencies): eisdemo_20260825
%     ac/report_architectures.txt, ac/ac_results.json (README S3.2-S3.4), and
%     for the corner axis eiscorner_20260825 corner_tables.txt S2.
%   - the ordered error budget (feedback capacitor, RE tap, calibration-impedance
%     mismatch, loop error vs noise gain 0.120/0.330/0.660/1.322 % at
%     R_f/|Z| = 7.3/20/40/80): same files, README S3.3-S3.5.
%   - the whole per-corner accuracy paragraph (0.023-0.041 % on the calibration
%     resistor; 0.70-1.40 % process and 0.88-1.46 % temperature on typical/N=20
%     at 1 kHz; phase 4.57-4.79 deg): eiscorner_20260825 corner_tables.txt S2.
%   - the coverage table tab:eis-coverage and the band-edge table
%     tab:eis-band-corner, with their captions: they are in the dropped
%     note_eis_coverage.tex, still on disk, sourced from eisdemo_20260825
%     coverage_map.txt and eiscorner_20260825 band/band_results.json.
%   - fig_eis_error.tex (error against frequency per load and gain code) is on
%     disk and no longer input.
%   - the full Monte Carlo architecture comparison (A / B_law / B_dc / C at
%     seven frequencies on four load-and-code configurations, with chi2
%     intervals): eisdeep_20260826 mc/, mc_report/ (README S7), superseding
%     eiscorner_20260825 mc_summary.txt S3.
%   - narrative summary: ~/vestarv/docs/afe_rev2_eis_lite.html.
% NOTHING MEASURED IS LOST.
%
% SECOND CUT 2026-08-26, OWNER'S EXPLICIT INSTRUCTION: "about two pages" (see
% note_eis_intro's header for the scope).  This fragment is now two paragraphs
% and one figure.
%
% WHAT THIS SECOND PASS REMOVED FROM THIS FRAGMENT, AND WHERE IT IS:
%   - THE WHOLE WAFER-CALIBRATION PARAGRAPH ("Across a wafer the spread is the
%     feedback resistance") and with it \input fig_eis_mccal.tex, which is left
%     on disk unedited and is no longer input by anything.  Everything that
%     paragraph carried leaves the rendered prose together, so no number in it
%     is left under-qualified: the process/mismatch/both decomposition
%     (10.91 / 0.43 / 11.12 % over 500/500/654 samples), the stored-law cost
%     (sigma = 10.2 % of magnitude, that run's own -32.1 .. +30.6 % over 649
%     dies, on top of the law's +9 % nominal offset at N = 0), the per-chip dc
%     calibration residue (sigma = 0.0006 %), the phase-escapes-the-argument
%     clause, the above-the-pole ratiometric comparison (6.85 +- 1.88 % against
%     2.51 +- 0.83 % on the fresh electrode at N = 63 and 5 kHz) and the
%     switchable-calibration-impedance requirement (0.0002 % matched against
%     0.18 % at 9.2x away).  Sources are unchanged and listed below:
%     ~/chips/castalia/simruns/eisdeep_20260826 S7 (mc/, mc_report/),
%     eiscorner_20260825 mc_summary.txt S2-S4, eisdemo_20260825
%     coverage_map.txt, the figure data in eisfigs2_20260826, and the narrative
%     ~/vestarv/docs/afe_rev2_eis_lite.html.  The section keeps ONE wafer
%     number, the sigma = 11.1 % of note_eis_intro's pinned-passive paragraph.
%     The per-chip gain constant this paragraph justified is a live requirement
%     elsewhere in the chapter: it is a row of tab_calibration_constants.
%   - the corner sentence of the lead paragraph ("Across the corner set the
%     error spans a factor of two, worst at 85 C on every load and gain code,
%     while phase moves far less, because this corner set does not move the
%     poles"): eiscorner_20260825 corner_tables.txt S2.  The corner claim the
%     section still makes is the band edge in the same paragraph, which
%     note_eis_intro's pinned-passive paragraph qualifies.
%   - the 0.030 % / 0.010 deg accuracy tier at and below 1 Hz -- below the
%     0.1 Hz usable band the paragraph states: eisdeep_20260826 S3.1
%     accuracy_envelope.txt.
%   - the closing sentence of the code-level paragraph ("One record, one die,
%     one frequency, one resistive load: a deterministic attribution, not a
%     distribution").  THE CAVEAT IS NOT DROPPED: tab_eis_codelevel's caption
%     carries it in bold, in a stronger form ("with no noise sources in the
%     deck"), and the sentence that quotes +0.343 % points at that table.
%     Do not delete it from the caption.
%   - "on a two-level excitation" was KEPT when the phase clause was tightened:
%     180*f/f_s is the square wave's term and a sinusoid on the identical grid
%     returns 0.0000 deg (eisdeep_20260826 S2.4), so the clause is a qualifier
%     and not a flourish.
% NOTHING MEASURED IS LOST BY THIS PASS EITHER.
%
% 2026-08-27, OWNER'S EXPLICIT INSTRUCTION: the code-level FIGURE became a
% TABLE.  This fragment now inputs tab_eis_codelevel.tex and its code-level
% paragraph points at Table~\ref{tab:eis-codelevel}; fig_eis_codelevel.tex is
% left ON DISK UNEDITED and is no longer input by anything.  The figure printed
% all ten of its scalars on its own face and cost a full page; the table states
% them, plus the +-0.573 deg gate and the -180 f/f_s offset the figure showed
% spatially, in a fraction of the space.  NOTHING MEASURED IS LOST.
%
% ----------------------------------------------------------------------
% THIRD CUT 2026-08-27, OWNER'S EXPLICIT INSTRUCTION: "two pages or fewer".
% THE SECTION'S LAST FIGURE IS GONE.  fig_eis_nyquist.tex is no longer input by
% anything and is left ON DISK UNEDITED apart from its own header note; the
% archive that holds it and its data is
% ~/chips/castalia/simruns/eisfigs2_20260826.  Its \input and its
% Figure~\ref were both in note_eis_intro.tex, whose header carries the
% per-fact account of what that caption held and why none of it moved into the
% prose.  Two of those facts are THIS fragment's business, and both were
% checked against the text below before being cut:
%   - the caption's 0.44 % / 1.03 deg below 20 kHz on three loads at gain code
%     0 is the same claim as this fragment's opening headline, which is
%     broader (all five loads, gain-ranged to the lowest code in compliance,
%     0.44 % / 1.04 deg, 0.01 Hz - 20 kHz).  The headline stands unchanged; the
%     caption's narrower restatement went.
%   - the caption's "gain code 0 covers 236 ohm to 150.9 kOhm" is 236.01 ohm x
%     the 639.3x = 2.81 decades per code that the coverage sentence below
%     already states.  Derived, so cut; the coverage headline is unchanged.
% NOTHING IN THIS FRAGMENT'S RENDERED PROSE CHANGED IN THIS PASS.  The prose
% cut of this pass fell entirely in note_eis_intro (the load names, which the
% figure's panel titles were the only consumer of, and one redundant
% band-spread clause) -- see its header.  This fragment's band sentence is what
% made that clause redundant, so do not cut "with the poly and the capacitor
% cornered as well it falls to 2.5 kHz on a worst-case die" without restoring
% something in its place.
% ----------------------------------------------------------------------
% ======================================================================
%
% SOURCE OF EVERY SURVIVING NUMBER:
%   eisdemo_20260825 ac/report_architectures.txt -- pureR 11.000 kOhm at N = 0:
%     +0.0294 % / 0.0000 deg, flat 0.1 Hz - 1 kHz, 0.0301 % at 5 kHz.
%   eisdemo_20260825 coverage_map.txt (README S5) -- the window: union
%     236.01 ohm - 14.789 MOhm = 4.80 decades, each code spanning 639.3x =
%     2.81 decades; |Z|min is where the square's PEAK current fills the
%     +-0.80 V compliance, |Z|max where the FUNDAMENTAL of the output falls to
%     one converter LSB (1.5934 mV).
%   eisdemo_20260825 ratio/band_edges.txt + eiscorner_20260825 corner_tables.txt
%     S3-S4 -- the phase gate is the binding one: f(0.573 deg) 5527 Hz nominal
%     at N = 0, 5290 Hz worst (ss), against a magnitude gate never below
%     8369 Hz.  Hence "0.1 Hz to 5.3 kHz, set by phase".
%   eisdeep_20260826 S5 (band_ff40/, band_ss40/) -- with the poly and the mimcap
%     cornered too and the ratio held at 40, that phase gate runs 3398 Hz (ss)
%     to 10.36 kHz (ff); with the calibration resistor sized from the NOMINAL
%     R_f it runs 2523 Hz to 15.36 kHz.  The 2.52 kHz quoted is that worst case,
%     and the ff/ss passive sections are about +-3 sigma of the 11.1 % process
%     spread, so it is a worst-case die and not a typical one.
%   eisdeep_20260826 S3.1 accuracy_envelope.txt -- best gain code in compliance
%     at each frequency, no calibration, architecture A, 20 points/decade:
%     0.44 % / 1.04 deg over 0.01 Hz - 20 kHz on all five loads; 0.33 % /
%     0.29 deg at and below 5 kHz; 0.030 % / 0.010 deg at and below 1 Hz.
%     "Ranging badly costs an order of magnitude": worst code still in
%     compliance reads +0.43 % on the calibration resistor at dc.
%   eiscorner_20260825 corner_tables.txt S2 -- the factor of two across corners,
%     85 C worst on magnitude for every load and code, ss best; phase moves far
%     less because this corner set does not move the poles.
%   eisdeep_20260826 S2, code_report/codeA_re_n20.json, csv/codelevel.csv --
%     256 conversions of the extracted converter (av_extracted, 4451 devices) at
%     869 565.2 S/s, 64 per excitation period -> f_exc = 13 586.9565 Hz exactly;
%     excitation from the real 12-bit ladder (code pair 2032/2064); both
%     channels through the on-chip mux, interleaved slot 0 (V_OUT) / slot 6 (RE);
%     load 11.000 kOhm at N = 20; max |V_OUT - vcm| 0.132 V of the +-0.80 V
%     window.  AC analysis of the same load and code: 11 032.43 ohm / +0.3603 deg.
%       Z_analog  +0.1262 % / -0.0127 deg
%       Z_sampled +0.9812 % / -2.8949 deg
%       Z_code    +0.3431 % / -2.9383 deg     (steps +0.126 / +0.855 / -0.638 pp)
%     Inter-channel gain mismatch +0.6359 % (628.00 codes/V on V_OUT, fundamental
%     92.7 LSB; 624.04 on RE, 7.7 LSB) predicts -0.632 pp against -0.638
%     measured -> 0.006 pp for quantisation, dither and aperture together.
%   eisdeep_20260826 S2.4 code_report/deskew_test.txt -- the -2.94 deg is the
%     excitation grid: the two slots sample T_conv/2 apart -> 180*f/f_s =
%     2.8125 deg.  Pure sine on the same interleaved grid 0.0000 % / 0.0000 deg.
%     +-0.573 deg caps an uncorrected interleaved read at 2768 Hz = 0.00318*f_s.
%   eisdeep_20260826 S7 mc/, mc_report/ -- R_f at N = 0, castalia_pm_25, nominal
%     corner, 37 C, seed 12345, exact chi2 95 % intervals:
%       process only  10.9074 % (n = 500)   mismatch only 0.4262 % (n = 500)
%       both          11.1190 % (n = 654)   quadrature sum 10.916 %
%     typ_n0 (n = 649) vs the 11.0 k standard, 0.1 Hz: A/B_dc 0.0033 +- 0.0006,
%       B_law -7.6271 +- 10.2423 with that run's OWN range -32.0596 .. +30.6335,
%       C -0.0234 +- 0.1791.  Quote ONE run's own sample range, never a union.
%     typ_n0 vs the MATCHED 100.8 k standard (n = 643), 0.1 Hz: C -0.0000 +-
%       0.0002 -- the "sigma = 0.0002 %" below.
%     fresh_n63 (n = 652) at 5 kHz: A 2.9948 +- 0.9646, B_dc 6.8508 +- 1.8784,
%       C 2.5124 +- 0.8301.
%     The stored-law sigma on measured magnitude, 10.2 % at N = 0, and the
%     +9.06 % nominal offset of the law at N = 0, are eisdemo_20260825
%     coverage_map.txt with eiscorner_20260825 mc_summary.txt S2/S3.
%     B_law and B_dc report BYTE-IDENTICAL phase in every table -- R_f is a real
%     scale factor and drops out of phase entirely (mc_summary.txt S4.5).
% The four firmware architectures of the source campaign are named in prose
% rather than lettered: A = "reading the potential at the reference electrode and
% the current at the transimpedance output"; B = "a stored feedback resistance";
% C = "ratiometric against a known impedance measured at the same frequency";
% D = "assuming the commanded potential reaches the electrode".
% Library cell names are kept out of the rendered prose.
%
% ----------------------------------------------------------------------
% CORRECTED 2026-09-05.  EVERY SMALL-SIGNAL NUMBER ABOVE WAS MEASURED WITHOUT
% THE 22 pF FEEDBACK CAPACITOR.  The replica bench castalia/tb_anatop_eisdemo
% had no C_F instance until 2026-09-05; the shipping transimpedance stage has
% one, and it dominates the phase.  Re-running eis_ac and eis_stb_tia on that
% bench with the capacitor added, nominal corner, everything else identical
% (before -> after, ~/chips/castalia/simruns/eisdemo_20260825/maestro_verify
% against the 2026-09-05 re-run, F17 report section 2):
%     err_1hz    0.003225 % -> 0.003225 %      (unchanged)
%     err_1k     0.30228 %  -> 0.30295 %       (+0.0007 pp)
%     err_max_5k 0.31074 %  -> 0.32640 %       (+0.016 pp)
%     ph_max_5k  0.26485 deg -> 1.04885 deg    (x3.96)
%     aol_db     79.574 dB  -> 79.574 dB       (unchanged)
%     pm         92.299 deg -> 83.988 deg
%     gm         28.020 dB  -> 9.761 dB        (the bench's own spec is > 10 dB)
%     fx         1.0642 MHz -> 22.154 MHz      (the loop now crosses at the
%                                               amplifier's own bandwidth,
%                                               because the feedback factor
%                                               tends to 1 above the pole)
% CONSEQUENCE FOR THIS FRAGMENT: the magnitude headline is unaffected (0.311 ->
% 0.326 % at 5 kHz).  The PHASE headline and the band edge are not: the 0.573
% deg gate is crossed a factor of four earlier than the section said, so the
% uncorrected band edge is well below 5.3 kHz and not at it.  The sentence in
% the prose now states the condition and the correction rather than the number
% alone.  NOTHING IS DELETED: the 0.44 % / 1.04 deg and 0.33 % / 0.29 deg
% envelopes stand as measured, with the condition attached.
%
% ADDED IN THE SAME PASS: the paragraph "Measured through the complete channel"
% and tab_eis_pixel.tex, from the square-wave campaign through the real pixel
% (six Maestro tests, 26 completed runs, F17 report,
% ~/chips/castalia/simarchive/tb_anatop_eissq/).  What each number is:
%   settled |Z| within 0.0068-0.15 % of the model on every non-railing load and
%     gain code, 0.0147 % on the typical load (20260905_sq_step/
%     table_sq_step.txt, 20260905_sq_rail/table_sq_rail.txt), and within
%     0.025 % across the five-corner set (20260905_sq_step/typ_10k_corners).
%   demodulator 0.01 % / 0.01 deg: err(z_cell) on the calibration load,
%     0.002-0.007 % and under 0.06 deg at every frequency
%     (20260905_sq_freq/demod_res.txt).
%   converter +0.160 % / 0.063 deg and 0.305 LSB rms code residual
%     (20260905_sq_code/table_sq_code.txt, codes_main.json); the 16-conversion
%     control at twice the frequency gives +0.195 % / -0.150 deg.
%   extraction within 0.15 % at 1 kHz and 0.40 % at 5 kHz, ratio-calibrated,
%     R_s supplied (20260905_sq_freq/extract_variants.txt, rows "cal true").
%     With R_s taken as Re(Z) at the top sweep frequency instead, that
%     estimator alone carries a +9.92 % structural bias on this load and the
%     extracted R_ct and C_dl degrade to -2.19 % and +4.86 %.
% WHY THE -2.94 deg OF tab:eis-codelevel IS NOT CONTRADICTED: that record
% interleaved two multiplexer slots, and the offset is the 180 f/f_s grid term
% the table already attributes.  The square-wave record holds the slot static,
% so the term is absent and the residual 0.063 deg is the converter alone.
% ----------------------------------------------------------------------

\textbf{The error is set by the noise gain \(1 + R_f/|Z|\) at the frequency
being measured, not by the electrode's chemistry.} Gain-ranged to the lowest
code in compliance and with no calibration applied, the channel holds
\SI{0.44}{\percent} of magnitude and \SI{1.04}{\degree} of phase from
\SI{0.01}{\hertz} to \SI{20}{\kilo\hertz} on all five loads, and
\SI{0.33}{\percent} and \SI{0.29}{\degree} at and below \SI{5}{\kilo\hertz};
ranging badly costs an order of magnitude. The
320-tap ladder covers \SI{236}{\ohm} to \SI{14.79}{\mega\ohm}, \num{4.80}
decades, each gain code spanning \num{2.81} of them between the
\(\pm\SI{0.80}{\volt}\) output compliance and one converter code of output
fundamental; adjacent codes overlap by more than a decade, so there are no
impedance seams. The usable band is \SI{0.1}{\hertz} to \SI{5.3}{\kilo\hertz},
set by a \(\pm\SI{0.573}{\degree}\) phase gate at the feedback pole and not by
magnitude; with the poly and the capacitor cornered as well it falls to
\SI{2.5}{\kilo\hertz} on a worst-case die. \textbf{Every figure in this
paragraph is a small-signal measurement on a replica that carries no feedback
capacitor.} Adding the \SI{22}{\pico\farad} the shipping stage has, and
changing nothing else, leaves the magnitude envelope alone
(\num{0.311} to \SI{0.326}{\percent} at \SI{5}{\kilo\hertz}) and multiplies
the phase error by four, \SI{0.265}{\degree} to \SI{1.049}{\degree}: the
uncorrected band edge is therefore well below \SI{5.3}{\kilo\hertz}, not at
it. The same capacitor moves the transimpedance loop crossing from
\SI{1.06}{\mega\hertz} to the amplifier's own \SI{22.2}{\mega\hertz} and the
gain margin from \SI{28.0}{\decibel} to \SI{9.76}{\decibel}, against a
\SI{10}{\decibel} specification; phase margin \SI{84.0}{\degree} still
passes.

\paragraph{Measured through the converter.} An \SI{11.000}{\kilo\ohm} load at
gain code 20, excited by the real reference ladder and read through the input
multiplexer and an extracted converter, demodulates from \num{256} converter
codes to \(+0.343\,\%\) of magnitude and \SI{-2.94}{\degree} of phase against
the small-signal analysis of the same load and gain code
(Table~\ref{tab:eis-codelevel}). Three demodulations of that one record
attribute the magnitude error completely, and \textbf{the converter's whole cost
is one number}: a \(+0.636\,\%\) inter-channel gain mismatch, which two points
of gain calibration remove. The phase term is the sampling grid and not the
converter: one converter serving two channels samples them half a conversion
apart, costing \(180\,f/f_s\) on a two-level excitation: deterministic, so it
needs no calibration;
uncorrected, the \(\pm\SI{0.573}{\degree}\) gate caps an interleaved read at
\SI{2.77}{\kilo\hertz}.

%% HAND-WRITTEN table fragment (maestro2tex house style; nothing here is
%% generated -- one transient record decoded by hand, not a Maestro results
%% database, and there is no maestro2tex config for any *_eis_* fragment).
%%
%% 2026-08-27, OWNER'S EXPLICIT INSTRUCTION: THIS TABLE REPLACES THE FIGURE
%% fig_eis_codelevel.tex, which took a full page to print ten scalars on the
%% face of a chart.  That fragment is LEFT ON DISK UNEDITED and is no longer
%% input by anything -- restoring it is putting an \input back and moving the
%% \label.  Its header carries the same note and the full provenance.  Nothing
%% measured is lost: every plotted value is a row here, at the precision of the
%% source data rather than of a bar top.
%%
%% WHAT THE FIGURE CARRIED SPATIALLY AND THIS TABLE CARRIES OTHERWISE:
%%   - filled vs open markers (measured vs corrected vs analytic control) ->
%%     the Reading column.
%%   - the +-0.573 deg gate band -> the gate column.
%%   - the dashed -180 f/f_s line the two interleaved reads sit on -> a row of
%%     its own, adjacent to them, marked analytic.
%%
%% Source of every number: ~/chips/castalia/simruns/eisfigs2_20260826/
%%   data/fig5a_attribution.csv (magnitude block, column step_delta_pp) and
%%   data/fig5b_phase.csv (phase block, columns dphase_vs_ac_deg and kind),
%%   reduced from ~/chips/castalia/simruns/eisdeep_20260826/
%%   code_report/codeA_re_n20.json -- the headline record, 256 conversions,
%%     mux interleaved slot 0 (transimpedance output) <-> slot 6 (RE), 294 us,
%%     869 565.2 S/s, 64 conversions per excitation period -> f_exc =
%%     13 586.9565 Hz exactly; excitation from the real 12-bit ladder (code pair
%%     2032/2064); 11.000 kOhm load at N = 20; max |V_OUT - vcm| 0.132 V of the
%%     +-0.80 V window; extracted converter, 4451 device lines.
%%   code_report/pair_F0_F6.json -- the two-record scheme (codeF0_n20 +
%%     codeF6_n20), the "two records, one grid" row.
%%   code_report/deskew_test.txt, csv/codelevel.csv, sampling_prediction.txt.
%% AC analysis of the same load and gain code: 11 032.43 ohm / +0.3603 deg.
%%   Z_analog  +0.1262 % / -0.0127 deg   (20 ns strobe grid, boxcar-decimated,
%%                                        no sampling -- the transient path)
%%   Z_sampled +0.9812 % / -2.8949 deg   (same analog nets at the apertures)
%%   Z_code    +0.3431 % / -2.9383 deg   (the ten-bit words themselves)
%% Steps +0.126 / +0.855 / -0.638 pp.  Inter-channel gain mismatch +0.6359 %
%% (628.00 codes/V on the current channel, fundamental 92.79 LSB; 624.04 on the
%% potential channel, fundamental 7.65 LSB) predicts -0.632 pp against -0.638
%% measured -> 0.006 pp for quantisation, dither and aperture together.
%% Aperture fitted at nominal + 0 ns, 0.78 LSB rms; the three per-cycle code
%% phasors agree to 0.0000 %.  Half-conversion grid offset 180*f/f_s = 2.8125
%% deg; deskewed -0.082 deg (codes -0.126); two records on one grid +0.083 deg
%% measured; pure sine on the same grids 0.000 deg.  +-0.573 deg caps an
%% uncorrected interleaved read at 2768 Hz = 0.00318*f_s.
%%
%% THE PURE-SINE ROW IS COMPUTED, NOT SIMULATED: an ideal excitation on the
%% analytic load with a single-pole Z_t fitted to the record's own zt_ac, point
%% sampled on the two interleaved grids and demodulated exactly as the extractor
%% does.  The same model driven with a two-code square returns +0.2727 % /
%% -2.9802 deg against the campaign's own sampling_prediction.txt value of
%% +0.2728 % / -2.9807 deg at 64 conversions per period -- four decimals, which
%% is what licenses the row.  The -180 f/f_s row is exact arithmetic on f/f_s.
%% Every other row is measured, save the deskewed row, which is a measured read
%% less that exact angle.
%%
%% Caveats that had to survive into the caption: both blocks are referred to the
%% AC ANALYSIS OF THE SAME LOAD AND GAIN CODE, whose own error against the
%% analytic load is +0.295 % and +0.405 deg (against the analytic load the codes
%% read +0.639 % and -2.533 deg); ONE record, ONE die, ONE frequency, ONE
%% resistive load, no noise sources in the deck -- a deterministic attribution
%% and not a distribution, so there is nothing behind this table to quote a
%% sigma from.  The record is demodulated over cycles 1-2; cycle 0 carries the
%% TIA's step response to the first DAC edge and the last cycle demodulates
%% 1.6 deg off in the boxcar path (eisdeep S10).  The complex-load code record
%% (codeC_rc_n0) saturated at 136 % of the compliance window and is excluded, so
%% this is a resistive-load result -- which is also why the magnitude
%% attribution is unambiguous.
%% Library cell names are kept out of the rendered prose.
\begin{table}[htbp]
  \centering
  \caption[{Impedance from converter codes: magnitude attribution and the phase
  mechanism.}]{Impedance demodulated from the ten-bit words of the extracted
  converter, on one record of \num{256} conversions at
  \SI{869565}{\per\second} with the input multiplexer interleaving the
  transimpedance output and the reference electrode, \num{64} conversions per
  excitation period at \SI{13.587}{\kilo\hertz}, the excitation synthesised by
  the real \num{12}-bit ladder into \SI{11.000}{\kilo\ohm} at gain code 20.
  The magnitude rows are three successive readings of that one record and
  accumulate to the total; the phase rows are independent readings of it, and
  \(-180\,f/f_s\) is the analytic half-conversion grid offset that the two
  interleaved readings sit on. The reference for both blocks is the
  small-signal analysis of the same load and gain code, itself \(+0.295\,\%\)
  and \SI{0.405}{\degree} from the analytic load. \textbf{One record, one die,
  one frequency into a resistive load}, with no noise sources in the deck: a
  deterministic attribution with no distribution behind it.}
  \label{tab:eis-codelevel}
  \begin{tabular}{llrcl}
    \toprule
    Quantity & Unit & Value & \(\pm\SI{0.573}{\degree}\) gate & Reading \\
    \midrule
    \multicolumn{5}{l}{\emph{Magnitude, against the small-signal analysis}} \\
    Transient path                 & points & \(+0.126\) & & measured \\
    Sampled system                 & points & \(+0.855\) & & measured \\
    Gain mismatch                  & points & \(-0.638\) & & measured \\
    Total, from codes              & points & \(+0.343\) & & measured \\
    \midrule
    \multicolumn{5}{l}{\emph{Phase, departure from the small-signal analysis}} \\
    From codes, interleaved        & \si{\degree} & \(-2.938\) & --- & measured \\
    Sampled, interleaved           & \si{\degree} & \(-2.895\) & --- & measured \\
    Predicted, \(-180\,f/f_s\)     & \si{\degree} & \(-2.813\) & --- & analytic \\
    Analog, no sampling            & \si{\degree} & \(-0.013\) & \checkmark & measured \\
    Deskewed by \(180\,f/f_s\)     & \si{\degree} & \(-0.082\) & \checkmark & corrected \\
    Two records, one grid          & \si{\degree} & \(+0.083\) & \checkmark & measured \\
    Pure sine, same grids          & \si{\degree} & \(0.000\)  & \checkmark & analytic control \\
    \bottomrule
  \end{tabular}
\end{table}


\paragraph{Measured through the complete channel.} A square-wave record taken
through the shipping channel, with the excitation entering at the reference
ladder and the true electrode current read in series with the working
electrode, separates the demodulator from the instrument.
\textbf{The demodulator, the two-code square and the boxcar contribute
\SI{0.01}{\percent} of magnitude and \SI{0.01}{\degree}}; the settled
impedance is within \SI{0.015}{\percent} of the model on the typical load,
inside \SI{0.15}{\percent} at every gain code, and inside
\SI{0.025}{\percent} across the corner set. The converter costs
\(+0.160\,\%\) and \SI{0.063}{\degree}, reconstructing its analog
input to \SI{0.305}{LSB} rms. What remains is the instrument, and it is two
terms that do not overlap (Table~\ref{tab:eis-pixel}): the feedback pole, which
is \(\arctan(2\pi f R_f C_F)\) and which the demodulator removes exactly by
dividing by \(R_f/(1 + \mathrm{j}\omega R_f C_F)\) rather than by \(R_f\),
and the control amplifier tracking the commanded potential, \SI{0.15}{\percent}
and \SI{0.09}{\degree} at \SI{1}{\kilo\hertz} rising to \SI{2.1}{\percent}
and \SI{8.4}{\degree} at \SI{50}{\kilo\hertz}. With the pole restored and the
record taken ratiometrically against a known impedance, charge-transfer
resistance and double-layer capacitance come out within
\SI{0.15}{\percent} at \SI{1}{\kilo\hertz} and \SI{0.40}{\percent} at
\SI{5}{\kilo\hertz}. \textbf{That last claim assumes the electrolyte
resistance is supplied}: taken as the real part at the top sweep frequency, as
the firmware note describes, it carries a \(+9.9\,\%\) bias on this load
before any measurement error and degrades both extracted parameters by an order
of magnitude.

%% VestaRV: hand-written TRM table fragment, maestro2tex house style; nothing here is
%% generated -- six transient records demodulated by script, not a Maestro
%% results database, and there is no maestro2tex config for any *_eis_*
%% fragment).  ADDED 2026-09-05.
%%
%% WHY IT EXISTS.  Every other impedance number in this section comes from a
%% small-signal replica bench that carried NO feedback capacitor, and from a
%% single code-level record.  This table is the same measurement taken through
%% the complete channel with the 22 pF present: one square-wave record per
%% frequency, the true electrode current read from a series 0 V source, and
%% three demodulations of each record so the error sources separate instead of
%% summing.  It is what tells the reader that the demodulator is exact and the
%% band edge is instrument, not algorithm.
%%
%% SOURCE OF EVERY NUMBER: ~/chips/castalia/simarchive/tb_anatop_eissq/
%%   20260905_sq_freq/demod_arc.txt (and demod_arc.json) -- the arc load,
%%     R_s 1 kOhm, R_ct 10 kOhm, C_dl 3.18 nF (corner 5.0 kHz), gain code 20,
%%     one recorded excitation period per frequency, converter idle.  The four
%%     demodulations of each record are z_cell (potential and current both
%%     taken at the cell), z_iwe (commanded potential over true current),
%%     z_raw (the firmware nominal-R_f law) and z_cf (the same with the
%%     feedback pole restored, Z_t = R_f/(1 + j w R_f C_F)).
%%   20260905_sq_freq/demod_res.txt -- the 11.000 kOhm calibration load at the
%%     same gain code: err(z_cell) 0.002 to 0.007 % and under 0.06 deg at every
%%     frequency, which is the check that the demodulator carries no bias.
%%   20260905_sq_code/table_sq_code.txt, codes_main.json -- the converter row:
%%     32 conversions per period at 27 173.913 Hz through the extracted
%%     converter, static multiplexer slot, code residual 0.305 LSB rms.
%%   20260905_sq_step/table_sq_step.txt -- the settled-impedance row.
%% The nominal-law phase column is atan(2 pi f R_f C_F) with R_f = 138 kOhm at
%% gain code 20 and C_F = 22 pF: 1.093 deg predicted against 1.159 measured at
%% 1 kHz, 5.45 against 5.69 at 5 kHz, 10.80 against 11.47 at 10 kHz.
%% Rendered prose and caption carry no em or en dashes.
\par\addvspace{\intextsep}
\noindent\parbox[b]{\textwidth}{%
  \captionof{table}[{Impedance error through the complete channel, by source.}]{Error
  of one square-wave record per frequency, demodulated four ways so the sources
  separate: at the cell, at the commanded potential, through the nominal
  feedback-resistance law, and through that law with the feedback pole
  restored. Conditions are the linear cell model at
  \(R_u = \SI{1}{\kilo\ohm}\), \(R_{\mathrm{ct}} = \SI{10}{\kilo\ohm}\),
  \(C_{\mathrm{dl}} = \SI{3.18}{\nano\farad}\), gain code 20, typical corner,
  \SI{37}{\celsius}, one recorded excitation period, converter idle. The
  nominal-law column is \(\arctan(2\pi f R_f C_F)\) to within \SI{6}{\percent}
  at every frequency, so it is the feedback pole and nothing else; what
  survives the correction is the control amplifier tracking the commanded
  potential, which is why the band edge exists at all.}%
  \label{tab:eis-pixel}%
  \par\nobreak
  \centering
  \begin{tabular}{lrrrrr}
    \toprule
    & \multicolumn{1}{c}{True \(|Z|\)}
    & \multicolumn{1}{c}{At the cell}
    & \multicolumn{1}{c}{\(+\) tracking}
    & \multicolumn{1}{c}{\(+\) nominal law}
    & \multicolumn{1}{c}{\(+\) law, pole restored} \\
    \midrule
    \multicolumn{6}{l}{\emph{Magnitude error}} \\
    \SI{1}{\kilo\hertz}  & \SI{10788.6}{\ohm} & \(+0.034\,\%\) & \(+0.179\,\%\) & \(-0.238\,\%\) & \(-0.256\,\%\) \\
    \SI{2}{\kilo\hertz}  & \SI{10221.4}{\ohm} & \(+0.049\,\%\) & \(+0.171\,\%\) & \(-0.235\,\%\) & \(-0.307\,\%\) \\
    \SI{5}{\kilo\hertz}  & \SI{7814.0}{\ohm}  & \(+0.031\,\%\) & \(-0.011\,\%\) & \(-0.108\,\%\) & \(-0.559\,\%\) \\
    \SI{10}{\kilo\hertz} & \SI{5003.7}{\ohm}  & \(-0.007\,\%\) & \(-0.568\,\%\) & \(+0.647\,\%\) & \(-1.135\,\%\) \\
    \SI{20}{\kilo\hertz} & \SI{2841.1}{\ohm}  & \(-0.011\,\%\) & \(-1.821\,\%\) & \(+4.488\,\%\) & \(-2.376\,\%\) \\
    \SI{50}{\kilo\hertz} & \SI{1480.0}{\ohm}  & \(+0.003\,\%\) & \(+2.123\,\%\) & \(+40.61\,\%\)  & \(+1.751\,\%\) \\
    \midrule
    \multicolumn{6}{l}{\emph{Phase error}} \\
    \SI{1}{\kilo\hertz}  & \SI{-10.259}{\degree} & \(+0.078\) & \(+0.170\) & \(+1.159\) & \(+0.066\) \\
    \SI{2}{\kilo\hertz}  & \SI{-19.702}{\degree} & \(+0.061\) & \(+0.195\) & \(+2.274\) & \(+0.089\) \\
    \SI{5}{\kilo\hertz}  & \SI{-39.783}{\degree} & \(+0.015\) & \(+0.294\) & \(+5.692\) & \(+0.244\) \\
    \SI{10}{\kilo\hertz} & \SI{-53.118}{\degree} & \(-0.007\) & \(+0.641\) & \(+11.472\) & \(+0.672\) \\
    \SI{20}{\kilo\hertz} & \SI{-55.985}{\degree} & \(-0.005\) & \(+2.078\) & \(+23.088\) & \(+2.206\) \\
    \SI{50}{\kilo\hertz} & \SI{-42.038}{\degree} & \(+0.002\) & \(+8.385\) & \(+52.348\) & \(+8.703\) \\
    \bottomrule
  \end{tabular}%
  \par\smallskip
  {\footnotesize Phase errors in degrees. Gain code 20 puts the feedback pole
  at \SI{52.4}{\kilo\hertz}.}%
}
\par\addvspace{\intextsep}


% HAND-WRITTEN fragment -- NOT generated by maestro2tex. Input from
% AnalogChapter.tex after note_eis_excitation.tex, last in the impedance
% section.  ADDED 2026-08-26 (owner request): "a few equations for how firmware
% turns the sequence of digitised values into R and C, just the equations".
%
% Nothing here is a new measurement.  Every equation is the algebra already
% implied by statements elsewhere in this chapter, written out:
%   (1) the demodulator: note_eis_intro, "an I/Q multiply and accumulate over an
%       integer number of excitation periods, no window, no detrending".  The
%       normalisation 2/MN is what makes X-hat the AMPLITUDE phasor: for
%       c_n = A cos(2 pi n/N + phi), sum c_n exp(-j 2 pi n/N) = (MN/2) A e^{j phi}.
%   (2) the transimpedance sign: note_ampab_tia sss:ampab-tia states
%       V_OUT = v_cm - I_WE R_f for A2, so in the fundamental bin (where the dc
%       common mode does not appear) I_WE-hat = -V_OUT-hat / R_f.  The leading
%       minus in eq:eis-z-rf is that inversion and nothing else.
%   (3) the LSB cancels because Z is a ratio of two phasors taken from the same
%       converter; likewise the 4/pi fundamental of the square (note_eis_excitation,
%       "the demodulator is ratiometric and the excitation's shape cancels").
%   (4) the ratiometric form is the architecture this section requires four times
%       over (feedback pole, code-pair error, quantisation bias, sigma = 11 % on
%       R_f).  It is architecture "C" of the source campaign.
%   (5) the cell model is note_models: R_ct || C_dl in series with R_u.
%       Z - R_u = R_ct/(1 + j w R_ct C_dl), so 1/(Z - R_u) = 1/R_ct + j w C_dl.
%       That identity is the whole reason the extraction is two divisions and no
%       fit: R_ct falls out of the real part and C_dl out of the imaginary part
%       of ONE inverted mid-band point.
%
% POTENTIAL CHANNEL, deliberate wording.  V_P-hat is written as "whatever the
% converter multiplexer presents as the potential channel" rather than as the RE
% tap, because on the shipping cell it CANNOT be the RE tap: ADC_MUX has no
% reference-electrode input (sss:pixtop-adc), so slot 2, the commanded potential,
% is what firmware actually demodulates.  That is architecture "D" of the source
% campaign, and the gap to "A" is the 0.02-0.2 % of magnitude quoted in
% note_eis_accuracy on this axis and the up-to-3.5 % of note_eis_excitation on
% the excitation axis.  Do NOT tighten this to "the reference electrode": it
% would state a connection the die does not have.
%
% No number is introduced here.  Every figure quoted is a back-reference.
% Rendered prose carries no em or en dashes.

\subsubsection{From Codes to Resistance and Capacitance}\label{sss:eis-firmware}

Firmware holds two code sequences per frequency, the potential channel and the
transimpedance output, taken on a grid of \(N\) samples per excitation period
over \(M\) whole periods. One multiply and accumulate per channel reduces each
to its fundamental phasor,
\begin{equation}\label{eq:eis-demod}
  \hat{X} \;=\; \frac{2}{MN}\sum_{n=0}^{MN-1} c_n\,
                \mathrm{e}^{-\mathrm{j}2\pi n/N},
  \qquad \mathrm{j}^2 = -1,
\end{equation}
giving \(\hat{V}_{\mathrm{P}}\) and \(\hat{V}_{\mathrm{OUT}}\). The converter LSB
is never applied: it cancels in every ratio below, as does the \(4/\pi\)
fundamental of the two-level excitation.

Since A2 holds the working electrode at the common mode across \(R_f\)
(Section~\ref{sss:ampab-tia}), the electrode current in that bin is
\(\hat{I}_{\mathrm{WE}} = -\hat{V}_{\mathrm{OUT}}/R_f\), so with a per-chip dc
calibration of the feedback resistance
\begin{equation}\label{eq:eis-z-rf}
  \hat{Z} \;=\; \frac{\hat{V}_{\mathrm{P}}}{\hat{I}_{\mathrm{WE}}}
          \;=\; -R_f\,\frac{\hat{V}_{\mathrm{P}}}{\hat{V}_{\mathrm{OUT}}},
\end{equation}
the minus sign being the transimpedance inversion. Measuring a known
\(Z_{\mathrm{cal}}\) at the same frequency and gain code removes \(R_f\)
altogether, and with it the feedback pole and the quantisation bias:
\begin{equation}\label{eq:eis-z-ratio}
  \hat{Z} \;=\; Z_{\mathrm{cal}}\,
    \frac{\hat{V}_{\mathrm{P}}\,\hat{V}^{\,\mathrm{cal}}_{\mathrm{OUT}}}
         {\hat{V}_{\mathrm{OUT}}\,\hat{V}^{\,\mathrm{cal}}_{\mathrm{P}}}.
\end{equation}

Against the cell model of Section~\ref{ss:analog-models},
\begin{equation}\label{eq:eis-cell}
  Z(\omega) \;=\; R_u \;+\; \frac{R_{\mathrm{ct}}}
                                 {1 + \mathrm{j}\omega R_{\mathrm{ct}} C_{\mathrm{dl}}},
\end{equation}
the extraction is two divisions and no curve fit. Take \(R_u\) as
\(\Re\,\hat{Z}\) at the top of the band, where the parallel branch has
collapsed, then invert one mid-band point:
\begin{equation}\label{eq:eis-extract}
  Y \;=\; \frac{1}{\hat{Z} - R_u}
    \;=\; \frac{1}{R_{\mathrm{ct}}} + \mathrm{j}\omega C_{\mathrm{dl}},
  \qquad
  R_{\mathrm{ct}} = \frac{1}{\Re\,Y},
  \qquad
  C_{\mathrm{dl}} = \frac{\Im\,Y}{\omega}.
\end{equation}
Both readings come from the same inversion, so a single frequency inside
\SI{0.1}{\hertz} to \SI{5.3}{\kilo\hertz} yields both, and the accuracy of each
is the accuracy of \(\hat{Z}\) at that frequency, which the calibration paragraphs above bound.

\noindent The potential phasor \(\hat{V}_{\mathrm{P}}\) is whatever the
converter multiplexer presents as that channel. On this cell that is the commanded potential,
not a reference-electrode tap, which no multiplexer slot carries
(the converter readout note records the slot map); the difference is worth
\SIrange{0.02}{0.2}{\percent} of magnitude here and up to \SI{3.5}{\percent} on
the excitation (the excitation and averaging note carries that figure).

\FloatBarrier

% Master include for Adc. HAND-WRITTEN: the fragments of this section are not
% maestro2tex output -- they are derived from the archived CSV named in each
% fragment header, simarchive/tb_anatop_pixel_adc/adc_transfer.csv.
%%
%% Usage from the TRM (which builds with cwd = latex/TRM/):
%     % Master include for Adc. HAND-WRITTEN: the fragments of this section are not
% maestro2tex output -- they are derived from the archived CSV named in each
% fragment header, simarchive/tb_anatop_pixel_adc/adc_transfer.csv.
%%
%% Usage from the TRM (which builds with cwd = latex/TRM/):
%     % Master include for Adc. HAND-WRITTEN: the fragments of this section are not
% maestro2tex output -- they are derived from the archived CSV named in each
% fragment header, simarchive/tb_anatop_pixel_adc/adc_transfer.csv.
%%
%% Usage from the TRM (which builds with cwd = latex/TRM/):
%     \input{include/analog/Adc.tex}
%
%% Needs pgfplots, booktabs, siunitx, graphicx and placeins -- all already in
%% packages-commands.tex. \providecolor needs xcolor (also present).

\providecommand{\MaestroRoot}{include/analog/}

% Categorical palette: slots 1-5 of the validated reference palette,
% in documented order. Dash patterns are a secondary encoding for print.
\providecolor{mtxA}{HTML}{2A78D6}
\providecolor{mtxB}{HTML}{EB6834}
\providecolor{mtxC}{HTML}{1BAF7A}
\providecolor{mtxD}{HTML}{EDA100}
\providecolor{mtxE}{HTML}{E87BA4}

% Placeholder float macro, used below for any figure of this section that has
% not been drawn yet.
\input{\MaestroRoot duckplaceholder.tex}

% Keep this block's floats inside this block (see placeins).
\FloatBarrier

\input{\MaestroRoot note_adc_intro.tex}
\IfFileExists{\MaestroRoot sch_adc_sar.tex}%
  {\input{\MaestroRoot sch_adc_sar.tex}}%
  {\duckplaceholder{Generic 10-bit charge-redistribution SAR converter:
    sampling network, capacitor DAC, comparator and successive-approximation
    register.}{fig:adc-sar}}
\input{\MaestroRoot fig_adc_capdac_layout.tex}
\input{\MaestroRoot note_adc_transfer.tex}
\input{\MaestroRoot fig_adc_transfer.tex}
\input{\MaestroRoot fig_adc_error.tex}
\input{\MaestroRoot fig_adc_range.tex}
\input{\MaestroRoot fig_adc_inl.tex}
\input{\MaestroRoot tab_adc_transfer.tex}
\input{\MaestroRoot tab_adc_corners.tex}
\input{\MaestroRoot note_adc_power.tex}
\input{\MaestroRoot tab_adc_power.tex}
\input{\MaestroRoot note_adc_mc.tex}

\FloatBarrier

%
%% Needs pgfplots, booktabs, siunitx, graphicx and placeins -- all already in
%% packages-commands.tex. \providecolor needs xcolor (also present).

\providecommand{\MaestroRoot}{include/analog/}

% Categorical palette: slots 1-5 of the validated reference palette,
% in documented order. Dash patterns are a secondary encoding for print.
\providecolor{mtxA}{HTML}{2A78D6}
\providecolor{mtxB}{HTML}{EB6834}
\providecolor{mtxC}{HTML}{1BAF7A}
\providecolor{mtxD}{HTML}{EDA100}
\providecolor{mtxE}{HTML}{E87BA4}

% Placeholder float macro, used below for any figure of this section that has
% not been drawn yet.
\input{\MaestroRoot duckplaceholder.tex}

% Keep this block's floats inside this block (see placeins).
\FloatBarrier

\input{\MaestroRoot note_adc_intro.tex}
\IfFileExists{\MaestroRoot sch_adc_sar.tex}%
  {\input{\MaestroRoot sch_adc_sar.tex}}%
  {\duckplaceholder{Generic 10-bit charge-redistribution SAR converter:
    sampling network, capacitor DAC, comparator and successive-approximation
    register.}{fig:adc-sar}}
\input{\MaestroRoot fig_adc_capdac_layout.tex}
\input{\MaestroRoot note_adc_transfer.tex}
\input{\MaestroRoot fig_adc_transfer.tex}
\input{\MaestroRoot fig_adc_error.tex}
\input{\MaestroRoot fig_adc_range.tex}
\input{\MaestroRoot fig_adc_inl.tex}
\input{\MaestroRoot tab_adc_transfer.tex}
\input{\MaestroRoot tab_adc_corners.tex}
\input{\MaestroRoot note_adc_power.tex}
\input{\MaestroRoot tab_adc_power.tex}
\input{\MaestroRoot note_adc_mc.tex}

\FloatBarrier

%
%% Needs pgfplots, booktabs, siunitx, graphicx and placeins -- all already in
%% packages-commands.tex. \providecolor needs xcolor (also present).

\providecommand{\MaestroRoot}{include/analog/}

% Categorical palette: slots 1-5 of the validated reference palette,
% in documented order. Dash patterns are a secondary encoding for print.
\providecolor{mtxA}{HTML}{2A78D6}
\providecolor{mtxB}{HTML}{EB6834}
\providecolor{mtxC}{HTML}{1BAF7A}
\providecolor{mtxD}{HTML}{EDA100}
\providecolor{mtxE}{HTML}{E87BA4}

% Placeholder float macro, used below for any figure of this section that has
% not been drawn yet.
\input{\MaestroRoot duckplaceholder.tex}

% Keep this block's floats inside this block (see placeins).
\FloatBarrier

\input{\MaestroRoot note_adc_intro.tex}
\IfFileExists{\MaestroRoot sch_adc_sar.tex}%
  {\input{\MaestroRoot sch_adc_sar.tex}}%
  {\duckplaceholder{Generic 10-bit charge-redistribution SAR converter:
    sampling network, capacitor DAC, comparator and successive-approximation
    register.}{fig:adc-sar}}
\input{\MaestroRoot fig_adc_capdac_layout.tex}
\input{\MaestroRoot note_adc_transfer.tex}
\input{\MaestroRoot fig_adc_transfer.tex}
\input{\MaestroRoot fig_adc_error.tex}
\input{\MaestroRoot fig_adc_range.tex}
\input{\MaestroRoot fig_adc_inl.tex}
\input{\MaestroRoot tab_adc_transfer.tex}
\input{\MaestroRoot tab_adc_corners.tex}
\input{\MaestroRoot note_adc_power.tex}
\input{\MaestroRoot tab_adc_power.tex}
\input{\MaestroRoot note_adc_mc.tex}

\FloatBarrier

% HAND-WRITTEN fragment -- NOT generated by maestro2tex, and Adc.tex is itself
% hand-written; kept separate so the section frame stays untouched. Input from
% AnalogChapter.tex immediately after Adc.tex.
%
% Salvaged from the retired note_calibration_seq.tex step 3 (converter zero) and
% note_calibration_residual.tex (the dither argument). Fabric from
% ~/vestarv/docs/afe_rev2_debug_mux.html §6.3: AFE_ADCSRCSEL[2:0], 0 = vout,
% 1 = vcm, 2 = vref, 3 = RE, 4 = ATB; §9 step 10 is the external ramp;
% per-channel slot 9 = adc_in, observable between conversions only.
% Numbers: tab_adc_transfer (1.535x slope, 0.441 V intercept, 1.591 mV LSB);
% DAC step 528.96 uV (tab_dacr2r12_mc_reference). No cell names in prose.

\subsubsection{Measuring this block on chip}

The converter is the one block whose gain has never been checked in system, and
the one that can be characterised end to end with no analog pad measurement.
Set \texttt{AFE\_ADCSRCSEL}~=~4 to route the test bus into the converter input,
drive an external ramp onto \texttt{atb\_d} with \texttt{AFE\_DBGDIREN}~=~1 and
\texttt{AFE\_DBGBUFEN}~=~0, and fit slope and intercept per die: the
\(1.535\times\) slope, \SI{0.441}{\volt} code-zero intercept and
\SI{1.591}{\milli\volt} effective LSB of Table~\ref{tab:adc-transfer} rest on a
single nominal extracted run with no chip-to-chip data behind them. For the zero
code, set \texttt{AFE\_ADCSRCSEL}~=~1 and average the common-mode rail; that
code is the origin of the current axis and carries the converter's offset and
the common-mode value together, which is the combination a current reading
needs. Averaging reaches the quantisation floor only if the input dithers, and
with the electrode isolated the amplifier output noise is tens of microvolts ---
so step the common-mode DAC over three of its \SI{529}{\micro\volt} codes, which
span one converter code. Observe the sampling node itself at per-channel slot 9
between conversions only. This test fabric is design intent; it is not yet in
the schematics.

% maestro2tex -- generated master include for TiaRprog.
%% Regenerate with tools/maestro2tex; do not edit by hand.
%%
%% Usage from the TRM (which builds with cwd = latex/TRM/):
%     % maestro2tex -- generated master include for TiaRprog.
%% Regenerate with tools/maestro2tex; do not edit by hand.
%%
%% Usage from the TRM (which builds with cwd = latex/TRM/):
%     % maestro2tex -- generated master include for TiaRprog.
%% Regenerate with tools/maestro2tex; do not edit by hand.
%%
%% Usage from the TRM (which builds with cwd = latex/TRM/):
%     \input{include/analog/TiaRprog.tex}
%
%% Needs pgfplots, booktabs, siunitx and placeins -- all already in
%% packages-commands.tex. \providecolor needs xcolor (also present).

\providecommand{\MaestroRoot}{include/analog/}

% Categorical palette: slots 1-5 of the validated reference palette,
% in documented order. Dash patterns are a secondary encoding for print.
\providecolor{mtxA}{HTML}{2A78D6}
\providecolor{mtxB}{HTML}{EB6834}
\providecolor{mtxC}{HTML}{1BAF7A}
\providecolor{mtxD}{HTML}{EDA100}
\providecolor{mtxE}{HTML}{E87BA4}

% Keep this block's floats inside this block (see placeins).
\FloatBarrier

\input{\MaestroRoot note_tiarprog_intro.tex}
\input{\MaestroRoot tab_TiaRprog_corners.tex}
\input{\MaestroRoot fig_tiarprog_r_vs_code.tex}
\input{\MaestroRoot tab_tiarprog_corners.tex}
\input{\MaestroRoot fig_tiarprog_err_vs_code.tex}
\input{\MaestroRoot fig_tiarprog_lin_n0.tex}
\input{\MaestroRoot tab_tiarprog_lin.tex}
\input{\MaestroRoot note_tiarprog_mc.tex}
\input{\MaestroRoot tab_tiarprog_mc.tex}
\input{\MaestroRoot fig_tiarprog_mc_seam.tex}
\input{\MaestroRoot note_tiarprog_scope.tex}

\FloatBarrier

%
%% Needs pgfplots, booktabs, siunitx and placeins -- all already in
%% packages-commands.tex. \providecolor needs xcolor (also present).

\providecommand{\MaestroRoot}{include/analog/}

% Categorical palette: slots 1-5 of the validated reference palette,
% in documented order. Dash patterns are a secondary encoding for print.
\providecolor{mtxA}{HTML}{2A78D6}
\providecolor{mtxB}{HTML}{EB6834}
\providecolor{mtxC}{HTML}{1BAF7A}
\providecolor{mtxD}{HTML}{EDA100}
\providecolor{mtxE}{HTML}{E87BA4}

% Keep this block's floats inside this block (see placeins).
\FloatBarrier

\input{\MaestroRoot note_tiarprog_intro.tex}
\input{\MaestroRoot tab_TiaRprog_corners.tex}
\input{\MaestroRoot fig_tiarprog_r_vs_code.tex}
\input{\MaestroRoot tab_tiarprog_corners.tex}
\input{\MaestroRoot fig_tiarprog_err_vs_code.tex}
\input{\MaestroRoot fig_tiarprog_lin_n0.tex}
\input{\MaestroRoot tab_tiarprog_lin.tex}
\input{\MaestroRoot note_tiarprog_mc.tex}
\input{\MaestroRoot tab_tiarprog_mc.tex}
\input{\MaestroRoot fig_tiarprog_mc_seam.tex}
\input{\MaestroRoot note_tiarprog_scope.tex}

\FloatBarrier

%
%% Needs pgfplots, booktabs, siunitx and placeins -- all already in
%% packages-commands.tex. \providecolor needs xcolor (also present).

\providecommand{\MaestroRoot}{include/analog/}

% Categorical palette: slots 1-5 of the validated reference palette,
% in documented order. Dash patterns are a secondary encoding for print.
\providecolor{mtxA}{HTML}{2A78D6}
\providecolor{mtxB}{HTML}{EB6834}
\providecolor{mtxC}{HTML}{1BAF7A}
\providecolor{mtxD}{HTML}{EDA100}
\providecolor{mtxE}{HTML}{E87BA4}

% Keep this block's floats inside this block (see placeins).
\FloatBarrier

\input{\MaestroRoot note_tiarprog_intro.tex}
\input{\MaestroRoot tab_TiaRprog_corners.tex}
\input{\MaestroRoot fig_tiarprog_r_vs_code.tex}
\input{\MaestroRoot tab_tiarprog_corners.tex}
\input{\MaestroRoot fig_tiarprog_err_vs_code.tex}
\input{\MaestroRoot fig_tiarprog_lin_n0.tex}
\input{\MaestroRoot tab_tiarprog_lin.tex}
\input{\MaestroRoot note_tiarprog_mc.tex}
\input{\MaestroRoot tab_tiarprog_mc.tex}
\input{\MaestroRoot fig_tiarprog_mc_seam.tex}
\input{\MaestroRoot note_tiarprog_scope.tex}

\FloatBarrier

% HAND-WRITTEN fragment -- NOT generated by maestro2tex. Input from
% AnalogChapter.tex immediately after TiaRprog.tex, which is generated.
%
% Salvaged from the retired note_calibration_seq.tex steps 6, 7 and 9 and
% note_calibration_paths.tex C4 (on-chip gain reference resistor; it needs no
% slot of its own -- both ends are per-channel slots 5 and 7, see
% ~/vestarv/docs/afe_rev2_debug_mux.html §11.4). Slot 4 = vout; slots 7/8
% straddle the C1 working-electrode isolation switch, and their difference under
% a forced current is that switch's i*Ron. Reading at the pad rather than
% through the SAR is §9.2 step 5: the SAR window truncates railing samples.
% Numbers: tab_tiarprog_corners (+-32-35 % across corners), 1.65 kOhm closed
% switch stack and 1.26 LSB bottom-code nonlinearity (ss:tiarprog),
% tab_tiarprog_mc (0.04-0.36 % matching). No cell names in prose.

\subsubsection{Measuring this block on chip}

Force \(\pm\) full-scale current into the working-electrode pin with no cell
present and read the transimpedance output at per-channel slot 4 on
\texttt{atb\_d}, one gain code at a time. Read it at the pad and not through the
converter: the converter window truncates the samples that rail, and the
two-point slope wanted here is the end-to-end constant in amperes per code,
absorbing the resistor's \(\pm35\,\%\) sheet tolerance
(Table~\ref{tab:tiarprog-corners}), the \SI{1.65}{\kilo\ohm} closed-switch stack
and the converter's own gain in a single number. Slots 7 and 8 straddle the
working-electrode isolation switch, so their difference under the same forced
current is that switch's \(i\!\cdot\!R_{\mathrm{on}}\), which calibration stores
rather than minimises. To refresh the gain without equipment, switch the on-chip
reference resistor between the counter-electrode node and the summing node, take
the same two-point slope, and scale the stored constant by the change in that
ratio: sheet resistance cancels, leaving matching alone at
\SIrange{0.04}{0.36}{\percent} (Table~\ref{tab:tiarprog-mc}). Choose the drive
amplitude per gain code, or the high-gain codes rail, and avoid the bottom code
entirely --- its drive nonlinearity is 1.26 LSB typical and no two-point slope
sees it. This test fabric is design intent; it is not yet in the schematics.

% maestro2tex -- generated master include for DacR2R12.
%% Regenerate with tools/maestro2tex; do not edit by hand.
%%
%% Usage from the TRM (which builds with cwd = latex/TRM/):
%     % maestro2tex -- generated master include for DacR2R12.
%% Regenerate with tools/maestro2tex; do not edit by hand.
%%
%% Usage from the TRM (which builds with cwd = latex/TRM/):
%     % maestro2tex -- generated master include for DacR2R12.
%% Regenerate with tools/maestro2tex; do not edit by hand.
%%
%% Usage from the TRM (which builds with cwd = latex/TRM/):
%     \input{include/analog/DacR2R12.tex}
%
%% Needs pgfplots, booktabs, siunitx and placeins -- all already in
%% packages-commands.tex. \providecolor needs xcolor (also present).

\providecommand{\MaestroRoot}{include/analog/}

% Categorical palette: slots 1-5 of the validated reference palette,
% in documented order. Dash patterns are a secondary encoding for print.
\providecolor{mtxA}{HTML}{2A78D6}
\providecolor{mtxB}{HTML}{EB6834}
\providecolor{mtxC}{HTML}{1BAF7A}
\providecolor{mtxD}{HTML}{EDA100}
\providecolor{mtxE}{HTML}{E87BA4}

% Keep this block's floats inside this block (see placeins).
\FloatBarrier

\input{\MaestroRoot note_dacr2r12_intro.tex}
\input{\MaestroRoot sch_dacr2r12_ladder.tex}
\input{\MaestroRoot tab_DacR2R12_corners.tex}
\input{\MaestroRoot fig_dacr2r12_inl_ideal.tex}
\input{\MaestroRoot fig_dacr2r12_dnl_ideal.tex}
\input{\MaestroRoot tab_dacr2r12_ideal.tex}
% note_dacr2r12_driver: HAND-WRITTEN. The bit driver was replaced on
% 2026-09-05 (1.0 V standard cell on the 2.5 V rail -> custom 2.5 V cell).
% Both the corner tables above and the Monte Carlo fragments below now
% read the new cell, so the fragment is the record of the change and of
% its mechanism, not a correction. Sources: F16 report, simarchive/
% tb_anatop_pixel_dacR2R12/20260905_driverfix/ and, for the five corner
% points, .../20260906_corners_driverfix/.
\input{\MaestroRoot note_dacr2r12_driver.tex}
\input{\MaestroRoot note_dacr2r12_mc.tex}
\input{\MaestroRoot fig_dacr2r12_mc_inl.tex}
\input{\MaestroRoot fig_dacr2r12_mc_dnl.tex}
\input{\MaestroRoot tab_dacr2r12_mc_ideal.tex}
\input{\MaestroRoot tab_dacr2r12_mc_scale.tex}
% EVERYTHING FROM HERE TO note_dacr2r12_scope CHARACTERISES THE DAC SUPPLY
% BLOCK anatop_pixel_dac_psu, WHICH WAS DELETED FROM anatop_pixel ON
% 2026-08-24; both ladders now run from the 2.5 V analog rail directly.
% The fragments are kept and moved behind note_dacr2r12_supply, which was
% retitled as the removed-block subsection and carries the deletion note, so
% that every \ref into them stays live and no reader takes a VddDac number
% for the tapeout channel. Ordering change 2026-09-06; the fragments and
% their labels are untouched.
\input{\MaestroRoot note_dacr2r12_supply.tex}
\input{\MaestroRoot sch_dacr2r12_psu.tex}
\input{\MaestroRoot fig_dacr2r12_vddac.tex}
\input{\MaestroRoot tab_dacr2r12_reference.tex}
\input{\MaestroRoot fig_dacr2r12_inl_wpsu.tex}
\input{\MaestroRoot fig_dacr2r12_dnl_ratio.tex}
\input{\MaestroRoot tab_dacr2r12_wpsu.tex}
\input{\MaestroRoot tab_dacr2r12_mc_wpsu.tex}
\input{\MaestroRoot fig_dacr2r12_mc_vddacpp.tex}
\input{\MaestroRoot tab_dacr2r12_mc_reference.tex}
\input{\MaestroRoot fig_dacr2r12_psu_load.tex}
\input{\MaestroRoot tab_dacr2r12_loadreg.tex}
\input{\MaestroRoot fig_dacr2r12_mc_loadreg.tex}
\input{\MaestroRoot tab_dacr2r12_mc_psu_load.tex}
\input{\MaestroRoot fig_dacr2r12_psu_line.tex}
\input{\MaestroRoot tab_dacr2r12_linereg.tex}
\input{\MaestroRoot tab_dacr2r12_mc_psu_line.tex}
\input{\MaestroRoot fig_dacr2r12_psu_zout.tex}
\input{\MaestroRoot tab_dacr2r12_zout.tex}
\input{\MaestroRoot tab_dacr2r12_mc_psu_zout.tex}
\input{\MaestroRoot fig_dacr2r12_psu_startup.tex}
\input{\MaestroRoot note_dacr2r12_scope.tex}

\FloatBarrier

%
%% Needs pgfplots, booktabs, siunitx and placeins -- all already in
%% packages-commands.tex. \providecolor needs xcolor (also present).

\providecommand{\MaestroRoot}{include/analog/}

% Categorical palette: slots 1-5 of the validated reference palette,
% in documented order. Dash patterns are a secondary encoding for print.
\providecolor{mtxA}{HTML}{2A78D6}
\providecolor{mtxB}{HTML}{EB6834}
\providecolor{mtxC}{HTML}{1BAF7A}
\providecolor{mtxD}{HTML}{EDA100}
\providecolor{mtxE}{HTML}{E87BA4}

% Keep this block's floats inside this block (see placeins).
\FloatBarrier

\input{\MaestroRoot note_dacr2r12_intro.tex}
\input{\MaestroRoot sch_dacr2r12_ladder.tex}
\input{\MaestroRoot tab_DacR2R12_corners.tex}
\input{\MaestroRoot fig_dacr2r12_inl_ideal.tex}
\input{\MaestroRoot fig_dacr2r12_dnl_ideal.tex}
\input{\MaestroRoot tab_dacr2r12_ideal.tex}
% note_dacr2r12_driver: HAND-WRITTEN. The bit driver was replaced on
% 2026-09-05 (1.0 V standard cell on the 2.5 V rail -> custom 2.5 V cell).
% Both the corner tables above and the Monte Carlo fragments below now
% read the new cell, so the fragment is the record of the change and of
% its mechanism, not a correction. Sources: F16 report, simarchive/
% tb_anatop_pixel_dacR2R12/20260905_driverfix/ and, for the five corner
% points, .../20260906_corners_driverfix/.
\input{\MaestroRoot note_dacr2r12_driver.tex}
\input{\MaestroRoot note_dacr2r12_mc.tex}
\input{\MaestroRoot fig_dacr2r12_mc_inl.tex}
\input{\MaestroRoot fig_dacr2r12_mc_dnl.tex}
\input{\MaestroRoot tab_dacr2r12_mc_ideal.tex}
\input{\MaestroRoot tab_dacr2r12_mc_scale.tex}
% EVERYTHING FROM HERE TO note_dacr2r12_scope CHARACTERISES THE DAC SUPPLY
% BLOCK anatop_pixel_dac_psu, WHICH WAS DELETED FROM anatop_pixel ON
% 2026-08-24; both ladders now run from the 2.5 V analog rail directly.
% The fragments are kept and moved behind note_dacr2r12_supply, which was
% retitled as the removed-block subsection and carries the deletion note, so
% that every \ref into them stays live and no reader takes a VddDac number
% for the tapeout channel. Ordering change 2026-09-06; the fragments and
% their labels are untouched.
\input{\MaestroRoot note_dacr2r12_supply.tex}
\input{\MaestroRoot sch_dacr2r12_psu.tex}
\input{\MaestroRoot fig_dacr2r12_vddac.tex}
\input{\MaestroRoot tab_dacr2r12_reference.tex}
\input{\MaestroRoot fig_dacr2r12_inl_wpsu.tex}
\input{\MaestroRoot fig_dacr2r12_dnl_ratio.tex}
\input{\MaestroRoot tab_dacr2r12_wpsu.tex}
\input{\MaestroRoot tab_dacr2r12_mc_wpsu.tex}
\input{\MaestroRoot fig_dacr2r12_mc_vddacpp.tex}
\input{\MaestroRoot tab_dacr2r12_mc_reference.tex}
\input{\MaestroRoot fig_dacr2r12_psu_load.tex}
\input{\MaestroRoot tab_dacr2r12_loadreg.tex}
\input{\MaestroRoot fig_dacr2r12_mc_loadreg.tex}
\input{\MaestroRoot tab_dacr2r12_mc_psu_load.tex}
\input{\MaestroRoot fig_dacr2r12_psu_line.tex}
\input{\MaestroRoot tab_dacr2r12_linereg.tex}
\input{\MaestroRoot tab_dacr2r12_mc_psu_line.tex}
\input{\MaestroRoot fig_dacr2r12_psu_zout.tex}
\input{\MaestroRoot tab_dacr2r12_zout.tex}
\input{\MaestroRoot tab_dacr2r12_mc_psu_zout.tex}
\input{\MaestroRoot fig_dacr2r12_psu_startup.tex}
\input{\MaestroRoot note_dacr2r12_scope.tex}

\FloatBarrier

%
%% Needs pgfplots, booktabs, siunitx and placeins -- all already in
%% packages-commands.tex. \providecolor needs xcolor (also present).

\providecommand{\MaestroRoot}{include/analog/}

% Categorical palette: slots 1-5 of the validated reference palette,
% in documented order. Dash patterns are a secondary encoding for print.
\providecolor{mtxA}{HTML}{2A78D6}
\providecolor{mtxB}{HTML}{EB6834}
\providecolor{mtxC}{HTML}{1BAF7A}
\providecolor{mtxD}{HTML}{EDA100}
\providecolor{mtxE}{HTML}{E87BA4}

% Keep this block's floats inside this block (see placeins).
\FloatBarrier

\input{\MaestroRoot note_dacr2r12_intro.tex}
\input{\MaestroRoot sch_dacr2r12_ladder.tex}
\input{\MaestroRoot tab_DacR2R12_corners.tex}
\input{\MaestroRoot fig_dacr2r12_inl_ideal.tex}
\input{\MaestroRoot fig_dacr2r12_dnl_ideal.tex}
\input{\MaestroRoot tab_dacr2r12_ideal.tex}
% note_dacr2r12_driver: HAND-WRITTEN. The bit driver was replaced on
% 2026-09-05 (1.0 V standard cell on the 2.5 V rail -> custom 2.5 V cell).
% Both the corner tables above and the Monte Carlo fragments below now
% read the new cell, so the fragment is the record of the change and of
% its mechanism, not a correction. Sources: F16 report, simarchive/
% tb_anatop_pixel_dacR2R12/20260905_driverfix/ and, for the five corner
% points, .../20260906_corners_driverfix/.
\input{\MaestroRoot note_dacr2r12_driver.tex}
\input{\MaestroRoot note_dacr2r12_mc.tex}
\input{\MaestroRoot fig_dacr2r12_mc_inl.tex}
\input{\MaestroRoot fig_dacr2r12_mc_dnl.tex}
\input{\MaestroRoot tab_dacr2r12_mc_ideal.tex}
\input{\MaestroRoot tab_dacr2r12_mc_scale.tex}
% EVERYTHING FROM HERE TO note_dacr2r12_scope CHARACTERISES THE DAC SUPPLY
% BLOCK anatop_pixel_dac_psu, WHICH WAS DELETED FROM anatop_pixel ON
% 2026-08-24; both ladders now run from the 2.5 V analog rail directly.
% The fragments are kept and moved behind note_dacr2r12_supply, which was
% retitled as the removed-block subsection and carries the deletion note, so
% that every \ref into them stays live and no reader takes a VddDac number
% for the tapeout channel. Ordering change 2026-09-06; the fragments and
% their labels are untouched.
\input{\MaestroRoot note_dacr2r12_supply.tex}
\input{\MaestroRoot sch_dacr2r12_psu.tex}
\input{\MaestroRoot fig_dacr2r12_vddac.tex}
\input{\MaestroRoot tab_dacr2r12_reference.tex}
\input{\MaestroRoot fig_dacr2r12_inl_wpsu.tex}
\input{\MaestroRoot fig_dacr2r12_dnl_ratio.tex}
\input{\MaestroRoot tab_dacr2r12_wpsu.tex}
\input{\MaestroRoot tab_dacr2r12_mc_wpsu.tex}
\input{\MaestroRoot fig_dacr2r12_mc_vddacpp.tex}
\input{\MaestroRoot tab_dacr2r12_mc_reference.tex}
\input{\MaestroRoot fig_dacr2r12_psu_load.tex}
\input{\MaestroRoot tab_dacr2r12_loadreg.tex}
\input{\MaestroRoot fig_dacr2r12_mc_loadreg.tex}
\input{\MaestroRoot tab_dacr2r12_mc_psu_load.tex}
\input{\MaestroRoot fig_dacr2r12_psu_line.tex}
\input{\MaestroRoot tab_dacr2r12_linereg.tex}
\input{\MaestroRoot tab_dacr2r12_mc_psu_line.tex}
\input{\MaestroRoot fig_dacr2r12_psu_zout.tex}
\input{\MaestroRoot tab_dacr2r12_zout.tex}
\input{\MaestroRoot tab_dacr2r12_mc_psu_zout.tex}
\input{\MaestroRoot fig_dacr2r12_psu_startup.tex}
\input{\MaestroRoot note_dacr2r12_scope.tex}

\FloatBarrier

% HAND-WRITTEN fragment -- NOT generated by maestro2tex. Input from
% AnalogChapter.tex immediately after DacR2R12.tex, which is generated.
%
% Salvaged from the retired note_calibration_seq.tex steps 2 and 8. Fabric from
% ~/vestarv/docs/afe_rev2_debug_mux.html: per-channel slot 2 = vref, gated by
% AFE_DBGISOEN because leakage on an R-2R output is a code-dependent error;
% chip-level slot 7 = VddDac, slot 6 = shared vcm; AFE_ADCSRCSEL = 2 reads vref
% through the SAR (§6.3). Timing: one conversion 1.15 us at the DESIGN rate
% (note_adc_intro; 4.600 us on the clock pattern as delivered, 2026-08-25), so
% 4096 codes x 16 averages = 75 ms per channel, 302 ms as delivered. Droop fit and its 94 % / 931 uV
% rms residual from sss:dacr2r12-supply; per-code table costs 4 KiB per channel
% and buys 0.07 mV rms. No cell names in prose.

\subsubsection{Measuring this block on chip}

\begin{enumerate}[nosep]
\item Select per-channel slot 2 with \texttt{AFE\_DBGISOEN}~=~1 and read the
  reference output on \texttt{atb\_d}. The isolation switch is required, not
  optional: leakage on an R-2R output is a code-dependent error.
\item Absolute anchor: take three codes with a voltmeter to fix the DAC's
  offset and step in volts. Nothing on the die can supply that scale; it
  anchors the channel's whole potential axis.
\item Sweep all 4096 codes into the converter with
  \texttt{AFE\_ADCSRCSEL}~=~2 (\SI{75}{\milli\second} per channel at sixteen
  averages and the design conversion rate, four times that on the clock
  pattern as delivered --- Section~\ref{sss:adc-power}) and fit the
  one-parameter droop of
  Section~\ref{sss:dacr2r12-supply}: it removes \SI{94}{\percent} of the
  \SI{25.8}{\milli\volt} reference movement. Do not store the per-code
  residual table --- through this converter each entry carries as much
  quantisation as the error it corrects (\SI{0.07}{\milli\volt}~rms for
  4\,KiB per channel) --- but run the sweep regardless: it is where a broken
  ladder shows.
\item Chip-level slot 7 reads the DAC supply node directly --- the only place
  the \SI{25.8}{\milli\volt} code-dependent excursion is confirmed or refuted
  on silicon. Chip-level slot 6 reads the shared common-mode DAC, forceable
  through \texttt{AFE\_VCMISOEN}.
\end{enumerate}
This test fabric is design intent; it is not yet in the schematics.


% Four-channel macro, hand-written (no maestro2tex config; the numbers are read from
% the executed psf of four transient Maestro tests by skill/quad_eval.py). It is last
% among the measured sections and before the calibration section: it assembles every
% block above it, and the calibration section refers to per-channel constants that only
% mean something once the channels are shown to be independent. Runs
% ~/chips/castalia/simruns/quad_20260905/, wrapper netlist and CDL
% ~/chips/castalia/simarchive/anatop_quad/.
\FloatBarrier
% VestaRV: hand-written TRM fragment, not maestro2tex output (no config behind it;
% the numbers come from four transient Maestro tests on the four-channel
% wrapper, evaluated on the executed psf).  ADDED 2026-09-05.
%
% WHAT IT IS.  Every other section of this chapter measures one block or one
% channel.  This one measures the four-channel macro: four complete channels,
% each with its own reference and common-mode DACs, its own feedback ladder,
% its own converter and its own local bias generator, sharing only the analog
% supplies and the two test pads.  It is TB-09 of the verification plan,
% adapted to the per-site DACs the pixel actually has.
%
% SOURCE OF EVERY NUMBER: F13 report, 2026-09-05.
%   Cell castalia/anatop_quad (82 instances, 150 nets, 38 terminals, schCheck
%     0 errors, 0 unbound and 0 unconnected instTerms, 260/260 stub audit);
%     bench castalia/tb_anatop_quad (62 instances, schCheck 0 errors).
%   Runs ~/chips/castalia/simruns/quad_20260905/<test>/eval_<test>.txt, read
%     from the psf by skill/quad_eval.py; netlist and full CDL of the wrapper
%     archived at ~/chips/castalia/simarchive/anatop_quad/.
%   Conditions: nominal corner, 37 C, castalia_typ_25, converter av_extracted,
%     20 MHz trigger clock at SAMPLESTEP 7 (1.150 us per conversion), ideal
%     2.5 V and 1.0 V supplies, one linear cell model per site.
%   top_op   10 uA out of every working electrode at gain code 0: raw 121 ->
%     code 633 on all four sites, site-to-site spread 0 codes, vout spread
%     42 uV, RE - V_PATTERN -0.56 mV and WE - vcm -0.59 mV on every site.
%   top_mix  +20 uA at code 0, -20 uA at code 0, +1 uA at code 63,
%     -0.2 uA at code 319, simultaneously: codes 769 / 265 / 756 / 267.
%   top_xtalk site 0 driven -20/+20 uA at a 4-conversion period, sites 1-3 at
%     zero: quiet sites read code 508 and 1.7511 / 1.7502 / 1.7511 mV pp, which
%     is 1.10 LSB.  The same measure on top_op, where NO neighbour steps at
%     all, is 2.10 / 2.13 / 2.10 mV pp, so the ripple is each channel's own
%     converter sampling its own transimpedance output, not coupling.  This
%     bench has ideal rails and per-site bias, so it has no crosstalk path to
%     measure; the number is reported as a bound, not as a crosstalk figure.
%   top_pwr  i_avdd 1.815 mA at 2.5 V for four channels, i_vdd10 14.2 uA at
%     1.0 V, p_total 4.55 mW = 1.14 mW per channel, averaged over the last
%     conversion.  i_avdd_pk 42.6 mA is the start-up inrush into an ideal
%     source with no series impedance and is not an operating number.
%   The constant code_err of -4.4 LSB on every site and every test is the
%     converter window centre (1.254 V) against vcm 1.2496 V plus the zero-code
%     convention: a property of the expected-code expression, not a channel
%     error, and it cancels in any vcm-calibrated reading.
% NOT MEASURED HERE, deliberately: corners and Monte Carlo (TB-09 is a nominal
% first pass), and crosstalk with a real supply and substrate network, which
% needs the layout.  Both are stated in the prose.
% Rendered prose and caption carry no em or en dashes.

\subsection{Four-Channel Macro}\label{ss:quad}

The four channels are independent by construction: each carries its own
reference and common-mode converters, its own feedback ladder, its own
ten-bit converter and its own local bias generator, and shares only the analog
supplies and the two test pads. Running all four at once is therefore a check
that the assembly is right, not a test of a shared resource.

\textbf{Site-to-site spread is zero codes.} With \SI{10}{\micro\ampere} drawn
from every working electrode at the bottom gain code, all four sites return code
\num{633}, their transimpedance outputs agreeing to \SI{42}{\micro\volt}, and
the control and transimpedance pedestals reproduce the single-channel figures
(\SI{-0.56}{\milli\volt} of applied-potential error, \SI{-0.59}{\milli\volt} at
the working electrode) rather than adding to them. Driven with a different
current, a different polarity and a different gain code per site,
simultaneously, the four codes land symmetrically about the common-mode code at
\num{+257}, \num{-247}, \num{+244} and \num{-245} for \SI{+20}{\micro\ampere},
\SI{-20}{\micro\ampere}, \SI{+1}{\micro\ampere} and
\SI{-0.2}{\micro\ampere} across \num{107}-fold of gain.

\textbf{The macro draws \SI{4.55}{\milli\watt}}, \SI{1.14}{\milli\watt} per
channel, of which the digital rail is \SI{14.2}{\micro\ampere}. The
neighbour-disturbance column of Table~\ref{tab:quad-tb09} is
\textbf{not a crosstalk measurement}: this bench has ideal supplies and per-site
bias, so it carries no coupling path, and the same \SI{1.1}{LSB} of ripple
appears with no neighbour stepping at all. It is each channel's own converter
sampling its own transimpedance output. A crosstalk figure needs the
post-layout supply and substrate network. Corners and Monte Carlo were not run
on the macro; every block below it carries its own.

%% HAND-WRITTEN table fragment (maestro2tex house style; nothing generated).
%% ADDED 2026-09-05.  Source: F13 report and
%% ~/chips/castalia/simruns/quad_20260905/<test>/eval_<test>.txt, read from the
%% executed psf by skill/quad_eval.py.  note_quad_tb09.tex's header carries the
%% per-test conditions and the two caveats (the disturbance column is not
%% crosstalk; the peak supply current is a start-up inrush into an ideal
%% source).  Codes are after the bit-9 inversion the converter macro needs.
\par\addvspace{\intextsep}
\noindent\parbox[b]{\textwidth}{%
  \captionof{table}[{Four-channel macro, all sites running simultaneously.}]{The
  four channels driven together at the nominal corner, \SI{37}{\celsius}, with
  the converter as its extracted view and a \SI{20}{\mega\hertz} trigger clock
  (\SI{1.150}{\micro\second} per conversion). The mixed-condition rows carry a
  different current, polarity and gain code per site at the same instant. The
  disturbance column is measured with only site~0 stepping, but the same figure
  appears when no site steps at all, so it is each channel's own converter
  sampling its own output and not coupling: this bench has ideal supplies and
  per-site bias and therefore no coupling path. Nominal only; corners and Monte
  Carlo live in the block sections.}%
  \label{tab:quad-tb09}%
  \par\nobreak
  \centering
  \begin{tabular}{lrrrr}
    \toprule
    & Site 0 & Site 1 & Site 2 & Site 3 \\
    \midrule
    \multicolumn{5}{l}{\emph{Identical drive, \SI{10}{\micro\ampere} at gain code 0}} \\
    Code returned                   & 633 & 633 & 633 & 633 \\
    Transimpedance output (\si{\volt}) & 1.449192 & 1.449195 & 1.449234 & 1.449195 \\
    Applied-potential error (\si{\milli\volt}) & \(-0.56\) & \(-0.56\) & \(-0.56\) & \(-0.56\) \\
    Working-electrode error (\si{\milli\volt}) & \(-0.59\) & \(-0.59\) & \(-0.59\) & \(-0.59\) \\
    \midrule
    \multicolumn{5}{l}{\emph{Mixed drive, simultaneous}} \\
    Current forced (\si{\micro\ampere}) & \(+20\) & \(-20\) & \(+1\) & \(-0.2\) \\
    Gain code                       & 0 & 0 & 63 & 319 \\
    Code, referred to common mode   & \(+257\) & \(-247\) & \(+244\) & \(-245\) \\
    \midrule
    \multicolumn{5}{l}{\emph{Neighbour disturbance, site 0 stepping \SI{40}{\micro\ampere}}} \\
    Quiet-site output (\si{\milli\volt} pp) & --- & 1.751 & 1.750 & 1.751 \\
    \midrule
    \multicolumn{5}{l}{\emph{Power, all four channels}} \\
    \SI{2.5}{\volt} supply          & \multicolumn{4}{r}{\SI{1.815}{\milli\ampere}} \\
    \SI{1.0}{\volt} supply          & \multicolumn{4}{r}{\SI{14.2}{\micro\ampere}} \\
    Total                           & \multicolumn{4}{r}{\SI{4.55}{\milli\watt}, \SI{1.14}{\milli\watt} per channel} \\
    \bottomrule
  \end{tabular}%
}
\par\addvspace{\intextsep}



% Per-chip calibration, hand-written and last: a system statement built on every block
% section above and referring to all of them by \ref, so it must follow them. Only the
% programmable knobs live here; the measurements are in the note_meas_*.tex fragments.
% Design intent only, no calibration has been run.
\FloatBarrier
% HAND-WRITTEN fragment -- NOT generated by maestro2tex. Input directly from
% AnalogChapter.tex, after the block sections: it refers to all of them by \ref.
%
% Replaces the six-fragment calibration section of 2026-08-21
% (note_calibration_{intro,paths,seq,residual}.tex + tab_calibration_errors.tex),
% which the owner judged far too long. The per-block measurement recipes moved
% into the block sections themselves as note_meas_*.tex; what remains here is
% only what firmware programs. Deleted fragments are recoverable from git.
%
% Sources of the numbers below, so they can be re-checked:
%   bias trim span / spread   tab_biasgen_trim; MC sigma/mu 12.5 % (bg_mc, Plan.4.Run.2)
%   A1 offset                 tab_ampab_mc_offset, tab_ampab_mc_tracking
%   zero-current offset       tab_px_mc_izero, sss:pixel-mc
%   reference droop           note_dacr2r12_supply (R^2 = 94 %, residual 931 uV rms
%                             against 25.70 mV pp), tab_dacr2r12_mc_reference
%   ladder INL                tab_dacr2r12_mc_ideal (1.49 DAC LSB mean = 0.788 mV)
%   converter step            tab_adc_transfer (1.591 mV effective LSB)
%   Rf tolerance / matching   tab_tiarprog_corners, tab_tiarprog_mc
%   quantisation floor        1.591/sqrt(12) = 0.459 mV; two conversions = 0.650 mV
%   potential residual        RSS(0.650, 0.931, 0.788) = 1.38 mV rms
%   current residual          RSS(0.650, 0.294) = 0.713 mV rms = 0.45 LSB on the
%                             measured 1.591 mV step. Published as 0.71 LSB until
%                             2026-08-25: the RSS is of two VOLTAGES (mV), and the
%                             unit was asserted, not converted.
%   range condition           0.815 V / (3 x 5.7176 mV) = 47.5
%   uncalibrated potential sigma  5.7176 mV (tab_ampab_mc_tracking)
% Library cell names are deliberately kept OUT of the rendered prose (owner
% request 2026-08-21); they live in these comments instead.

\subsection{Per-Chip Calibration}\label{ss:calibration}

Every correction is a stored digital constant except one, the bias generator's
6-bit trim code. Table~\ref{tab:calibration-constants} is the complete set of
knobs: what each corrects, how it is measured, and when it is applied. The
measurement itself is described in each block's section, under \emph{Measuring
this block on chip}.

Only two constants are absolute --- one voltage at the reference DAC output, one
current forced at the working electrode --- and both need the production tester.
The other eight are zeros and ratios the channel measures against itself through
its own converter, which is what lets the power-up refresh run with no equipment.

Order matters twice. At test, apply the bias trim before anything else is
measured: it moves every transconductance behind every constant that follows. In
a reading, subtract the zero code, scale by the gain constant for the active gain
code, then subtract the drop across the working-electrode switch and, where an
impedance sweep has supplied it, the uncompensated electrolyte resistance.
Potential-axis corrections apply where the waveform is tabulated, not per sample.

\input{\MaestroRoot tab_calibration_constants.tex}

\paragraph{Residual.} Each zero is the difference of two conversions, so its
floor is \SI{0.65}{\milli\volt} rms, two quantisations of the converter's
\SI{1.591}{\milli\volt} step; the common-mode DAC has to dither the input to
reach it. On the potential axis the residual is \SI{1.38}{\milli\volt} rms ---
that floor, the \SI{0.93}{\milli\volt} droop-fit residual and
\SI{0.79}{\milli\volt} of uncorrected ladder INL in quadrature --- against
\SI{5.71}{\milli\volt} of uncalibrated \(\sigma\). That holds over the
specified \SIrange{0.4}{2.1}{\volt} window only: outside it the tracking error
is no longer a constant and no stored constant removes it
(Section~\ref{sss:pixtop-mc}). On the current axis it is
\SI{0.71}{\milli\volt} rms --- \SI{0.45}{LSB} on the converter's measured
\SI{1.591}{\milli\volt} step (Table~\ref{tab:adc-transfer}) --- plus
\SIrange{0.1}{0.4}{\percent} of reading, against \SIrange{12.1}{167.3}{LSB}
uncalibrated, which are on the nominal \SI{2.441}{\milli\volt} step.
A constant removes the current zero only
while that zero stays inside the converter window, which requires
\(R_f/R_{\mathrm{cell}} < 47.5\); the \SI{560}{\kilo\ohm} tap into a
\SI{10}{\kilo\ohm} electrode (Section~\ref{sss:pixel-mc}) is outside it and those
samples are recoverable by nothing.

\paragraph{Out of reach.} The reference DAC's \(-6.6\)~LSB of DNL is lost command
resolution, up to \SI{3.5}{\milli\volt} at the worst code, and no constant
creates a code the ladder does not produce. A reference-electrode junction
potential is indistinguishable from A1's offset, so only a zero taken \emph{wet}
against that electrode absorbs it. The temperature dependence of every constant
here is unquantified: no block in this chapter has a temperature-resolved Monte
Carlo, so how far a power-up re-measurement tracks a \SI{37}{\celsius} operating
point is not known.

\FloatBarrier

% Orphans: the fragments in this directory that no \input reaches. 227 .tex files
% here, 193 reachable from this file, 34 not. They are kept with their own provenance
% headers, and restoring one is putting an \input back. Recompute the list with a
% transitive \input walk from this file, not by eye.
%
%   Assembled channel, 17 (the commented-out block above):
%     note_pixtop_{intro,potential,zero,cv,adc,mux,comp,noise,power,mc,open}
%     tab_pixtop_{summary,mc}  fig_pixtop_{cv,noise,mc}  note_meas_pixtop
%   Impedance spectroscopy, 9: note_eis_{coverage,excitation,open,rules} and
%     fig_eis_{error,settling,mccal,nyquist,codelevel}. tab_eis_codelevel replaced
%     fig_eis_codelevel, and the section renders no figure at all.
%   Build helpers, 7: preview_{Adc,AmpAB,BiasGenCascWideSwing,DacR2R12,Electrode,
%     Pixel,TiaRprog} -- standalone wrappers for compiling one block section on its
%     own; never part of the manual.
%   Unused schematic, 1: sch_tb_slew (the slew bench, never written up).
%
% Nothing here is superseded: every number in these fragments is still traceable to the
% run named in its header. They are out of the manual because the chapter was
% shortened, not because they are wrong.
